\chapter{Wortbildung}\label{chapter:Wortbildung}

In diesem Kapitel geht es um die Analyse und Bildung komplexer Wortstämme nach bestimmten Modellen bzw. Regeln. Unter\is{Wortbildung} Wortbildung versteht man die systematische Abwandlung von Grundbedeutungen und Stammformen. Aus dem maskulinen Substantiv \textit{Freund} können wir \textit{Tierfreund} bilden, aber auch \textit{freundlich}, \textit{Freundschaft} oder \textit{befreunden}. Um die unterschiedlichen wortbildenden Bestandteile komplexer Wörter beim Lesen schneller zu erfassen, verwenden wir in diesem Kapitel häufig eine Bindestrichschreibung, also \textit{Tier-freund}, \textit{freund-lich}, \textit{Freund-schaft} und \textit{be-freund}(\textit{en}). Die verbale Endung \textit{-en} setzen wir meist in Klammern, da es sich nicht um einen wortbildenden Bestandteil, sondern um eine Flexionsendung (vgl. Abschnitt~\ref{subsection:Verbflexion}) handelt. 

Bei Wortbildungsprozessen können wie bei dem maskulinen Substantiv \textit{Tier-freund} die Wortart und die Flexionsklasse unverändert bleiben, sie können aber auch wie bei dem femininen Substantiv \textit{Freund-schaft}, bei dem Adjektiv \textit{freund-lich} und bei dem Verb \textit{be-freund}(\textit{en}) verändert werden. 

In dem folgenden Abschnitt~\ref{section:Wortbildung_Einführung} grenzen wir die Wortbildung von der Wortformenbildung (Flexion) und von der Wortgruppenbildung (Syntax) ab. Danach werden in Abschnitt~\ref{section:Einheiten Wortbildung} die Einheiten der Wortbildung vorgestellt: lexikalische Stammformen wie \textit{Freund}, sogenannte Derivationsaffixe wie das Präfix \textit{be-} in \textit{be-freund}(\textit{en}) und das Suffix \textit{-schaft} wie in \textit{Freund-schaft}, Verbpartikeln wie \textit{an} in \textit{jemand freundet sich mit jemandem an} und Konfixe wie \textit{polit} in \textit{Polit-freund}. Danach geht es in Abschnitt~\ref{section:Komposition} um die Bildung von Komposita wie \textit{Tier-freund}. Hierbei wird ein Hauptaugenmerk auf den Prinzipien der Binarität und der Rechtsköpfigkeit liegen. In Abschnitt~\ref{section:Derivation} geht es unter dem Stichwort \textit{Derivation} um die Strukturiertheit komplexer Wörter mit Präfixen wie \textit{be-freund}(\textit{en}), mit Suffixen wie \textit{Freund-schaft} aber auch mit einem Wechsel des Stammvokals wie bei \textit{fallen} versus \textit{fällen}. Schließlich beschäftigen wir uns unter dem Stichwort \textit{Konversion} mit Wortbildungen ohne Formsignal, so \zB das Substantiv (\textit{das}) \textit{Grün} gegenüber dem Adjektiv \textit{grün}. Danach geht es in Abschnitt~\ref{section:Vpt Bildung} um Partikelverben wie \textit{an-freund}(\textit{en}) oder eben \textit{freundet an} als eigenen Wortbildungstyp. In Abschnitt~\ref{section:andere Wb-Arten} diskutieren wir verschiedene Wortbildungstypen im Spannungsfeld zwischen Komposition und Derivation. Abschließend erläutern wir in Abschnitt~\ref{section:Produktivität} die Begriffe \textit{Produktivität}, \textit{Transparenz} sowie \textit{Regel-} und \textit{Analogiebasiertheit}. Auf die Zusammenfassung in Abschnitt~\ref{section:Zfsg_Wortbildung} folgen in Abschnitt~\ref{section:BspAnalysenWB} vier Beispielanalysen zum Üben.   


\section{Wortbildung -- Gestalt(ung) vor Syntax und Flexion}\label{section:Wortbildung_Einführung}

Der wesentliche Unterschied\is{Wortbildung!versus Flexion} zwischen Wortbildung und Flexion besteht darin, dass letztere relativ strikt vorhersagbar und transparent ist: \textit{Freund} bildet ein Paradigma aus acht Flexionsformen, die sich durch die Kombination der zwei Numerus- mit den vier Kasuskategorien ergeben. Die Wortbildung ist hingegen weniger vorhersagbar und kann weniger transparent sein. Man kann mit jemandem \textit{be-freundet} sein, aber weniger gut \textit{be-feindet}. Das wäre \textit{ver-feindet}, wobei jedoch \textit{ver-freundet} nicht geht. \textit{Freund-chen} hat eine andere Wortqualität als \textit{Feind-chen}. Das ist auch bei \textit{An-feind-ung} im Vergleich mit dem möglichen, aber viel seltener gebrauchten \textit{An-freund-ung} der Fall. Manche Wortbildungen wie \textit{Freund-schaft-er} gibt es heute anders als den \textit{Gesell-schaft-er} nicht mehr. Das liegt auch daran, dass wir die Bedeutung von \textit{Freund-schaft} heute anders auf \textit{Freund} und \textit{-schaft} beziehen als bei \textit{Gesell-schaft} auf \textit{Geselle} und \textit{-schaft}.

Flexion ist kein wortbildender, kein wortartverändernder Prozess, sondern die klar geregelte Formung von Wortstämmen als Instanzen von Wortarten für ganz bestimmte Funktionen in Wortgruppen. Die Wortbildung geht der Flexion voraus. Formale Flexionsmittel folgen den Mitteln der Wortbildung. Schauen wir uns ein Beispiel an: \textit{freund}-\textit{lich}-\textit{es}. Zuerst bilden wir aus \textit{Freund} mit Hilfe des Suffixes \textit{-lich} das Adjektiv \textit{freundlich}. Danach wird das Adjektiv für den jeweiligen Gebrauch in einer Wortgruppe flektiert, \zB \textit{ein freundlich-es Lächeln} oder \textit{mit einem freundlich-en Lächeln}. Flexionsmittel kommen anders als Wortbildungsmittel nur ganz am Ende von Wortstämmen zum Einsatz. Diese absolute Endstellung sichert ihre schnelle Erkennbarkeit, die wesentlich ist, um die jeweilige syntaktische Einbindung zu erkennen. Nur selten können wir Mittel der Wortbildung wie das \textit{-s} in dem Adverb \textit{abend-s} auf Flexionsmittel wie die Genitivform \textit{eines Abend-s} zurückführen.

Die Wortbildung\is{Wortbildung!Ökonomie} ist ökonomisch, weil wir auf jeweils bekannte Bausteine, auf Wortstämme und Affixe sowie auf die Muster ihrer Verbindung zurückgreifen, anstatt Wörter ganz neu zu erfinden. Sie ist auch ökonomisch, weil die Produkte der Wortbildung im Rahmen der bestehenden Wortarten und Flexionsparadigmen gebraucht werden. Hier gilt im Besonderen, was \citet[42]{Humboldt1836} in der Sprache als \gqq{geistige Thätigkeit} bezeichnet: Sie ist \gqq{immer zugleich auf etwas schon Gegebenes gerichtet, nicht rein erzeugend, sondern umgestaltend}. 

Mit jener\is{Morphologie} Umgestaltung und Gestaltetheit von Wortstämmen beschäftigt sich die wortbildende Morphologie. Sie ist neben der Syntax und der Flexionsmorphologie der zweite große Bereich der Grammatik im traditionellen Sinne.

Im Gegensatz zu syntaktischen Strukturen sind die morphologischen weniger explizit, dafür aber kompakter und assoziativer. Mit dem Begriff \textit{Holz-kiste} wird ein Typ von Kiste bezeichnet, der in irgendeiner Beziehung zu Holz steht. Erst durch den Kontext verstehen wir, ob eine Kiste mit Holz, aus Holz oder für Holz gemeint ist. Durch Wortgruppen wird das von einem Substantiv Bezeichnete beschrieben: \textit{eine Kiste aus Holz} oder \textit{eine hölzerne Kiste} versus \textit{eine Kiste mit Holz}. Formal zeigt sich die Besonderheit der Wortbildung darin, dass ihre Einheiten nicht über syntaktisch motivierte Flexionsendungen verfügen: \textit{Holz-Kiste} versus \textit{hölzern-e Kiste}.    

Schnittstellen zur Syntax gibt es dann, wenn eine flektierte (also syntaktische) Wortform wie \textit{freundliche} aus \textit{der freundliche Mensch} sang- und klanglos in die Kategorie Substantiv wie \textit{der Freundliche} überführt wird oder wenn die Wortgruppe (VP) \textit{ja sagen} zum Substantiv \textit{Jasag-er} wird. Eine stabile Schnittstelle zur Syntax sind auch Partikelverben wie \textit{ab-reißen}. Denn ihre Struktur gleicht der Struktur von sogenannten Resultativkonstruktionen, deren Prädikat aus einem Verb und einem Adjektiv besteht wie \textit{kaputt reißen}. Auf historische Übergänge von der Syntax zur Morphologie werden wir in den einzelnen Abschnitten eingehen.  

Die hier skizzierten Übergänge und Schnittstellen machen die Wortbildung so interessant. Es geht um komplexe Wortstämme, um ihre Bedeutung und ihre Form im Spannungsfeld zwischen lautlicher Struktur und syntaktischer Verwendung.\footnote{In diesem Sinne bezeichnen \citet[1]{SpencerZwicky1998} die Morphologie als \gqq{conceptual centre of linguistics}. Ein Beispiel für die Relevanz der lautlichen Form ist die Ableitung von Adjektiven zu Substantiven durch \textit{-heit} versus \textit{-keit}. Ableitungen mit \textit{-heit} basieren auf endbetonten oder einsilbigen Adjektiven wie \textit{Gesundheit} oder \textit{Freiheit} im Gegensatz zu \textit{Brüderlichkeit}. Wir kommen in Abschnitt~\ref{section:Produktivität} hierauf zurück.} Den einzelnen Bildungen liegen Muster zugrunde, die mehr oder weniger produktiv sein können, worum es im Abschnitt~\ref{section:Produktivität} gehen wird.

Bevor es im nächsten Abschnitt~\ref{section:Einheiten Wortbildung} um die Einheiten der Wortbildung geht, geben wir in dem folgenden Exkurs~\ref{Exkurs:Wb_Schule} eine kurze Übersicht über die schulische Relevanz der Wortbildung. 

\newpage
\begin{Exkurs}{Wortbildung in der Schule}\label{Exkurs:Wb_Schule}

Das Thema Wortbildung gehört zwar mit der Komposition und der Derivation zum Kanon des schulischen Grammatikunterrichts. Aber im Deutschunterricht der Sekundarstufen I und II ist die Wortbildung ein vernachlässigter Lernbereich.\footnote{Cf. \citet[33]{Ulrich2021}.} Bei Wortbildung als Unterrichtsgegenstand geht es darum, dass Schülerinnen und Schüler praktisch lernen, komplexe Wörter auf der Grundlage einfacher Wörter zu bilden, also um Wortschatzerweiterung. Diese Wortschatzerweiterung dient unter anderem dazu, neue Dinge zu benennen (\textit{Fahrradstraße}), Dinge präziser zu benennen (\textit{Kirschsaft} statt \textit{Saft}) und textuelle Zusammenhänge wie in (\ref{ea:Wortbildung_Text}) durch Kursivdruck angezeigt, herzustellen.

\ea \label{ea:Wortbildung_Text} Die Polizei kann eine Person für längstens 24 Stunden \textit{verhaften}. Spätestens dann muss die Person dem Staatsanwalt übergeben werden. Innerhalb von 48 Stunden seit der \textit{Verhaftung} muss dieser entscheiden, ob er die Person freilassen oder in \textit{Haft} behalten will.\footnote{Tages-Anzeiger, 18.08.2011, S. 13.}
\z 

Die Erweiterung des Wortschatzes und textueller Fähigkeiten bedarf der sogenannten \textit{morphologischen Bewusstheit}. Unter morphologischer Bewusstheit wird seit \citet[194]{Carlisle1995} die Fähigkeit verstanden, die morphologische Struktur von Wörtern zu reflektieren und zu manipulieren, d.\,h. Wortfamilien zu bilden, Wortstämme zu identifizieren, Wortbausteine zu segmentieren und aus den bekannten einzelnen Bausteinen die Bedeutung komplexer Wörter zu erschließen. 
Nach \citet[92]{Seiffert2023} spielen die morphologische Bewusstheit und die mit ihr verbundene Wortbildungskompetenz in den folgenden sechs Lernbereichen des Deutschunterrichts eine wichtige Rolle.

\begin{enumerate}

\item Beim Lesen, und zwar bei der Dekodierung komplexer Wörter und dem Textverständnis.

\item Beim Rechtschreiben, wo insbesondere die Morphemkonstanz (vgl. Abschnitt~\ref{sec:MorphPrin}) eine wesentliche Rolle spielt, aber auch bei der Getrennt- und Zusammenschreibung (vgl. Abschnitt~\ref{sec:Getrennt_Zusammen}) und bei der satzinternen Großschreibung (vgl. Abschnitt~\ref{sec:satzinternGroß}).

\item Bei der Wortschatzarbeit.

\item Bei der Textproduktion, insbesondere bei der Herstellung textueller Zusammenhänge und bei stilistischer Variation. 

\item Bei der Textrezeption, bei der es auch um die Reflexion über künstlerischen Sprachgebrauch, über Mehrdeutigkeiten und über Neuschöpfungen geht.

\item Bei der Sprachreflexion über Varianten, \zB \textit{Rindsbraten} versus \textit{Rinderbraten} und bei der Beurteilung alternativer Ausdrucksmöglichkeiten wie \zB bei \textit{Flüchtling} versus \textit{Flüchtende} oder \textit{Geflüchtete}.

\end{enumerate}

Zukünftigen Lehrerinnen und Lehrern empfehlen wir die unterhaltsame Lektüre eines Artikels von \citet{Donalies2003Wbfrei} mit dem Titel \textit{Gebt endlich die Wortbildung frei! Über sinnige und unsinnige Kritik an der Wortbildung}. Denn es ist ein Plädoyer für einen freien, spielerisch-kreativen, aber nicht unkritischen Umgang mit Neubildungen und mit der Kritik an ihnen.  

\end{Exkurs}


  
\section{Einheiten der Wortbildung}\label{section:Einheiten Wortbildung}

Wir beginnen mit einem komplexen Wort wie in (\ref{ea:Bioeinkaufserlebnis}), welches es gilt, schrittweise in seine Bestandteile und damit in seine Konstituenten zu zerlegen. So gelangen wir schrittweise an die kleinsten Bausteine komplexer Wörter. 

\ea \label{ea:Bioeinkaufserlebnis} Bioeinkaufserlebnis
\z

Das Substantiv in (\ref{ea:Bioeinkaufserlebnis}) können wir in \textit{Bioeinkaufs} und \textit{Erlebnis} zerlegen. Die gegenüber dem Substantiv \textit{Bioeinkauf} besondere Substantivform \textit{Bioeinkaufs} verfügt über ein \textit{-s} als sogenanntes Fugenelement. Das Substantiv \textit{Bioeinkauf} könnten wir weiter zerlegen in \textit{bio} und das Substantiv \textit{Einkauf}.

Alternativ können wir das Substantiv \textit{Bioeinkaufserlebnis} auch zerlegen in \textit{bio} und das Substantiv \textit{Einkaufserlebnis}. Dieses Substantiv können wir weiter zerlegen in die besondere Substantivform \textit{Einkaufs} mit dem Fugenelement \textit{-s} und das Substantiv \textit{Erlebnis}.

Das Substantiv \textit{Einkauf} geht auf den komplexen Verbstamm \textit{einkauf-} zurück. Dieser kann zerlegt werden in die Verbpartikel \textit{ein} und den Verbstamm \textit{kauf-}. Das Substantiv \textit{Erlebnis} lässt sich zerlegen in den komplexen Verbstamm \textit{erleb-} und das Suffix \textit{-nis}. Der Verbstamm \textit{erleb-} kann weiter zerlegt werden in das Präfix \textit{er-} und den Verbstamm \textit{leb-}.

Somit haben wir das komplexe Wort \textit{Bioeinkaufserlebnis} in seine kleinsten bedeutungs- und funktionstragenden Bausteine zerlegt: \textit{bio}, \textit{ein}, \textit{kauf}, \textit{er}, \textit{leb} und \textit{nis}. Solche Wortbausteine nennt\is{Morphem} man Morpheme.

\begin{Definition}{Morphem}\label{def:Morphem}
    
\noindent Morpheme sind die kleinsten, nicht weiter zerlegbaren bedeutungs- oder funktionstragenden Einheiten der Sprache. 
\end{Definition}



Einem Morphemverständnis zufolge sind \textit{bio}, \textit{ein}, \textit{kauf}, \textit{er}, \textit{leb} und \textit{nis} eigenständige Einheiten des Lexikons. Diese Lexikoneinheiten werden durch bestimmte Regeln miteinander kombiniert, so \zB das Präfix \textit{er-} mit dem Verbstamm \textit{leb-} zu \textit{erleb-}. Dieser Ansatz ist morphembasiert.\footnote{So \zB die \textit{Distributed Morhology}. Cf. \citet{HarleyNoyer1999}.}

In den folgenden Abschnitten werden wir komplexe Wörter schrittweise in ihre Morpheme zerlegen. Denn dieses Vorgehen erlaubt Einsicht in die Struktur bestehender komplexer Wörter und in das Prinzip der Binarität. Bei der Komposition und der Suffigierung ermöglicht es darüber hinaus Einsicht in das Prinzip der Rechtsköpfigkeit.

Damit schließen wir uns jedoch dem morphembasierten Ansatz nicht an. Denn es ist durchaus fraglich, Präfixe wie \textit{er-} und Suffixe wie \textit{-nis} mit kaum fassbaren Bedeutungen genauso wie lexikalische Morpheme wie \textit{kauf} oder \textit{leb} als eigenständige Bestandteile des Lexikons zu klassifizieren. Stattdessen werden wir -- wie bei der Flexion (vgl. Abschnitt \ref{section:Flektierbare}) -- primär wortbasiert\footnote{Wortbasiert ist \zB der konstruktionsgrammatische Ansatz von \citet{Booij2010}.} vorgehen. Damit gehen wir nicht davon aus, dass Präfixe und Suffixe eigenständige Einheiten des Lexikons sind, sondern im Falle von Produktivität Bestandteile von schematischen Konstruktionen wie \zB [er-X]\textsubscript{Verb}. X ist auf Verbstämme festgelegt. Die so gebildeten transitiven Verben drücken aus, dass ein Resultat durch das vom Basisverb ausgedrückte Geschehen angestrebt oder erreicht wird: so \zB \textit{etwas erarbeiten}, \textit{etwas erfassen}, \textit{etwas erkaufen}, \textit{etwas errechnen} oder \textit{etwas ergoogeln}.

In den einzelnen Abschnitten werden wir auf wortbasierte Erklärungen zurückkommen. In dem folgenden Abschnitt beschäftigen wir uns mit verschiedenen Stammformen lexikalischer Morpheme.


\subsection{Stammformen der Wortbildung}\label{subsection:Stämme}

In der Wortbildung spielen drei verschiedene Stammformen eine Rolle: lexikalische Stammformen (vgl. Abschnitt~\ref{section:Wortbegriffe}), Kompositionsstammformen und Derivationsstammformen. 

Lexikalische Stammformen\is{Stammformen!lexikalisch} haben wir im Abschnitt~\ref{section:Wortbegriffe} als bedeutungstragende Grundbausteine eines Wortes minus Flexionsaffix definiert. Deshalb spricht man auch von Flexionsstammformen. Diese Stammformen bilden den Ausgangspunkt für die Ableitung von flektierten Wortformen. Bei nicht flexionsfähigen Stämmen kann einfach von lexikalischer Grundformen gesprochen werden. Das Substantiv \textit{Erlebnis} in \textit{Bioeinkaufserlebnis} ist solch eine lexikalische Stammform bzw. Flexionsstammform. Flexionsstammformen dienen als Zweitglieder von Komposita. In \textit{Bioeinkaufserlebnis} ist \textit{Erlebnis} eine komplexe Flexionsstammform. Denn sie besteht aus mehr als einem Morphem. Einfache, d.\,h. nur aus einem Morphem bestehende Flexionsstammformen sind \textit{Kind} in \textit{Kleinkind} oder \textit{blau} in \textit{hellblau}.

Auch das Erstglied\is{Stammformen!Komposition} \textit{bioeinkaufs-} ist eine Stammform. Diese ist jedoch für ihre Funktion als Erstglied eines Kompositums durch das Fugenelement \textit{-s} ausgezeichnet. Andere Fugenelemente sind \textit{-e} bei \textit{bade-} in \textit{Badeerlebnis}, \textit{-es} bei \textit{tages-} in \textit{Tageslicht} oder \textit{-er} bei \textit{kinder-} in \textit{Kindertraum}. Neben den Fugenelementen \textit{-s}, \textit{-e}, \textit{-es} und \textit{-er} gibt es noch \textit{-n} wie bei \textit{fugen} in \textit{Fugenelement}, \textit{-en} wie bei \textit{herren} in \textit{Herrentag} und äußerst selten \textit{-ens} wie bei \textit{herzens} in \textit{Herzensangelegenheit}.  

Mit \citet[96]{Wurzel1970} bezeichnen wir alle Erstglieder von Komposita als Kompositionsstammformen, auch dann, wenn diese wie \textit{tag-} in \textit{Tagtraum} kein Fugenelement aufweisen. Nach \citet[54]{Ortner_Komposita1991} entspricht die Kompositionsstammform in ca. 73\% der Fälle der lexikalischen Stammform, weist also kein Fugenelement auf. Kompositionsstammformen können sich auch wie \textit{sprach-} in \textit{Sprachkurs} oder \textit{schul} in \textit{Schulbeginn} durch Kürzungen gegenüber den lexikalischen Stammformen \textit{Sprache} und \textit{Schule} auszeichnen.\footnote{\citet[247]{Eisenberg2020_Wort} bezeichnet die Kürzung als \gqq{Subtraktionsfuge}, da hier der \textit{e}-Laut abgezogen wird.}

Um die Systematik von Fugenelementen\is{Fugenelement} geht es im Abschnitt~\ref{subsection:Fuge}. An dieser Stelle sei vorweggenommen, dass Fugenelemente grundsätzlich als Bestandteile der entsprechenden Stammform und nicht als Flexionsaffixe analysiert werden. Denn wie bei \textit{Liebesbrief} muss es keine entsprechende Flexionsform geben. \textit{Liebes} ist keine Flexionsform des femininen Substantivs \textit{Liebe}.\footnote{Die Genitivform von (\textit{die}) \textit{Liebe} lautet (\textit{der}) \textit{Liebe}. Die Flexion der Substantive wird in Abschnitt \ref{subsection:Substantivflexion} dargestellt.} Wenn auch Pluralformen die Kompositionsstammform wie in \textit{Bücherregal} semantisch motivieren, sind sie, wie \textit{Rinderbraten} zeigt, nicht durchgängig semantisch motiviert. Für die Zugehörigkeit des Fugenelements zur Kompositionsstammform spricht auch, dass bei verkürzten Kompositionsstammformen wie bei \textit{sprach-} in \textit{Sprachkurs} die Kürzung das Erstglied, also \textit{Sprache} betrifft. In Koordinationsellipsen wie \textit{Tages- und Abendkarte} verbleibt das Fugenelement beim Erstglied.

Um Derivationsstammformen handelt\is{Stammform!Derivation} es sich, wenn ein Stamm durch Derivationsaffixe, insbesondere durch Suffixe, abgewandelt wird. So ist \textit{freund-} in \textit{freundlich} eine Derivationsstammform. Sie unterscheidet sich nicht von der lexikalischen Stammform. Um formal ausgezeichnete Derivationsstammformen handelt es sich bei \textit{häus-} in \textit{häuslich}, \textit{blum-} in \textit{blumig} und \textit{hoffent-} in \textit{hoffentlich}. Viele Derivationsstammformen gehen auf lautliche Anpassungen (vgl. Abschnitt~\ref{subsection:Assimilation}) der lexikalischen Stammform an Suffixe zurück. So handelt es sich beim heutigen \textit{Vöglein} um die Hebung des Stammvokals [o] zu [{øː}] als Anpassung an ein folgendes langes [i], wie es beim mittelhochdeutschen  \textit{vogelîn} der Fall war.\footnote{In Abschnitt \ref{sec:Vokale} werden die Vokale des Deutschen und ihre Transkription vorgestellt.}

Derivations- und Kompositionsstammformen sind identisch, wenn die Suffixe von ursprünglichen lexikalischen Stammformen abstammen. Dann basiert die Ableitung historisch auf dem Muster der Komposition wie \textit{-haft} in \textit{frühlingshaft} oder \textit{-tum} in \textit{Heldentum}.\footnote{Das heutige Suffix \textit{-haft} geht auf ein gleichlautendes Adjektiv mit der Bedeutung \gq{gebunden/haftend} zurück. Das Suffix \textit{-tum} geht auf das Substantiv \textit{tuom} mit der Bedeutung \gq{Urteil/Entscheidung} zurück. Cf. \citet[390, 500]{Wilmanns1896Wb}.}


Die verschiedenen\is{Allomorph} Stammformen eines Morphems sind Allomorphe dieses Morphems. Allomorphe sind formale Varianten eines abstrakten Morphems, das heißt Varianten einer minimalen Einheit aus Form und Bedeutung. In Anlehnung an \citet[22--33]{Fuhrhop1998} stellen wir die Allomorphe des Morphems \textit{maus} in dem Stammparadigma in der Tabelle~\ref{tab:Stammparadigma} dar. 


\begin{table}[!htbp]
  \centering
  \begin{tabular}{l l l}
    \lsptoprule
    Stammform & Allomorph & Vorkommen im Wort \\
    \hline
    Flexionsstammform & \textit{maus}  & \textit{Maus}, \textit{Wühl-maus} \\
    & \textit{mäus-} & \textit{Mäus-e}, \textit{Mäus-en} \\
    
    Derivationsstammform & \textit{maus-}  & \textit{maus-haft} \\
    
    Kompositionsstammform & \textit{mäuse-} & \textit{Mäuse-busshard}, \textit{Mäuse-jagd} \\
    & \textit{mause-} & \textit{Mause-loch}, \textit{Mause-falle} \\
    \lspbottomrule 
  \end{tabular}
  \caption{Stammparadigma: Allomorphe des Morphems \textit{maus}}
  \label{tab:Stammparadigma}
\end{table}

Die Allomorphe eines Morphems lassen sich nicht regulär bilden. Sie müssen im Laufe des Spracherwerbs gelernt werden. 

\subsection{Derivationsaffixe}\label{subsection:Derivationsaffixe}

Viele\is{Derivationsaffix} Derivationsaffixe wie \textit{er-} in \textit{er-leben} oder \textit{-schaft} in \textit{Freund-schaft} haben sich aus lexikalischen Stammformen entwickelt.\footnote{So geht \textit{er-} auf die althochdeutsche Präposition \textit{ur} mit der Bedeutung \gq{aus} zurück, \textit{-schaft} auf das Substantiv \textit{scaf} mit der Bedeutung \gq{Beschaffenheit/Eigenschaft}. Cf. \citet[387]{Wilmanns1896Wb}.} Mitunter ist diese Entwicklung heute noch nachvollziehbar, wie bei \textit{-los} in \textit{ahnungslos}.\footnote{\citet[151]{EisenbergSayatz2002} bezeichnen diese Affixe folglich als \gqq{Kompositionsaffixe}.}

Derivationsaffixe kommen\is{Derivationsaffix} nur zusammen mit Stammformen oder mit Konfixen (vgl. Abschnitt~\ref{subsection:Konfixe}) vor. Sie können alleine nicht als syntaktisches Wort gebraucht werden, sie sind nicht derivierbar und nicht flektierbar. Hierauf zielt der \textit{Affix}begriff (von lat. \textit{affixum} \gq{Angeheftetes}) ab. Derivationsaffixe dienen als Operatoren der Abwandlung von Stammformen. In Bezug auf Derivationsaffixe wird die Stammform auch Basis genannt. Eine Basis kann einfach sein wie \textit{leb}(\textit{en}) in \textit{er-leb}(\textit{en}) oder komplex wie \textit{er-leb}(\textit{en}) in \textit{Erleb-nis}.   

Derivationsaffixe werden bezüglich ihrer Stellung zur Basis in Präfixe\is{Präfix} und Suffixe\is{Suffix} unterteilt. Präfixe stehen wie \textit{er-} in \textit{er-leb}(\textit{en}) vor der Basis. Suffixe stehen wie \textit{-nis} in \textit{Erleb-nis} nach der Basis. Viel seltener sind\is{Zirkumfix} Zirkumfixe, die ihre Basis wie \textit{Ge-e} in \textit{Ge-schimpf-e} oder \textit{be-ig} in \textit{be-grad-ig}(\textit{en}) umschließen. 

Die Ableitung durch Suffixe\is{Suffix} ist bei der Bildung von Substantiven und Adjektiven zentral. Wichtige Substantivsuffixe werden in der Tabelle~\ref{tab:Substantivsuffixe} aufgeführt, wichtige Adjektivsuffixe in der Tabelle~\ref{tab:Adjektivsuffixe}.

\begin{table}[!htbp]
  \centering
  \begin{tabular}{l l}
    \lsptoprule
    Substantivsuffixe & Beispiele \\
    \hline
    
    \textit{-chen} & \textit{Freund-chen}, \textit{Päck-chen}\\
    
    \textit{-er} & \textit{Bäck-er}, \textit{Mal-er}\\
    
    \textit{-}(\textit{er})\textit{ei} & \textit{Bäcker-ei}, \textit{Back-erei}\\
    
    \textit{-heit} & \textit{Frei-heit}, \textit{Gleich-heit}\\
    
    \textit{-keit} & \textit{Brüderlich-keit}, \textit{Ähnlich-keit}\\
    
    \textit{-nis} & \textit{Erleb-nis}, \textit{Finster-nis}\\
    
    \textit{-sal} & \textit{Trüb-sal}, \textit{Lab-sal}\\
    
    \textit{-schaft} & \textit{Freund-schaft}, \textit{Gesell-schaft}\\
    
    \textit{-tum} & \textit{Reich-tum}, \textit{Beamten-tum}\\
    
    \textit{-ung} & \textit{Liefer-ung}, \textit{Bestell-ung}\\
    
    \lspbottomrule 
  \end{tabular}
  \caption{Wichtige Substantivsuffixe}
  \label{tab:Substantivsuffixe}
\end{table}


\begin{table}[!htbp]
  \centering
  \begin{tabular}{l l}
    \lsptoprule
    Adjektivsuffixe & Beispiele \\
    \hline
    \textit{-bar} & \textit{trink-bar}, \textit{mach-bar}\\
    
    \textit{-}(\textit{e})\textit{n} & \textit{gold-en}, \textit{papier-(e)n}\\
    
    \textit{-haft} & \textit{schad-haft}, \textit{oberlehrer-haft}\\
    
    \textit{-ig} & \textit{stein-ig}, \textit{staub-ig}\\
    
    \textit{-isch} & \textit{neid-isch}, \textit{verführer-isch}\\
    
    \textit{-lich} & \textit{glück-lich}, \textit{bild-lich}\\
    
    \textit{-los} & \textit{angst-los}, \textit{brot-los}\\
    
    
    \lspbottomrule 
  \end{tabular}
  \caption{Wichtige Adjektivsuffixe}
  \label{tab:Adjektivsuffixe}
\end{table}

Suffixe legen die Wortart des abgeleiteten Wortes fest. Hierbei kommt es häufig, wie aus den Beispielen ersichtlich wird, zu einem Wechsel der Wortart. Viele Suffixe verbinden sich nur oder hauptsächlich mit Stammformen bestimmter Wortarten. So verbindet sich \textit{-ung} hauptsächlich mit Verbstämmen wie \textit{liefer-} zu \textit{Lieferung} und \textit{bestell-} zu \textit{Bestellung}. Hingegen verbindet sich \textit{-ig} mit Substantiven wie \textit{Stein} zu \textit{steinig} und \textit{Staub} zu \textit{staubig}.

Präfixe spielen\is{Präfix} eine zentrale Rolle bei der Verbbildung. Wichtige Verbpräfixe werden in der Tabelle~\ref{tab:Verbpräfixe} aufgeführt.

\begin{table}[!htbp]
  \centering
  \begin{tabular}{l l}
    \lsptoprule
    Verbpräfixe & Beispiele \\
    \hline
    \textit{be-} & \textit{be-antworten}, \textit{be-wirten}\\
    
    \textit{durch-} & \textit{durch-wandern}, \textit{durch-denken}\\ 
    
    \textit{ent-} & \textit{ent-kommen}, \textit{ent-kernen}\\
    
    \textit{er-} & \textit{er-leben}, \textit{er-bleichen}\\
    
    \textit{hinter-} & \textit{hinter-gehen}, \textit{hinter-ziehen}\\ 

    \textit{miss-} & \textit{miss-brauchen}, \textit{miss-fallen}\\
    
    \textit{über-} & \textit{über-treiben}, \textit{über-wintern}\\ 

    \textit{um-} & \textit{um-armen}, \textit{um-gittern}\\ 

    \textit{unter-} & \textit{unter-stützen}, \textit{unter-treiben}\\ 

    \textit{ver-} & \textit{ver-stehen}, \textit{ver-schrotten}\\
    
    \textit{zer-} & \textit{zer-brechen}, \textit{zer-fleischen}\\
    
    \lspbottomrule 
  \end{tabular}
  \caption{Wichtige Verbpräfixe}
  \label{tab:Verbpräfixe}
\end{table}

Verbpräfixe verbinden sich grundsätzlich mit Verbstämmen. Aber sie können auch, wie aus den Beispielen ersichtlich wird, aus Adjektiv- und Substantivstämmen Verben ableiten.

Derivationsaffixe werden auch als gebundene\is{Morphem!gebunden} Morpheme bezeichnet, weil sie erstens nur in Anbindung an eine Stammform vorkommen und zweitens als nicht weiter zerlegbare Einheiten aus Form und Bedeutung bzw. Funktion gelten. In diesem Sinne werden sie zwar als gebundene Morpheme von den lexikalischen Stämmen als freie Morpheme abgegrenzt. Als Morpheme sind beide jedoch selbstständige Bestandteile des Lexikons, was wesentliche Unterschiede verwischt. Auf dieses (theoretische) Problem gehen wir in dem folgenden Exkurs~\ref{ExkursMorphemLexikon} ein, bevor es im nächsten Abschnitt~\ref{subsection:VPartikeln} um Verbpartikeln geht.   


\begin{Exkurs}{Gebundene Morpheme als Einheiten des Lexikons?}
\label{ExkursMorphemLexikon}

Die Annahme, dass gebundene Morpheme\is{Morphem} genauso wie lexikalische Morpheme selbstständige Einheiten des Lexikons sind, ist problematisch.
Denn Affixe wie \textit{-nis} können nur zusammen mit lexikalischen Stämmen wie \textit{geheim} in \textit{Geheimnis} gebraucht werden. Sie dienen zur Bildung bestimmter Wortarten, \textit{-nis} zur Bildung von Substantiven, meist von Neutra. Aber sie sind selbst keine Vertreter dieser Wortart, d.\,h. \textit{-nis} ist kein Neutrum. Bedeutungstragend ist das Produkt einer Derivation wie \textit{Geheimnis}. Aber welche Bedeutung sollte \textit{-nis} zukommen? Bestimmte Formen lexikalischer Morpheme wie \textit{geheim} oder \textit{Geheimnis} werden auf syntaktischer Ebene mit anderen Wortformen kombiniert. Affixe treten auf syntaktischer Ebene nicht selbstständig in Erscheinung, weil sie Bestandteile von Wörtern sind. Sie dienen bzw. dienten der Ableitung neuer Wörter aus bestehenden Wörtern. 

Genauso wenig wie \textit{-nis} ein Neutrum ist, sind \textit{-ig} oder \textit{-bar} Adjektive. Sie operieren über substantivischen und verbalen Stammformen, indem sie diese zu Adjektiven wie \textit{steinig} oder \textit{essbar} ableiten. Diesen Ableitungen kommt eine Bedeutung zu. Aus der Differenz zwischen den Wörtern \textit{Stein} und \textit{steinig}, \textit{Staub} und \textit{staubig} sowie zwischen \textit{essen} und \textit{essbar}, \textit{trinken} und \textit{trinkbar} kann von den jeweiligen Basen und der konkreten Bedeutung der Ableitungen abstrahiert werden. Solche Abstraktionen kann man sich als ein Schema vorstellen, das über eine Leerstelle X und ein Affix verfügt. Die Bedeutung von [X-bar]\textsubscript{Adjektiv} trägt das Schema und nicht das Affix allein.

\citet[12]{Motsch2004}: schreibt deshalb: \gqq{Affixe sind [\ldots{}] keine Lexikoneinheiten mit einer eigenen semantischen Repräsentation, sondern lediglich Indikatoren für semantische Muster}.\footnote{Diese Sicht geht auf \citet[47]{Paul1920Wb} zurück.}

Aus diesem Blickwinkel sind Derivationsaffixe Formsignale von Bildungsmustern mit Leerstellen, die durch Regeln von bestimmten Stammformen besetzt werden. Nur selten verselbständigt sich ein Derivationsaffix zu einer lexikalischen Stammform. Dies ist \zB bei \textit{-ismus} in \textit{Sozialismus} oder \textit{Kapitalismus} der Fall. Das Suffix hat sich als (\textit{der}) \textit{Ismus} mit der Bedeutung \gq{bloße Theorie} oder \gq{Gedankenmodell} zu einem Wort entwickelt.

Dem Prozess der Verselbständigung entgegengesetzt ist der Verlust der Affixhaftigkeit. Dies ist \zB bei dem Suffix \textit{-t} in \textit{Fahrt} und \textit{Naht} der Fall. Nur der historische Blick gewährt Einblick in einen alten regulären deverbalen Ableitungsprozess mittels \textit{-t}. \citet[228]{Eisenberg2020_Wort} bezeichnet solche heute nicht mehr aktiven Einheiten als \gqq{morphologischen Rest}.    

\end{Exkurs}



\subsection{Verbpartikeln}\label{subsection:VPartikeln}

Im verbalen Bereich spielen neben den Verbpräfixen die\is{Verbpartikel} Verbpartikeln eine zentrale Rolle. Verbpartikeln tragen im Gegensatz zu Präfixen den Hauptakzent: \textit{EINkaufen} versus \textit{verKAUFen}. Wie bereits erwähnt, bilden sie in Verberst- und Verbzweitsätzen die rechte Klammer, stehen also von ihrer Basis getrennt. Die wichtigsten Verbpartikeln werden in der Tabelle~\ref{tab:Verbpartikeln} aufgeführt. Zur Kennzeichnung der klammerbildenden Funktion geben wir als Beispiele nicht die infiniten Verbformen, sondern die klammerbildende flektierte 3. Person Singular an.

\begin{table}[!htbp]
  \centering
  \begin{tabular}{l l}
    \lsptoprule
    Verbpartikeln & Beispiele \\
    \hline
    \textit{ab} & \textit{fährt \ldots{} ab}, \textit{stumpft \ldots{} ab}\\
    
    \textit{an} & \textit{fängt \ldots{} an}, \textit{kettet \ldots{}\ldots{} an}\\
    
    \textit{auf} & \textit{hebt \ldots{} auf}, \textit{tischt \ldots{} auf}\\

    \textit{aus} & \textit{wandert \ldots{} aus}, \textit{nüchtert \ldots{} aus}\\

    \textit{durch} & \textit{feiert \ldots{} durch}, \textit{macht \ldots{} durch}\\ 
    
    \textit{ein} & \textit{kauft \ldots{} ein}, \textit{schüchtert \ldots{} ein}\\
    
    \textit{mit} & \textit{kommt \ldots{} mit}, \textit{bringt \ldots{} mit}\\ 
    
    \textit{nach} & \textit{fragt \ldots{} nach}, \textit{plappert \ldots{} nach}\\ 

    \textit{über} & \textit{wirft \ldots{} über}, \textit{kocht \ldots{} über}\\ 

    \textit{um} & \textit{rührt \ldots{} um}, \textit{zieht \ldots{} um}\\ 

    \textit{vor} & \textit{macht \ldots{} vor}, \textit{sagt \ldots{} vor}\\
    
    \textit{zu} & \textit{klappt \ldots{} zu}, \textit{zahlt \ldots{} zu}\\
    
    \lspbottomrule 
  \end{tabular}
  \caption{Wichtige Verbpartikeln}
  \label{tab:Verbpartikeln}
\end{table}


Verbpartikeln haben sich aus verschiedenen Wortarten entwickelt. Ihr Prototyp sind heute gleichlautende Präpositionen. Im Unterschied zu den Präpositionen haben die Verbpartikeln jedoch ihre Kasusrektion verloren.

Auf Partikelverben als Schnittstelle zwischen Wortbildung und Syntax haben wir bereits hingewiesen. Wichtige Aspekte der Bildung und Analyse von Partikelverben werden im Abschnitt \ref{section:Vpt Bildung} dargestellt. Vorweggenommen sei hier erstens, dass u.\,a. \textit{durch}, \textit{über} und \textit{um} auch Präfixe sein können (vgl. Tabelle \ref{tab:Verbpräfixe}), und zweitens, dass wir Adjektive wie \textit{krank} in \textit{krankschreiben} und Adverbien wie \textit{empor} in \textit{emporsteigen} nicht zu den Verbpartikeln zählen. Hier handelt es sich um originär syntaktische Strukturen. 



\subsection{Konfixe}\label{subsection:Konfixe}
    
Eine\is{Konfix} Zwischenstellung zwischen gebundenen Affixen und lexikalischen Stämmen nehmen meist aus dem Englischen, Lateinischen oder Griechischen entlehnte Morpheme wie \textit{bio} in \textit{Bio-einkaufserlebnis}, \textit{cyber} in \textit{Cyber-angriff} oder \textit{homo} in \textit{homo-phob} ein. Sie werden Konfixe genannt.\footnote{Den Begriff hat \citet[49]{Schmidt1987} eingeführt, der Konfixe als Affixe bestimmt, welche jedoch nicht der Derivation, sondern der Komposition dienen. Der Begriff und das Konzept sind in Wortbildungslehren üblich geworden. Eine Übersicht bietet \citet{Donalies2009Konfix}.} Konfixe werden trotz einer lexikalischen Bedeutung nicht frei, sondern nur gebunden gebraucht. Häufig können sie keiner Wortart zugeordnet werden und sind folglich alleine nicht flexionsfähig. Das haben sie mit Derivationsaffixen gemeinsam.

Konfixe wie\is{Konfix!versus Affix} \textit{bio} und \textit{homo} können mit anderen Konfixen wie \textit{top} und \textit{phob} zu \textit{Bio-top} und \textit{homo-phob} kombiniert werden. Das unterscheidet sie grundlegend von Affixen, die nicht miteinander kombiniert werden können. Manche Konfixe können wie Stammformen sowohl als Erstglied als auch als Zweitglied so wie \textit{phil}(\textit{o}) in \textit{Philo-soph} und \textit{biblio-phil} gebraucht werden.\footnote{\citet[125]{MüllerPO2000} grenzt die Konfixkategorie auf solche positionsfreien Morpheme ein. Positionsfeste Einheiten zählt er zu den Derivationsaffixen. Der Preis dieser Klassifizierung ist allerdings hoch, da somit Affixe wie \textit{auto-} mit Affixen wie \textit{-mat} kombiniert werden können, was wir im Abschnitt \ref{subsection:Derivationsaffixe} ausgeschlossen haben.}

Wie Stammformen können viele Konfixe auch mit heimischen Stammformen wie in \textit{Homo-ehe} oder \textit{Bio-einkauf} kombiniert werden. 

Konfixe können wie Stammformen und wie Suffixe die Wortart der Bildung festlegen: so ist \textit{Philo-soph} ein maskulines Substantiv und \textit{biblio-phil} ein Adjektiv.

Anders als typische Suffixe und genauso wie Stammformen verfügen Konfixe über lexikalische Bedeutung. Durch das Konfixkompositum \textit{Velo-drom} wird nicht die Bedeutung von \textit{velo} (\gq{Fahrrad}) abgewandelt, sondern die von \textit{drom} auf eine Bahn für Fahrräder eingeschränkt. Aufgrund ihrer lexikalischen Bedeutung könnten diese Konfixe auch als untypische, da nur gebunden vorkommende Stammformen, die sich bevorzugt mit anderen gebundenen Stammformen verbinden, bezeichnet werden.

Mit \citet[63]{FleischerBarz2012} und \citet[22]{Donalies2005WbÜb} gehen wir davon aus, dass manche Konfixe, wie andere Stammformen auch, mit Derivationsaffixen kombiniert werden können. So werden Konfixe wie \textit{techn}, \textit{therm} und \textit{phob} zu den Adjektiven \textit{techn-isch}, \textit{therm-isch} und \textit{phob-isch} (\gq{eine Phobie betreffend}) abgeleitet. Bei anderen Konfixen wie \textit{phil} oder \textit{drom} ist eine Affixableitung nicht möglich: *\textit{phil-isch}, *\textit{drom-isch}.\footnote{\citet[254]{Eisenberg2020_Wort} setzt einen engen Konfixbegriff an und zählt nur letztere, nicht derivierbare zu den Konfixen. Derivierbare Einheiten bestimmt er als gebundene lexikalische Stämme.}


Konfixe wie \textit{öko} und \textit{bio} werden auch als syntaktische Wörter gebraucht und zwar mit spezialisierter Bedeutung prädikativ wie Adjektive (\ref{ea:öko_Prädikativ}), wie Substantive zur Bezeichnung von Personen (\ref{ex:öko_Person}) und teils wie unzählbare Stoffnamen (\ref{ex:öko_Stoffname}). Zur Bezeichnung von Personen wird das Konfix wie in (\ref{ex:öko_Person}) sogar flektiert.

\ea \label{ea:öko_Stammform}
\ea \label{ea:öko_Prädikativ} Der neue Burger ist öko, bio, artgerecht.\footnote{die tageszeitung, 12.04.2014.}
\ex \label{ex:öko_Person} Und wer fliegt, ist unter Ökos eh bis zum nächsten Sommer unten durch.\footnote{Hamburger Morgenpost, 08.04.2007.}
\ex \label{ex:öko_Stoffname} Eltern müssen für Babybrei kein Öko verwenden.\footnote{Braunschweiger Zeitung, 04.05.2013.}
\z 
\z 

Manche Konfixe dienen wie \textit{-loge} zur Bildung ganzer Reihen, d.\,h. sie kommen mit gleicher Funktion mit verschiedenen Konfixen und Stämmen vor, \zB \textit{Astro-loge}, \textit{Bio-loge}, \textit{Geo-loge}, \textit{Öko-loge} usw. Je mehr ein Konfix zur Bildung solcher Reihen dient, umso mehr ähnelt es einem Derivationsaffix. So könnte ein Konfix wie \textit{-loge} auch als untypisches, da mit anderen Konfixen kombinierbares Derivationsaffix bestimmt werden. Es verbindet sich mit Konfix- und selbst mit heimischen Substantivstämmen wie in \textit{Biero-loge} und \textit{Knasto-loge} zum Ausdruck von Fachleuten. Vergleichbar ist die Leistung von \textit{-loge} mit der des Suffix \textit{-er} in \textit{Gärt-ner} und \textit{Logik-er}.  

Viele Konfixe lateinischen und griechischen Ursprungs weisen neben ihrer gebundenen, nicht flexionsfähigen Stammform auch formal gekennzeichnete Kompositionsstammformen auf. Diese zeichnen sich durch ein \textit{-o} aus, das entweder wie bei \textit{thermo} in \textit{Thermo-jacke} Bestandteil der Entlehnung ist oder aber wie bei \textit{polito} in \textit{Polito-loge} oder \textit{Wahlo} in \textit{Wahlo-mat} einem Fugenelement gleicht. 

Unter einen weiten\is{Konfix} Konfixbegriff fallen auch nicht entlehnte Wortsegmente, die eine lexikalische Bedeutung tragen, aber die heute nicht mehr frei gebraucht werden können wie \textit{zimper} in \textit{zimper-lich} oder \textit{Zimper-lise}, \textit{lotter} in \textit{lotter-haft} und \textit{Lotter-wirtschaft}, \textit{schwieger} in \textit{Schwieger-sohn} oder \textit{stief} in \textit{stief-lich} und \textit{Stief-sohn}.


\section{Komposition}\label{section:Komposition}

Deutsch ist eine besonders kompositionsfreudige Sprache.\footnote{Cf. \citet[2]{Schluecker2012}.} Die Komposition ist eine Wortbildungstechnik auf der Ebene von Stammformen. Ihr Input sind Kompositionsstammformen wie \textit{sonnen-} oder \textit{bade-} und lexikalische Stammformen bzw. Flexionsstammformen wie \textit{tag}. Ihr Output sind wie \textit{Sonnentag} und \textit{Badetag} lexikalische Stammformen bzw. Flexionsstammformen, deren Erstglied den Hauptakzent trägt.

Das Prinzip\is{Rechtsköpfigkeit} der Verbindung ist die Einbettung. Die Kompositionsstammformen werden in die lexikalische Stammform eingebettet. Der einbettende Stamm steht rechts vom eingebetteten Stamm und legt die Wortart des Kompositums fest. Beim Prototyp, dem substantivischen Kompositum, legt der rechts stehende Stamm auch das Genus und die Flexionsklasse fest. \textit{Sonnentag} und \textit{Badetag} sind ungeachtet ihrer Erstglieder maskuline Substantive der starken Flexionsklasse (vgl. Abschnitt~\ref{subsubsection:st_sw_SFL}), weil \textit{Tag} maskulin ist und der starken Flexionsklasse angehört. Den rechtsstehenden Stamm bezeichnet man auch als Kopf, worauf wir im Abschnitt~\ref{subsection:Binarität} genauer eingehen\is{Komposition}. Semantisch legt das Zweitglied oder eben der Kopf den Typ fest. Die Begriffe \textit{Sonnentag} und \textit{Badetag} bezeichnen bestimmte Typen von Tagen.  

\begin{Definition}{Komposition}
    
\noindent Komposition ist ein Wortbildungsprozess, bei dem zwei lexikalische Morpheme oder Morphemketten miteinander verbunden werden. Das Erstglied wird als Kompositionsstammform in das Zweitglied eingebettet. Das Zweitglied ist der semantische und grammatische Kopf des Kompositums.
\end{Definition}


Die Stammformen eines Kompositums können intern wie bei \textit{Sitzungsvertagung} aus \textit{sitzungs-} und \textit{vertagung} komplex sein. Neben substantivischen und verbalen Kompositionsstammformen gibt es als Erstglieder auch Konfixe wie in \textit{Öko-tag}, Adjektive wie in \textit{Faul-tag}, bestimmte Muster mit Präpositionen wie in \textit{Vor-tag}, aber auch Adverbien wie in \textit{Sofort-hilfe}, Pronomen wie in \textit{Wir-gefühl} und im Randbereich Subjunktionen wie in \textit{Dass-Satz}. Möglich sind ebenso Wortgruppen wie \textit{Kino im Freien} als Erstglied:\textit{Kino-im-Freien-Tag}.


\subsection{Prototyp Determinativkompositum und Ableitungen}\label{subsection:DetKomp}

Etablierte Komposita\is{Kompositum} verfügen häufig über eine relativ fixierte, teilweise sogar über eine lexikalisierte Lesart. Relativ fixiert ist die Lesart von \textit{Sonnentag} als Tag, an dem die Sonne scheint. Möglich wäre aber auch ein anderes Verständnis, \zB als Tag, an dem die Sonne angebetet wird. 

Die andere Lesart\is{Lexikalisierung} als \gqq{die Zeit zwischen zwei aufeinanderfolgenden Durchgängen der Sonne durch den Meridian}\footnote{Digitales Wörterbuch der deutschen Sprache, \url{https://www.dwds.de/wb/Sonnentag}, abgerufen am 07.03.2021.} wird als \textit{lexikalisiert} bezeichnet. Denn diese Gesamtbedeutung lässt sich nicht aus den Bestandteilen herleiten.

Die Lesart von \textit{Badetag} ist relativ fixiert, als ein Tag, an dem gebadet wird. Ein \textit{Geburtstag} ist nicht primär der Tag einer Geburt, sondern deren Jahrestag. Auch hier geht die Fixierung der Lesart mit Lexikalisierung einher. In jedem Fall geht es bei \textit{Sonnen-}, \textit{Bade-} und \textit{Geburtstag} um bestimmte Typen von Tagen.

Bei dieser\is{Determinativkompositum} Art von Komposita ist grammatisch nur ausgewiesen, dass das Erstglied \gqq{in irgend welcher Art}, wie \citet[8--9]{Paul1920Wb} schreibt, das Zweitglied näher bestimmt. Deshalb werden sie als \textit{Determinativkomposita} (lat. \textit{determinare} \gq{bestimmen}) bezeichnet. Das bestimmende Erstglied ist das\is{Determinativkompositum!Determinans} \textit{Determinans} oder Bestimmungswort, das bestimmte Zweitglied ist das\is{Determinativkompositum!Determinatum} \textit{Determinatum} oder Grundwort. Das Determinatum bestimmt sowohl semantisch als auch formal die Gesamtbildung.

Bei nicht gebräuchlichen Komposita wie \textit{Hundetag} oder \textit{Ekeltag}, muss die jeweilige Beziehung hergestellt werden. Bei \textit{Hundetag} könnte es sich analog zum \textit{Kindertag} um eine Art Ehrentag für Hunde handeln oder um einen Tag, an dem Hunde frei in den Parks herumlaufen dürfen oder an dem sich um ausgesetzte Hunde gekümmert wird usw. Bei \textit{Ekeltag} ist nicht klar, ob das Determinans verbal oder substantivisch ist. Es könnte um einen Tag gehen, an dem der Ekel bekämpft oder gefeiert wird oder aber um einen Tag, an dem sich analog zu einem \textit{Wohlfühltag} geekelt wird oder vor dem sich geekelt wird usw. Das Assoziationspotential wird durch etablierte Bildungen mit fixierten Lesarten und wie in (\ref{ea:Hundetag}) und (\ref{ex:Ekeltag}) durch den Kontext eingeschränkt.

\ea \label{ea:Kompositumadhoc}
\ea \label{ea:Hundetag} Das Erika-Kino hat im Fernsehen groß angekündigt, den ersten Hundetag in der Geschichte des europäischen Kinos einzuführen. Da dürfen die Vierbeiner dann einmal im Jahr mitkommen.\footnote{Niederösterreichische Nachrichten, 01.03.2012.}
\ex \label{ex:Ekeltag} Die Menschen müssen das Recht haben, findet Dr. Blaul, ohne Ankündigung und ohne Begründung einen Tag ihrem Betrieb fernzubleiben, wenn sie plötzlich Ekel empfinden, sei es über die Firma, über einen Kollegen, über sich selbst oder worüber auch immer. Ein freier Ekeltag soll uns helfen, daß wir uns wieder fangen können, ohne gleich fliehen zu müssen.\footnote{Wilhelm Genazino: \textit{Die Liebesblödigkeit}. München, Wien: Carl Hanser Verlag 2005, S. 33.}
\z 
\z 

Determinativkomposita sind\is{Determinativkompositum!Unterspezifiziertheit} semantisch unterspezifiziert. Selbst lexikalisierte Lesarten können kontextuell überschrieben werden.\footnote{\citet[13]{Schluecker2012} führt das folgende Beispiel für \textit{Blumentopf} auf. Die lexikalisierte Bedeutung ist \gq{Topf, in den Blumen eingepflanzt werden}. Der abweichende Kontext: \textit{Den Brei kochte er in einem kleinen Topf, der mit Blumen bemalt war. Dieser Blumentopf\ldots{}}} Jene Unterspezifiziertheit hat \citet[2]{Heringer1984Fischfrau} in seinem Aufsatz \textit{Sinn aus dem Chaos} am Beispiel der zahlreichen Lesarten des Kompositums \textit{Fischfrau} thematisiert. Gemeint sein kann eine Frau, die Fisch verkauft, die mit einem Fisch zusammen ist, die im Sternbild Fisch geboren ist, die von einem Fisch abstammt usw.

Es wurde oft versucht, die Lesarten im Sinne von Umschreibungen (Paraphrasen) als verschiedene Typen von Komposita zu systematisieren. Hierbei wird die Relation zwischen Determinans und Determinatum beschrieben. \citet[141]{FleischerBarz2012} führen die folgenden, in der Tabelle \ref{tab:DeterminantivkompositaRel} wiedergegebenen Relationen auf.\footnote{\citet[133--139]{OrtnerOrtner1984} ermitteln durch Paraphrasen 47 semantische Beziehungen zwischen Erst- und Zweitgliedern. \citet[7]{Motsch2004} sieht für die Bildung der jeweiligen Typen spezielle semantische Muster vor.}

\begin{table}[!htbp]
  \centering
  \begin{tabular}{l l}
    \lsptoprule
    Relation & Beispiele \\
    \hline
    lokal (Ort) & \textit{Straßenlaterne}, \textit{Seehotel}\\
    
    temporal (Zeit) & \textit{Sommerferien}, \textit{Abendessen}\\
    
    final (Zweck) & \textit{Weinglas}, \textit{Hundenapf}\\

    kausal (Grund) &\textit{Schmerzensschrei}, \textit{Kratzwunde}\\

    komparativ (vergleichend) & \textit{Beifallssturm}, \textit{Patchworkfamilie}\\ 
    
    possessiv (zugehörig) & \textit{Schneckenhaus}, \textit{Gemeindeschule}\\
    
    partitiv (Teil-Ganzes) & \textit{Haustür}, \textit{Henkelkorb}\\ 
    
    Material & \textit{Lederschuh}, \textit{Goldkette}\\ 

    Oberbegriff & \textit{Auswahlverfahren}, \textit{Eichenbaum}\\ 

    \lspbottomrule 
  \end{tabular}
  \caption{Semantische Relationen Determinativkomposita}
  \label{tab:DeterminantivkompositaRel}
\end{table}

Ohne Frage handelt es sich bei den aufgeführten Relationen um musterhafte Beziehungen, wie sie in etablierten Komposita niedergelegt sind. Mitunter ist die Musterhaftigkeit auch strukturell geprägt. So werden Komposita wie \textit{feder-leicht} oder \textit{aal-glatt} mit substantivischem Erst- und adjektivischem Zweitglied präferiert komparativ als \gq{leicht wie eine Feder} bzw. \gq{glatt wie ein Aal} interpretiert und im Sinne von \gq{sehr leicht} bzw. \gq{sehr glatt} gebraucht. 

Aber grammatisch vorgegeben sind die in der Tabelle \ref{tab:DeterminantivkompositaRel} aufgeführten Beziehungen nicht. \citet[3]{Heringer1984Fischfrau} schreibt diesbezüglich: \gqq{Da ist aber nichts, kein Zeichen für diese Relation}. Die morphologische Struktur der Komposita sagt uns nichts über die konkrete Beziehung: \textit{Gold-preis} bezeichnet den Preis von Gold, \textit{Gold-kette} eine Kette aus Gold und \textit{Goldwaage} eine Waage für Gold. \citet[109, 117]{Coseriu1970} zeigt, dass es bei Nicht-Lexikalisierung unser Weltwissen und der sprachliche Kontext sind, die die Interpretation steuern. 

Formal festgelegt\is{Determinativkompositum!Unterspezifiziertheit} ist, dass das Determinatum der formale und semantische Kopf der Bildung ist. Die Bedeutung des Kopfs steht in einer unterspezifizierten Relation zu der Bedeutung des Determinans. Die Bedeutung des Kopfs legt die Klasse, die Kategorie fest. Das Determinans schränkt die Bedeutung des Kopfs auf eine Unterklasse, eine Unterkategorie ein: eine \textit{Haustür} ist eine Unterkategorie von \textit{Tür}, ein \textit{Hundetag} eine Unterkategorie von \textit{Tag}.  

Die wesentliche Restriktion besteht darin, dass die semantischen Bestandteile zu einem einheitlichen Gesamtkonzept miteinander verbunden werden. Dies erlaubt die Benennung einer Klasse von Entitäten.\footnote{\citet[88]{Olsen2001Kopulativ} bezeichnet dies als \gqq{Principle of Ontological Coherence}. Es geht darum, dass Komposita keine Beschreibung von Dingen leisten, sondern wie andere Substantive Klassen von Entitäten bedeuten und zur Bezeichnung eines Mitglieds der Klasse dienen. \citet[271--275]{Buecking2010} zeigt, dass und wie auch neu gebildete Komposita wie \textit{Blautee} als Bezeichnung für Klassen generisch interpretiert werden. Bei \textit{Blautee} handelt es sich um einen Typ von Tee, der mit dem Träger der Eigenschaft \textit{blau} in einer integralen Beziehung steht, \zB als \gq{Tee, der in einer blauen Büchse aufbewahrt wird}.} Je geläufiger, alltäglicher uns die entsprechende Klasse erscheint, umso leichter gelingt die Interpretation auch ohne Kontext.

Durch\is{Determinativkompositum} die Komposition festgelegt ist, dass das Determinans nicht auf einen konkreten Vertreter der bedeuteten Klasse referiert, sondern auf das entsprechende Konzept. Deshalb ist das Determinans normalerweise nicht mehr durch Attribute modifizierbar.\footnote{Bei sogenannten schiefen Attributen wie bei \textit{deutsche Sprachgeschichte} für \gq{Geschichte der deutschen Sprache} beziehen wir das Attribut \textit{deutsche} auf das Determinans \textit{Sprach-}. Bei Interesse verweisen wir auf \citet[178]{Maienborn2020}, die jenen schiefen Bezug als sekundären Effekt analysiert. Auslöser ist die pragmatisch sparsamste Anbindung unterspezifizierter Relationen wie der zwischen einem relationalen Adjektiv wie \textit{deutsch} und einem Bezugswort sowie der zwischen Determinans und Determinatum in Komposita. Regulär vorgesehen ist der Bezug von \textit{deutsche} auf das Kompositum \textit{Sprachgeschichte}, nur ist die Interpretation als \gq{deutsche Geschichte der Sprache} nicht die sparsamste Interpretation, was den Regelverstoß motiviert.} Während durchaus vorstellbar ist, dass sich an einem \textit{Hundetag} um streunende Hunde gekümmert wird, kann dies nicht durch *\textit{streunender Hundetag} ausgedrückt werden.

Die Bildung von Unterkategorien durch die Komposition unterscheidet sich von der syntaktischen Modifikation. Während durch ein Adjektivattribut wie in \textit{der schnelle Zug} ein bestimmter Zug genauer beschrieben wird, handelt es sich bei dem Kompositum \textit{Schnellzug} um eine Unterkategorie von Zug. Deshalb kann ein konkreter Vertreter des Typs Schnellzug durchaus langsam sein, nicht aber -- ohne Umdeutung -- der als schnell beschriebene Zug: \#\textit{der langsame schnelle Zug}.

Neben der klassifizierenden Benennungsfunktion können Komposita aber auch beschreibend gebraucht werden, \zB wenn wir in einem Restaurant eine bestimmte Frau, die gerade Fisch isst, als \textit{Fischfrau} bezeichnen. Das Gleiche geschieht in dem Beleg (\ref{ea:Zwickermann}) durch \textit{Zwickermann}, der zuvor als \textit{ein Mann mit Kneifer} beschrieben wird.

\ea \label{ea:Zwickermann} Und ein blonder Mann mit Kneifer ladet mich ein -- mit Zähnen wie eine Maus und so einem widerlich kleinen Mund, der glänzt feucht und macht den Zwickermann nackend.\footnote{Irmgard Keun: \textit{Das kunstseidene Mädchen}. [1932]. Bastei Lübbe Verlag, 1982, S. 106.}
\z 


\subsubsection{Rektionskompositum}\label{subsubsection:RektKomp}

Wenn wir die Beziehung zwischen Determinans und Determinatum auf verbale Valenzverhältnisse (vgl. Einführung ins Kapitel~\ref{chapter:Valenz}) zurückführen können, dann ist die Interpretation vorgezeichneter. So ist ein \textit{Roman-leser} jemand, der Romane liest und ein \textit{Zigarren-raucher} jemand, der Zigarren raucht. Das Determinans wird sinnvoll auf die Rolle der zweiten Ergänzung des Verbs bezogen, das dem Determinatum zugrunde liegt. Dieser Typ von Determinantivkomposita wird als\is{Rektionskompositum} \textit{Rektionskompositum} bezeichnet. Denn in der Paraphrase regiert das Verb das Determinans als Akkusativobjekt (vgl. Abschnitt \ref{subsection:Akkusativobj}). 

Die präferierte Interpretation ist nicht auf ursprünglich als Akkusativobjekte angelegte Ergänzungen beschränkt. So kann eine sinnvolle Relation des Determinans zum Determinatum auch auf eine ursprünglich kategorial regierte erste Verbergänzung im Nominativ bezogen werden wie bei \textit{Zug-abfahrt} oder \textit{Preis-anstieg}. Ebenso kann die Beziehung des Determinans zum Determinatum auf Dativ- (vgl. Abschnitt~\ref{subsection:Datobj}), auf Präpositionalobjekte (vgl. Abschnitt~\ref{subsection:Pobj}) und auf Direktiva (vgl. Abschnitt~\ref{subsection:Direktivum}) zurückgeführt werden: \textit{Helfers-helfer} für jemanden, der Helfern hilft, \textit{Literatur-beschäftigung} als Name für die Tätigkeit, wenn man sich mit Literatur beschäftigt, und \textit{Berg-steiger} für jemanden, der auf Berge steigt.    

Die Rückführung der Beziehung zwischen Determinans und Determinatum auf verbale Rektionsverhältnisse bietet häufig eine sinnvolle Interpretationsbasis. Sie ist aber nicht zwingend. Ohne Kontext wissen wir nicht, ob \textit{Pistolen-räuber} einen Typ von Räuber benennt, der Pistolen raubt (regiertes Akkusativobjekt) oder der mit Pistolen (etwas anderes) raubt (nicht regiertes Adverbial). Folglich ist es nicht die verbale Rektion, sondern es sind sämtliche, auch adverbial ausgedrückte Umstände und Spezifizierungen eines Verbszenarios, die ein starkes Interpretationspotential für Determinativkomposita bieten. 

Im weiteren\is{Rektionskompositum} Sinne zählt \citet[71]{Olsen1985} auch Komposita wie \textit{Juwelen-dieb} und \textit{Professoren-sohn} zu den Rektionskomposita, da das Determinans eine Leerstelle des Determinatums sättigt. Denn ein Dieb ist per definitionem Dieb von etwas, so wie ein Sohn immer Sohn von jemandem ist. In diesem weiteren Sinne der Leerstellensättigung können auch adjektivische Komposita wie \textit{anpassungs-fähig}, \textit{arbeits-willig} und \textit{ausgeh-bereit} zu den Rektionskomposita gezählt werden. Denn bei der syntaktischen Paraphrase sind die Erstglieder präpositionale Ergänzungen des adjektivischen Kopfs: \textit{anpassungsfähig} paraphrasiert als \textit{fähig zur Anpassung}, \textit{arbeitswillig} als \textit{willig zur Arbeit} und \textit{ausgehbereit} als \textit{bereit zum Ausgehen}. Aber bei der Komposition geschieht jene Sättigung gerade nicht in Form einer Präpositionalphrase, sondern in Form von Kompositionstammformen. Das wird durch das Fugenelement \textit{-s} bei \textit{arbeits} in \textit{arbeitswillig} deutlich. Diese Form hat nichts mit den Flexionsformen des Substantivs \textit{Arbeit} zu tun und auch nichts mit der Rektion von \textit{willig}. 

Der auf ursprüngliche Akkusativobjekte bezogene und auf semantische Leerstellen ausgedehnte Begriff der Rektionskomposita ist kein wortbildender Untertyp der Komposition. Hinter dem etablierten Begriff verbergen sich sinnstiftende Interpretationsroutinen von Determinativkomposita.


\subsubsection{Possessivkompositum}\label{subsubsection:PossKomp}

Als besondere Determinativkomposita sehen wir auch Bildungen wie \textit{Rot-kehlchen} oder \textit{Lang-finger}. Denn auch hier determiniert das Erstglied als Determinans das Zweitglied als Determinatum, das die Wortart und Flexionsklasse des Kompositums festlegt. Aber die Gesamtbildung wird metaphorisch und metonymisch auf etwas Drittes bezogen. Metonymisch, d.\,h. auf einer Teil-Ganzes-Beziehung beruhend, ist der Bezug, wenn dieses Dritte Träger des vom Kompositum Bezeichneten ist. So bezeichnen wir mit \textit{Rotkehlchen} eine ganz bestimmte Vogelart, die sich u.\,a. durch ein rotes Kehlchen auszeichnet. Ein \textit{Langfinger} ist nur im übertragenen, das heißt im metaphorischen Sinne der Träger langer Finger. Es geht um die Bezeichnung eines Diebes, weil er zum Diebstahl lange Finger macht. Mit \textit{Dumm-kopf} bezeichnen wir keine Art von Kopf, sondern einen Menschen, dessen Kopf, stellvertretend fürs Denken, dumm ist. Im weitesten Sinne besitzt oder verfügt das Bezeichnete über das vom Kompositum primär Bedeutete. Deshalb bezeichnet man diese Komposita auch als\is{Possessivkompositum} Possessivkomposita.

Da das Determinatum als formaler Kopf nicht mit dem Bezeichneten übereinstimmt, das Bezeichnete also \textit{außerhalb} liegt, nennt man solche Bildungen seit \citet[235]{Bloomfield1935} auch \textit{exo}zentrisch. Letztlich handelt es sich um eine besondere Art der Lexikalisierung von Determinativkomposita.


\subsubsection{Kopulativkompositum}\label{subsubsection:KopKomp}

Komposita wie \textit{Dichter-komponist}, \textit{Hemd-bluse}, \textit{Bett-sofa} oder \textit{nass-kalt} können wir als gewöhnliche Determinativkomposita \zB wie folgt interpretieren: \textit{Dichterkomponist} als Komponist, der für Dichter komponiert, \textit{Hemdbluse} als Bluse, die einem Hemd ähnelt, \textit{nasskalt} als kalt, das durch Nässe stark gefühlt wird und \textit{Bettsofa} als Sofa, das zu oder aus einem Bett gemacht werden kann.

Unter Beibehalt der morphologischen Rechtsköpfigkeit verstehen wir diese Komposita aber auch koordinativ: also \textit{Dichterkomponist} als Dichter und Komponist in einem, \textit{nasskalt} als nass und zugleich kalt usw. Komposita mit gleichrangig interpretierbaren Bestandteilen bezeichnet man als \textit{Koordinativ-} oder als\is{Kopulativkompositum} \textit{Kopulativkomposita}. Selten ist die koordinative gegenüber der determinativen Lesart formal durch die Abwesenheit eines sonst vorhandenen Fugenelements gekennzeichnet: so bei \textit{Herr-gott} gegenüber \textit{Herren-hut} oder bei \textit{Prinz-regent} gegenüber \textit{Prinzen-erzieher}.

\newpage
Die Abschwächung des Determinativverhältnisses liegt darin begründet, dass beide Stämme außerhalb des Kompositums denselben direkten Oberbegriff haben. Die Begriffe \textit{Dichter} und \textit{Komponist} fallen unter den direkten Oberbegriff \gq{Künstler}, \textit{nass} und \textit{kalt} unter den direkten Oberbegriff \gq{Lufteigenschaft} sowie \textit{Bett} und \textit{Sofa} unter den direkten Oberbegriff \gq{Liegestatt}.\footnote{Ganz anders sieht es aus, wenn wie bei \textit{Beeren-frucht} oder \textit{Eichen-baum} der Oberbegriff selbst das Determinatum ist. Dann liegt ein klares, nämlich verdeutlichendes oder explikatives Determinantionsverhältnis vor. Dieses nutzen wir \zB dann, wenn das Erstglied ungebräuchlich geworden ist wie bei \textit{Ren-tier} oder \textit{Lind-wurm} (\textit{lint} althochdeutsch \gq{Schlange}) oder wenn wir bewusst eine von gewissen Festlegungen abweichende Klassifizierung vornehmen wie bei \textit{Wal-fisch}.} Diese begriffliche Gleichrangigkeit kann auf das Kompositum übertragen werden. Die unterspezifizierte Relation zwischen Determinans und Determinatum wird koordinativ interpretiert. Die Voraussetzung hierfür ist genauso wie bei den Determinantivkomposita, dass sie \textit{ein} kohärentes Konzept bilden. Das heißt, das Erst- und das Zweitglied sind trotz unterschiedlicher Bedeutung bezeichnungsidentisch. Ihre Bedeutung trifft zu gleichen Teilen auf \textit{einen} Referenten zu.\footnote{Bei Interesse verweisen wir auf die Untersuchung von \citet{Olsen2001Kopulativ}.}


Eine völlige Aufhebung des Determinationsverhältnisses und somit pure Reihung liegt bei Eigennamen wie \textit{Mecklenburg-Vorpommern}, bei von Eigennamen abgeleiteten Adjektiven wie \textit{deutsch-argentinisch} sowie bei Farbadjektiven wie \textit{schwarz-rot-gelb} vor. Diese lassen sich nicht mehr auf der Schablone von Determinativkomposita bestimmen. Während \textit{blaugrau} durchaus als eine Art blaues Grau verstanden werden kann, ist dies bei \textit{schwarz-rot-gelb} nicht möglich. Hier handelt es sich ursprünglich einfach um das Verschmelzen von zwei syntaktisch benachbarten Wörtern zu einem, was man als \textit{Univerbierung} bezeichnet.  


\subsection{Fugenelemente}\label{subsection:Fuge}

Wie wir im Abschnitt~\ref{subsection:Stämme} gezeigt haben, können Kompositionsstammformen wie \textit{fugen-} in \textit{Fugen-element} formal durch ein\is{Fugenelement} Fugenelement ausgezeichnet sein. 

Historisch lassen sich Fugenelemente auf Genitivflexive beziehen, heute auf Genitivflexive im Singular (\ref{ea:FE_Gen_Sgl_Syn}) und im Plural (\ref{ex:FE_Gen_Pl_Syn}). Entsprechende Konstruktionen zeichneten sich durch die Voranstellung der Genitivattribute aus.

\ea \label{ea:FE_Genitiv_Synt}
\ea \label{ea:FE_Gen_Sgl_Syn} des Dieb\textit{es} Gut
\ex \label{ex:FE_Gen_Pl_Syn} der Mensch\textit{en} Wille
\z 
\z 

\newpage
Diese syntaktischen Konstruktionen (\ref{ea:FE_Genitiv_Synt}) wurden in Analogie zu älteren, schon bestehenden Komposita\footnote{\citet[408, 597--600]{Grimm1826Gr2} schreibt von \textit{eigentlichen Komposita}, welche aus dem Urgermanischen ererbt sind. Es handelt sich um flexionslose Stammverbindungen, in denen das Erstglied ein stammbildendes Suffix als nicht flexivischen \gqq{compositionsvocal} trägt, \zB \textit{tagawerch}. Seit dem Althochdeutschen entfällt mit dem stammbildenden Suffix auch der Kompositionsvokal oder bleibt relikthaft als reduziertes \textit{e} wie in \textit{Tage-werk}. Die später mit ursprünglich flexivischem Element gebildeten Komposita wie \textit{Tages-licht} bezeichnet Grimm als \textit{uneigentliche} Komposita, da sie auf Wortgruppen wie \textit{des Tages Licht} zurückgehen. \citet[6]{Paul1920Wb} zweifelt einen qualitativen Unterschied an, da auch die älteren Komposita aus syntaktischen Konstruktionen mit anschließender Analogiebildung hervorgegangen sein können.} reanalysiert (\ref{ea:FE_Genitiv}).

\ea \label{ea:FE_Genitiv}
\ea \label{ea:FE_Gen_Sgl} das Dieb\textit{es}gut
\ex \label{ex:FE_Gen_Pl} der Mensch\textit{en}wille
\z 
\z 

Die ursprünglichen Flexive werden als Fugenelement reanalysiert, das heißt als grammatisches Signal für die wortbildende Verbindung zweier Stammformen. Nur die Pluralkennzeichnung kann teilweise noch semantisch motiviert sein.\footnote{Mitunter haben die Fugenelemente \textit{-er} und (\textit{-e})\textit{n} die gleiche Herkunft wie die entsprechenden Pluralsuffixe, sind aber nicht aus ihnen abgeleitet. Im Abschnitt~\ref{subsubsection:Numerus} haben wir die Herkunft der heutigen Pluralflexive \textit{-er} (\textit{Lämmer}) und \textit{-en} (\textit{Augen}) bei einigen Neutra angesprochen. Es handelt sich um ursprünglich stammbildende Suffixe, die sowohl zu Pluralflexiven als auch zu Fugenelementen reanalysiert wurden. Cf. \citet[175--176]{Wegener2005}. Nach \citet[6, 10--12]{Paul1920Wb} gehen die Fugenelemente auf Genitivflexive zurück, die erst später mit Pluralflexiven zusammengefallen sind.}

Fugenelemente, die auf\is{Fugenelement!paradigmatisch} Flexionsformen bezogen werden können, heißen \textit{paradigmatisch}, eben weil sie einer Form aus dem jeweiligen Flexionsparadigma entsprechen. So steht die Kompositionsstammform \textit{diebes-} in \textit{Diebes-gut} in paradigmatischer Beziehung zur Genitivform \textit{Diebes} und die Kompositionsstammform \textit{häuser-} in \textit{Häuser-kampf} steht in paradigmatischer Beziehung zu den Pluralformen (außer Dativ) \textit{Häuser}. 

Manche Fugenelemente\is{Fugenelement!unparadigmatisch} wie \textit{-en} in \textit{Schwanen-hals} oder \textit{Hahnen-kamm} entsprechen heute nicht mehr einer Flexionsform, zeigen aber die ursprüngliche Zugehörigkeit dieser Substantive zur Klasse der schwachen Maskulina (vgl. Abschnitt~\ref{subsubsection:st_sw_SFL}) an. Insbesondere das Fugenelement \textit{-s} wurde in Analogie zu anderen Komposita auf feminine Stämme übertragen, die keine entsprechende Flexionsform aufweisen: \zB bei \textit{Liebes-brief} oder \textit{Geburts-tag} im Gegensatz zu den Genitivformen \textit{Liebe} und \textit{Geburt}. Diese Fugenelemente bezeichnet man als unparadigmatisch.

Auch paradigmatische\is{Fugenelement!versus Flexion} Fugenelemente setzen wir nicht mit den Flexiven gleich. Erstens besteht bei diesen Fugenelementen nicht die für die Flexion typische Variation zwischen \textit{-s} und \textit{-es} wie bei \textit{das Gut des Dieb-}(\textit{e})\textit{s} im Gegensatz zu *\textit{Diebs-gut}. 

Zweitens lässt sich eine semantische Genitivfunktion nicht durchgängig auf die Kompositionsstammform mit paradigmatischem Fugenelement beziehen. So ist \zB die lexikalisierte Bedeutung von \textit{Gottes-furcht} nicht zu verstehen als \textit{Gottes Furcht} oder \textit{Furcht Gottes}, sondern als \gq{Furcht vor Gott}. 

Drittens muss eine formal mit dem Plural identische Kompositionsstammform nicht semantisch motiviert sein. Bei \textit{Bücher-regal} kann von einer semantischen Pluralmotivation als Regal für Bücher ausgegangen werden. Bei \textit{Rinder-braten} oder \textit{Kinder-wagen} ist das nicht der Fall. Umgekehrt sollte in einem \textit{Feder-bett} mehr als eine Feder sein. 

Und viertens ist die \textit{s}-Fuge trotz semantischer Pluralmotivation dann ausgeschlossen, wenn sie mit der flexivischen Pluralform identisch ist: *\textit{Autos-haus} versus \textit{Bücher-geschäft}. Die \textit{s}-Fuge ist das einzige noch produktive unparadigmatische Fugenelement, aber dieses darf gerade nicht mit einer entsprechenden flexivischen Pluralform kollidieren.

Selten dienen Fugenelemente der semantischen Differenzierung des Erstglieds. Während es in \textit{Land-leben} um \textit{Land} als dörfliche Gegend geht, handelt es sich bei \textit{Landes-präsident} um \textit{Land} als Staat. Bei \textit{Länder-spiel} könnte der Plural das Fugenelement motivieren. In \textit{Herzens-wunsch} geht es anders als bei \textit{Herz-medikament} um ein metaphorisches Verständnis von \textit{Herz}. 

Bei manchen Komposita gibt es Varianten bezüglich des Fugenelements. Teilweise ist diese Variation räumlich bedingt. So wird im nord- und mitteldeutschen Sprachgebiet eher vom \textit{Rinder-} und \textit{Schweine-braten} gesprochen, im süddeutschen Raum eher vom \textit{Rinds-} und \textit{Schweins-braten}. Mitunter stehen sich ältere, stilistisch als gehoben empfundene Formen und jüngere Bildungen gegenüber: \textit{Monden-schein} vs. \textit{Mond-schein} oder \textit{Meeres-rauschen} vs. \textit{Meer-rauschen}. Häufig betrifft die Variation entlehnte Wörter und beruht auf der An- und Abwesenheit der \textit{s}-Fuge: \textit{Interessen}(\textit{s})\textit{-vertretung}, \textit{Advent}(\textit{s})\textit{-kerze}, \textit{Subjekt}(\textit{s})\textit{-pronomen}, \textit{Respekt}(\textit{s})\textit{-person}, \textit{Praktikum}(\textit{s})\textit{-bericht}.\footnote{Eine Übersicht über entsprechende Schwankungsfälle geben \citet[217--219]{NueblingSzczep2009}.}

Den Fugenelementen\is{Fugenelement!Motivation} lässt sich, wie \citet[525, 532]{Fuhrhop1996} schreibt, \gqq{keine eindeutige Funktion} und \gqq{keine klare Systematik} zuschreiben, denn die Mehrheit der Kompositionsstammformen weist kein Fugenelement auf.\footnote{\citet[105]{Kuerschner2003} untersucht ein Zeitungskorpus und ermittelt, dass 58\% der Komposita kein Fugenelement aufweisen. In Texten der erzählenden Prosa sind nach \citet[54]{Ortner_Komposita1991} 72,8\% der Komposita ohne Fuge gebildet.} Praktisch heißt das, die Kompositionsstammformen mit Fugenelement müssen gelernt werden. Dennoch kann die Anwesenheit bestimmter Fugenelemente systematisch durch verschiedene lautliche (Silbigkeit/Lautkonstanz), morphologische (Flexionsklasse/Komplexität), semantische (Plural) und lexikalische (Wortart) Eigenschaften motiviert sein. Solche Regularitäten und Tendenzen stellen wir bezüglich einzelner Fugenelemente in den folgenden Abschnitten dar. Insbesondere aus Sicht des Fremd- oder Zweitsprachenlernens kann ein systematischer Zugang auch aus sprachpraktischen Gründen sinnvoll sein.


\subsubsection{\textit{s}-Fuge}\label{subsubsection:s-Fuge}

Abgeleitete Substantive auf \textit{-ling}, \textit{-ung}, \textit{-heit},(\textit{-ig})\textit{keit} und \textit{-schaft} bilden die Kompositionsstammform immer mit der \textit{s}-Fuge: \textit{Frühlings-gefühl}, \textit{Versammlungs-freiheit}, \textit{Freiheits-wille}, \textit{Müdigkeits-attake}, \textit{Verwandtschafts-besuch}. \citet[210--211]{Fuhrhop2000} sieht die Fuge als Öffner von suffigierten Stämmen für weitere, dem Suffix folgenden Wortbildungsprozesse. Dies gilt insbesondere für die Derivation: \textit{frühlings-haft}, \textit{versammlungs-mäßig}. 

Auch von Partikelverben abgeleitete Substantive bilden ihre Kompositionsstammform mit der \textit{s}-Fuge: \textit{Absichts-erklärung} versus \textit{Sicht-achse}, \textit{Abfahrts-zeit} versus \textit{Fahrt-bericht}, \textit{Einkaufs-erlebnis} versus \textit{Kauf-erlebnis}. Die \textit{s}-Fuge kann als Grenzsignal zwischen morphologisch komplexen Erst- und Zweitgliedern verstanden werden. So besteht \textit{Abfahrtszeit} nicht aus \textit{ab} und \textit{fahrtzeit}, sondern aus \textit{abfahrt} und \textit{zeit}. Wenn das Erstglied sowohl verbal als auch substantivisch sein kann, signalisiert die \textit{s}-Fuge wie bei \textit{Einkaufserlebnis} oder \textit{Versuchsperson} den substantivischen Charakter. 

Nach \citet[217]{NueblingSzczep2010} bedingt nicht allein die morphologische Komplexität die Fuge, sondern die lautliche Qualität des Erstglieds. Je besser es der prototypischen deutschen Wortform mit betonter und unbetonter Silbe (vgl. Abschnitt~\ref{subsection:Betonung}) wie bei \textit{ANruf} und \textit{ABfall} entspricht, umso seltener wird eine Fuge gebraucht. Je stärker das Erstglied vom Prototyp abweicht, zum Beispiel durch die Abfolge einer unbetonten und einer betonten Silbe wie bei \textit{BeRUF} und \textit{VerFALL}, umso eher wird die \textit{s}-Fuge gebraucht: \textit{Berufs-wunsch} vs. \textit{Anruf-beantworter}, \textit{Verfalls-datum} vs. \textit{Abfall-verwertung}. 

Wird ein substantivierter Infinitiv wie (\textit{das}) \textit{Essen}, \textit{Leben} oder \textit{Wissen} als Erstglied gebraucht, so trägt die Kompositionsstammform das Fugenelement \textit{-s} wie in \textit{Essens-marke}, \textit{Lebens-hunger} und \textit{Wissens-durst}. Das Fugenelement verdeutlicht hier den substantivischen Charakter der Infinitive, der ansonsten nur syntaktisch ausmachbar ist wie in \textit{er möchte das Wissen gebrauchen} versus \textit{er möchte das wissen}.


\subsubsection{(\textit{e})\textit{n}-Fuge}\label{subsubsection:n-Fuge}

Bei Feminina, die auf -\textit{e} enden, ist eine \textit{n}-Fuge wie bei \textit{Blumen-beet}, \textit{Straßen-laterne} oder \textit{Küchen-tisch} erwartbar. Die \textit{n}-Fuge schließt die offene Silbe (vgl. Abschnitt~\ref{subsec:Silbenkonstituenten}). Bei \textit{Blumen-beet} oder \textit{Bienen-honig} ist eine Pluralmotivation möglich, bei \textit{Küchen-tisch} nicht. 

Bei den schwachen Maskulina wie \textit{Affe} (vgl. Abschnitt~\ref{subsubsection:st_sw_SFL}) ist -(\textit{e})\textit{n} als Fugenelement vorhersagbar. Deshalb heißt es \textit{Affen-art}, \textit{Hasen-pfote} oder \textit{Menschen-menge}. Hier ist die Fuge morphologisch motiviert, da sie paradigmatisch ist.



\subsubsection{\textit{e}-Fuge}\label{subsubsection:e-Fuge}

Verbstämme bilden Kompositionsstammformen ohne Fugenelement oder aber mit dem Fugenelement \textit{-e}. Dieses Fugenelement ist phonologisch motiviert, da es primär bei Verbstämmen steht, die auf einen stimmhaften Verschlusslaut (vgl. Abschnitt~\ref{subsection:Plosive} unter dem Stichwort \textit{Plosive}) auslauten. So heißt es im bundesdeutschen Standard \textit{Trage-tasche}, \textit{Schiebe-tür} und \textit{Bade-wasser} aber \textit{Wasch-tasche}, \textit{Dreh-tür} und \textit{Trink-wasser}. Das Fugenelement verhindert die Auslautverhärtung und sichert somit die lautliche Stammkonstanz. Allerdings gibt es Abweichungen in beide Richtungen. Trotz stimmhaftem Verschlusslaut fehlt die Fuge bei \textit{Schlag-stock} oder \textit{Bind-faden}, im Schweizerdeutschen auch \textit{Bad-zimmer} und \textit{Trag-tasche}. Ohne stimmhaften Verschlusslaut ist die \textit{e}-Fuge in \textit{Warte-zimmer} oder \textit{Halte-verbot} vorhanden. Ohne Fugenelement ist teilweise nicht zu klären, ob die Kompositionsstammform substantivisch oder verbal ist, \zB bei \textit{Streik-posten} oder \textit{Trauer-gesellschaft}. Hierauf kommen wir im Abschnitt~\ref{subsection:Konversion} unter dem Stichwort \textit{Konversion} zurück. 

Bei Bildungen wie \textit{Pflege-fall} oder \textit{Folge-kosten} kann das Erstglied analog zu \textit{Bade-wasser} und \textit{Zeige-finger} als verbale Kompositionsstammform mit phonologisch motivierter \textit{e}-Fuge analysiert werden. Möglich ist aber auch, wie \citet[543]{Fuhrhop1996} zeigt, der Bezug auf die durch das Suffix -\textit{e} abgeleiteten Substantive.


Bei Substantiven ist die \textit{e}-Fuge oft durch die gleichlautende Pluralform motiviert: \textit{Städte-tag} versus \textit{Stadt-mauer} oder \textit{Tage-buch} versus \textit{Tag-traum}, \textit{Ärzte-kammer} versus \textit{Arzt-praxis}. Dass es sich allenthalben um eine Motivation handelt, aber keine durchgängig semantisch motivierten Pluralformen, wird bei \textit{Schweine-braten}, der aus einem Stück vom Schwein besteht, und umgekehrt bei \textit{Topf-lappen}, der sicher für mehrere Töpfe gedacht ist, deutlich.


\subsubsection{\textit{er}-Fuge}\label{subsubsection:er-Fuge}

Die \textit{er}-Fuge bildet Kompositionsstammformen von Substantiven, die auch ihren Plural mit -\textit{er} bilden. Bei Komposita wie \textit{Länder-spiel} im Gegensatz zu \textit{Landes-spiel}, \textit{Bücher-regal} im Gegensatz \textit{Buch-preis} und \textit{Häuser-kampf} im Gegensatz zu \textit{Haus-tür} ist die Kompositionsstammform semantisch durch den Plural motiviert. Nicht paradigmatisch wird die \textit{er}-Fuge bei Kardinalzahlen wie in \textit{Zweier-WG} oder \textit{Dreier-bündnis} gebraucht.


\subsubsection{\textit{o}-Fuge}\label{subsubsection:o-Fuge}

Viele Konfixe verfügen über besondere Kompositionsstammformen, die auf \textit{-o} enden, so \zB \textit{thermo} in \textit{Thermo-dynamik} (versus \textit{therm-isch}) oder \textit{chrono} in \textit{Chrono-logie} (versus \textit{Chron-ist}). Die entsprechende Kompositionsstammform auf \textit{-o} wird auch dann verwendet, wenn das Zweitglied dem nativen, d.\,h. heimischen Wortschatz angehört wie bei \textit{Thermo-jacke} oder \textit{Chrono-forscher} in (\ref{ea:Chronoforscher}). 

\ea \label{ea:Chronoforscher} Die anstehende Zeitumstellung bringt nach Ansicht von Wissenschaftlern die innere Uhr aus dem Rhythmus. Dies berge gesundheitliche Risiken, so der Chronoforscher Professor Horst-Werner Korf aus Frankfurt.\footnote{ÄrzteZeitung, 25.03.2010.}
\z 
 
Analogisch wird diese Fuge aber auch bei nativen Erstgliedern mit einem Konfix als Zweitglied gebraucht: \textit{Filzo-kratie}, \textit{Knasto-loge}, \textit{Spielo-thek}, \textit{Wahlo-mat}.


\subsubsection{Minus-\textit{e}-Fuge}\label{subsubsection:Minus-Fuge}

Selten wird wie bei \textit{Schul-bus} oder \textit{Sprach-kurs} die Kompositionsstammform durch Kürzung des \textit{e}-Lautes gebildet. So heißt es auch \textit{Kirsch-saft} aber \textit{Kirsche}, \textit{Kutsch-fahrt} aber \textit{Kutsche} und \textit{Woll-knäuel} aber \textit{Wolle}.  

Bei Komposita wie \textit{Tusch-kasten} oder \textit{Tauf-becken} sollte jedoch nicht von einer \textit{e}-Tilgung der substantivischen Stämme \textit{Tusche} und \textit{Taufe} ausgegangen werden, sondern von verbalen Kompositionsstammformen analog zu \textit{Mal-kasten} oder \textit{Wasch-becken}. 



\subsection{Binarität trotz Mehrgliedrigkeit}\label{subsection:Binarität}
 
 Komposita\is{Kompositum!Binarität} sind grundsätzlich zweigliedrig, das heißt, sie bestehen aus zwei unmittelbaren Konstituenten. Die erste Konstituente bildet die Kompositionsstammform, welche in die zweite Konstituente, die lexikalische Stammform eingebettet ist. Exemplarisch stellen wir die Konstituenz von \textit{Geburtstag} mit Fugenelement in der Abbildung~\ref{Abb:Geburtstag} und von \textit{Waschtag} ohne Fugenelement in Abbildung~\ref{Abb:Waschtag} im\is{Kompositum!Baumdiagramm} Baumdiagramm dar. 
 
 \begin{figure}
	\centering
	\begin{forest}
		[Stammform \\Substantiv Mask\\\textit{Geburtstag} [Kompositionsstammform\\ Substantiv Fem\\\textit{geburts}][Stammform \\ Substantiv Mask\\\textit{Tag}]]	
	\end{forest}
\caption{Binäre Struktur von Komposita mit Fugenelement}\label{Abb:Geburtstag}
\end{figure}

	
	\begin{figure}
		\centering
	\begin{forest}
		[Stammform\\Substantiv Mask\\ \textit{Waschtag}[Kompositionsstammform \\ Verb\\\textit{wasch}][Stammform\\ Substantiv Mask\\\textit{Tag}]]
	\end{forest}
	\caption{Binäre Struktur von Komposita ohne Fugenelement}\label{Abb:Waschtag}
\end{figure}


Das Determinatum \textit{Tag} bezeichnet man analog zum syntaktischen Kopf (vgl. Einführung ins Kapitel~\ref{chap:Phrasen}) als morphologischen Kopf.\footnote{Der Begriff geht auf \citet[235]{Bloomfield1935} zurück: \gqq{Since a blackbird is a kind of bird, and a door-knob a kind of knob, we may say that these compounds have the same function as their head members.}} Das Determinatum ist sowohl der semantische als auch der grammatische Kopf. Als grammatischer Kopf legt das Determinatum die Wortart sowie bei Substantiven die Flexionsklasse fest.

\largerpage[2]
Deshalb schreiben wir im Baumdiagramm den Kopf je nach Wortart groß oder klein. Die Kompositionsstammform schreiben wir grundsätzlich klein, denn ihre Wortart wird im Kompositionsprozess von der Wortart des Kopfes überschrieben. Ob aus \textit{geburts} ein groß zu schreibendes Substantiv oder ein klein zu schreibendes Adjektiv wird, legt der Kopf fest: \textit{Geburtstag} versus \textit{geburtsreif}.

Der Kopf eines\is{Kompositum!Rechtsköpfigkeit} Kompositums\is{Rechtsköpfigkeit} ist der rechts stehende lexikalische Stamm. Der Kopf \gqq{legt}, wie \citet[338--339]{Olsen1991} schreibt, \gqq{in Form eines allgemeinen Prinzips fest, daß komplexe Wörter ihre morphologischen Eigenschaften vom rechten Bestandteil ererben}. Dies gilt, wie wir im Abschnitt~\ref{subsection:Suffigierung} sehen werden, auch für die Suffigierung. Deshalb wird auch vom \textit{Rechtsköpfigkeitsprinzip} gesprochen.\footnote{Dies geht auf die \textit{Righthand Head Rule} von \citet[248]{Williams1981} zurück: \gqq{In morphology we define the head of a morphologically complex word to be the righthand member of the word.}}

Auf den ersten Blick besteht ein Kompositum wie \textit{Kindergeburtstag} aus mehr als zwei Konstituenten, nämlich aus den Kompositionsstammformen \textit{kinder-}, \textit{geburts-} und der lexikalischen Stammform \textit{Tag}. Wir zerlegen jedoch \textit{Kindergeburtstag} schrittweise\is{Kompositum!Binarität} binär zuerst in \textit{kinder-} und \textit{Geburtstag} und anschließend weiter in \textit{geburts-} und \textit{Tag}. Die Konstituentenstruktur stellen wir in\is{Kompositum!Baumdiagramm} Abbildung~\ref{Abb:Kindergeburtstag} dar.

	\begin{figure}
		\centering
	\begin{forest}
		[Stammform\\Substantiv Mask\\ \textit{Kindergeburtstag}[Kompositionsstammform \\ Substantiv Neutr\\\textit{kinder}][Stammform\\ Substantiv Mask\\\textit{Geburtstag} [Kompositionsstammform\\Substantiv Fem \\\textit{geburts}] [Stammform \\ Substantiv Mask\\ \textit{Tag} ] ]]
	\end{forest}
	\caption{Binarität komplexer Komposita}\label{Abb:Kindergeburtstag}
\end{figure}

Theoretisch wäre auch eine andere Analyse möglich. Wir könnten \textit{Kindergeburtstag} auch primär in \textit{kindergeburts} und \textit{Tag} zerlegen und anschließend \textit{Kindergeburt} in \textit{kinder} und \textit{Geburt}. Diese Struktur würde jedoch dem Verständnis von \textit{Kindergeburtstag} zuwiderlaufen. Denn es geht nicht um einen Typ von Tag, der in einer Beziehung zu dem Konzept von Kindergeburt steht, sondern um einen Typ von Geburtstag, der in Beziehung zum Konzept von Kind steht.

\largerpage
Betrachten wir nun noch das Kompositum \textit{Kindergeburtstagsfeier}, so wird klar, dass es nach der gleichen Regel wie \textit{Geburtstag} und wie \textit{Kindergeburtstag} strukturiert ist, nämlich hierarchisch und\is{Kompositum!Binarität} binär. Hier liegt allerdings zwischen der Bildung von \textit{Kindergeburtstag} und der entsprechenden Kompositionsstammform ein Zwischenschritt, da die Komposition lexikalische Stämme und keine Kompositionsstammformen hervorbringt. Die Konstituentenstruktur\is{Kompositum!Baumdiagramm} stellen wir in der Abbildung~\ref{Abb:Kindergeburtstagsfeier} dar.

\begin{figure}
	\centering
	\begin{forest}
[Stammform \\ Substantiv Fem \\ \textit{Kindergeburtstagsfeier} [Kompositionsstammform \\ Substantiv Mask \\ \textit{kindergeburtstags} [Stammform \\ Substantiv Mask \\ \textit{Kindergeburtstag} [Kompositionsstammform \\ Substantiv Neutr \\ \textit{kinder}]  [Stammform \\ Substantiv Mask \\ \textit{Geburtstag} [Kompositionsstammform \\ Substantiv Fem \\ \textit{geburts}] [Stammform \\ Substantiv Mask \\ \textit{Tag}]  ]    ]  ] [Stammform \\ Substantiv Fem \\ \textit{Feier}] ]
	\end{forest}
	\caption{Analyse von \textit{Kindergeburtstagsfeier}}\label{Abb:Kindergeburtstagsfeier}
\end{figure}

\largerpage
Die Regel, die zwei substantivische\is{Rekursivität} Stämme zu einem Kompositum miteinander verbindet, kann genauso auf einfache wie auf komplexe Stämme angewendet werden. Der Output der Regel wie \textit{Geburtstag} kann wie bei der Bildung von \textit{Kindergeburtstag} und weiter zu \textit{Kindergeburtstagsfeier} zu ihrem Input werden. Bildungsprozesse, deren Output auch wieder ihr Input werden kann, heißen \textit{rekursiv} (lat. \textit{recurrere} \gq{zurücklaufen}). Auch eine äußerst komplexe Bildung wie der bekannte \textit{ Donaudampfschifffahrtsgesellschaftskapitän} basiert auf der rekursiven Anwendung derselben Regel, die zwei substantivische Stämme unabhängig von ihrer internen Komplexität miteinander verbindet.\footnote{\citet[221]{Augst2001} ermittelt in einer empirischen Studie, dass der Anteil von fünf- und mehrgliedrigen Komposita in Zeitungstexten ca. 1\% ist. Nach \citet[13, 20]{Ortner_Komposita1991} sind in Abhängigkeit von der Textsorte bis zu 10\% der Substantiv-Substantiv-Komposita dreigliedrig und 1,5\% viergliedrig. Der Prototyp ist ein zweigliedriges Kompositum wie \textit{Apfelbaum} oder \textit{Haustür}.}

Allerdings wird\is{Ambiguität!morphologisch} die Reihenfolge der jeweiligen Regelanwendung nicht durch die Regel festgelegt. Deshalb könnten wir \textit{Kindergeburtstagsfeier} auch in \textit{kinder-} und \textit{geburtstagsfeier} als unmittelbare Konstituenten zerlegen. Diese Analysemöglichkeit stellen wir in\is{Kompositum!Baumdiagramm} Abbildung~\ref{Abb:AlternativeKgbf} dar.

\begin{figure}
	\centering
	\begin{forest}
[Stammform \\ Substantiv Fem \\ \textit{Kindergeburtstagsfeier} [Kompositionsstammform \\ Substantiv Neutr \\ \textit{kinder}] [Stammform \\ Substantiv Fem \\ \textit{Geburtstagsfeier} [Kompositionsstammform \\ Substantiv Mask \\ \textit{geburtstags} [Stammform \\ Substantiv Fem \\ \textit{Geburtstag} [Kompositionsstammform \\ Substantiv Fem \\ geburts] [Stammform \\ Substantiv Mask\\ \textit{Tag}]  ]  ] [Stammform \\ Substantiv Fem \\ \textit{Feier}]     ]     ]  
	\end{forest}
	\caption{Alternative Analyse von \textit{Kindergeburtstagsfeier}}\label{Abb:AlternativeKgbf}
\end{figure}

Für die in den Abbildungen~\ref{Abb:Kindergeburtstagsfeier} und \ref{Abb:AlternativeKgbf} dargestellten Analysen gibt es jeweils einen guten Grund. Sowohl \textit{Kindergeburtstag} als auch \textit{Geburtstagsfeier} sind gängige Bildungen, wodurch beide Zwischenstufen plausibel sind. Die Abbildung~\ref{Abb:Kindergeburtstagsfeier} steht für das Verständnis, dass es sich um einen Typ von Feier handelt, der durch eine Beziehung zu dem Konzept von Kindergeburtstag bestimmt ist. Die Abbildung~\ref{Abb:AlternativeKgbf} steht für das Verständnis, dass es sich um einen Typ von Geburtstagsfeier handelt, der durch eine Beziehung zu dem Konzept von Kind bestimmt ist. Der Bedeutungsunterschied beruht auf verschiedenen Strukturen.

Mitunter sind die semantischen Auswirkungen struktureller Ambiguität jedoch viel größer als bei \textit{Kindergeburtstags-feier} versus \textit{Kinder-geburtstagsfeier}. Das ist der Fall bei den verschiedenen Analysemöglichkeiten des Kompositums \textit{Frauenfilmfestival}. So kann es sich bei \textit{Frauenfilmfestival} um ein Festival mit Bezug zu Frauenfilmen handeln. Möglich ist aber auch die Lesart als Filmfestival, das in irgendeiner Beziehung zu Frauen steht. Die erste Lesart beruht auf der Strukturanalyse in\is{Kompositum!Baumdiagramm} Abbildung~\ref{Abb:Frauenfilmfestival1}, die zweite Lesart auf der Strukturanalyse in Abbildung~\ref{Abb:Frauenfilmfestival2}.

	\begin{figure}
		\centering
	\begin{forest}
		[Stammform\\ Substantiv Neutr\\ \textit{Frauenfilmfestival} [Kompositionsstammform\\Substantiv Mask\\ \textit{frauenfilm} [Stammform\\ Substantiv Mask\\ \textit{Frauenfilm} [Kompositionsstammform\\Substantiv Fem \\ \textit{frauen}] [Stammform\\ Substantiv Mask\\ \textit{Film}]  ]   ] [Stammform\\Substantiv Neutr\\ \textit{Festival}]]
	\end{forest}
	\caption{\textit{Festival mit Bezug zu Frauenfilmen}}\label{Abb:Frauenfilmfestival1}
\end{figure}


	\begin{figure}
		\centering
	\begin{forest}
		[Stammform\\ Substantiv Neutr\\ \textit{Frauenfilmfestival} [Kompositionsstammform\\ Substantiv Fem\\ \textit{frauen} ] [Stammform\\ Substantiv Neutr\\ \textit{Filmfestival} [Kompositionsstammform\\ Substantiv Mask\\ \textit{film}] [Stammform\\ Substantiv Neut\\ \textit{Festival}]  ]  ]
	\end{forest}
	\caption{\textit{Filmfestival mit Bezug zu Frauen}}\label{Abb:Frauenfilmfestival2}
\end{figure}

Bevor wir uns im nächsten Abschnitt~\ref{section:Derivation} mit der Derivation beschäftigen, diskutieren wir in dem folgenden Exkurs~\ref{subsection:Übergänge} Bildungen wie \textit{Einfamilienhaus}, die wir nicht ohne Weiteres binär in \textit{Einfamilien-haus} zerlegen können. Ebenso gehen wir auf Bildungen wie (\textit{der}) \textit{Wir-schaffen-das-Tonfall} ein, deren Erstglied keine Kompositionsstammform, sondern eine Wortgruppe ist. Als Übergangsbereich zur Derivation führen wir Konfix-Bildungen wie \textit{Spielosoph} sowie Bildungen wie \textit{scheißegal} oder \textit{knitterarm} an, deren Erst- oder Zweitglied einem Prä- bzw. Suffix ähnelt.

\newpage
\begin{Exkurs}{Übergänge: Syntax/Morphologie und Komposition/Derivation}\label{subsection:Übergänge} 

Im Abschnitt~\ref{subsection:Fuge} haben wir paradigmatische Fugenelemente wie in \textit{Diebes-gut} und die Binarität aus Determinans und Determinatum historisch auf Flexionsformen innerhalb von Wortgruppen (\textit{des Dieb-es Gut}) bezogen. Syntaktische Nachbarn wurden zu einem Wort verbunden, also univerbiert\is{Univerbierung}. Aus syntaktischer Nachbarschaft wurde vor dem Hintergrund älterer Komposita qualitativ (grammatisch und begrifflich) etwas Neues, nämlich ein Muster zur Bildung von Komposita. Dessen Bestandteile können für sich genommen als syntaktische Wörter gebraucht werden. Sie sind wortfähig. Um Abweichungen hiervon geht es im Folgenden.

Nicht wortfähige Erstglieder wie *\textit{Einfamilie} in \textit{Einfamilienhaus} oder *\textit{Zweiraum} in \textit{Zweiraumwohnung} sind syntaktisch im Sinne von \textit{für eine Familie} bzw. \textit{mit zwei Räumen} motiviert. Seit \citet[14--15]{Henzen1965} werden sie auch als\is{Zusammenbildung} \textit{Zusammenbildung} bezeichnet. Diese Bezeichnung verdeckt, dass die Gesamtstruktur trotz nicht wortfähiger Erstglieder eindeutig zu den Determinativkomposita gehört. Auch die interne Struktur der Erstglieder ist nicht syntaktischer Natur. Das zeigt sich daran, dass die syntaktisch motivierte Flexion von \textit{eine} in \textit{eine Familie} bzw. von \textit{Räume} in \textit{zwei Räume} gerade nicht in die Komposita \textit{Einfamilienhaus} und \textit{Zweiraumwohnung} eingeht. 

Anders als bei \textit{Einfamilienhaus} kann auch das Determinatum nur wortartig sein, ohne eigenständig als Wort gebraucht zu werden. Das ist bei \textit{Hemmer} in \textit{Appetithemmer}, \textit{Steiger} in \textit{Bergsteiger}, \textit{Scheuche} in \textit{Vogelscheuche} oder \textit{Blüher} in \textit{Frühlingsblüher} der Fall. Diese Zweitglieder benutzen wir nicht oder kaum als selbständige Wörter. Die syntaktische Motivation dieser Komposita kann wie folgt paraphrasiert werden: \textit{was den Appetit hemmt}, \textit{wer auf Berge steigt}, \textit{was Vögel} (\textit{ver})\textit{scheucht} und \textit{was im Frühling blüht}. Dennoch folgen die entsprechenden Komposita dem etablierten Muster genauso wie \textit{Romanleser} und \textit{Gelegenheitsraucher}. Die Erstglieder werden unflektiert in die Zweitglieder eingebettet. Auf ähnliche Bildungen wie \textit{grünäugig} und \textit{Dickhäuter} gehen wir in dem Exkurs~\ref{subsubsection:DerivationoderKomposition} ein, da wir sie anders analysieren.     


Bei einer besonderen Art substantivischer\is{Phrasenkompositum} Komposita kommen wie in (\ref{ea:Zitat_Subst}) syntaktische Strukturen bis hin zum Satz als Erstglied vor. Ihre Gesamtstruktur entspricht dem rechtsköpfigen Muster der Determinativkomposita.\footnote{Für \citet[43]{Hein2015} besteht der wesentliche Unterschied zum Prototyp des Kompositums darin, dass Phrasenkomposita heute auf substantivische Zweitglieder beschränkt sind. Wir sehen hierin gerade die Nähe zum Prototyp des Kompositums. Die Abweichung vom Prototyp ergibt sich aus der phrasalen Struktur des Erstglieds.}


\ea \label{ea:Zitat_Subst} Sie verweigert den kanzlerischen Wir-schaffen-das-Tonfall und zerfleddert die Alles-wird-gut-Komödie zu einem kabarettistischen Bilderbogen.\footnote{die tageszeitung, 12.12.2015.}
\z 

Aufgrund der internen syntaktischen Struktur der Erstglieder werden diese Bildungen als \textit{Phrasenkomposita} bezeichnet. Entsprechende Phrasen werden analog zu Kompositionsstammformen gebraucht. Wie in den Belegen dominiert die Zusammenschreibung mit Bindestrich.\footnote{Cf. \citet[349--353]{Hein2015}.} Es handelt sich besonders wie in (\ref{ea:Zitat_Subst}) um Zitate (\textit{wir schaffen das}) oder im weitesten Sinne um etablierte Ausdrücke (\textit{alles wird gut}). Diese phrasalen Erstglieder sind bereits vorfabriziert und oft als Ganzes, also wie Wortstämme und häufig gebrauchte Wortformen, abgespeichert.\footnote{Eine Übersicht mit Beispielen bietet \citet[171]{Meibauer2003}.} Eine Vorprägung ist jedoch keine Bedingung für Phrasenkomposita. Wir bilden sie auch ad hoc wie in (\ref{ea:adhocPhrasKomp}), das heißt, innerhalb von Komposita verpacken wir jene Phrasen, als ob es eine entsprechende wortähnliche Prägung gäbe.

\ea \label{ea:adhocPhrasKomp} Das Bild ist ein großes Versprechen. Eine dunkle Tänzerin wirbelt durch die Luft, in ihrem Blick dieses unwiderstehliche Es-gibt-hier-noch-ein-kleines-Geheimnis-Lächeln [\ldots{}]\footnote{die tageszeitung, 09.07.2001.}
\z 

Wesentlich ist der wortbildende Effekt, durch den einzelne flektierte Formen des Erstgliedes genauso wie Kompositionsstammformen syntaktisch nicht mehr zugänglich sind. Deshalb kann mit \textit{das} in (\ref{ea:adhocPhrasKomp*das}) kaum \textit{ein kleines Geheimnis} wieder aufgenommen werden.

\ea \label{ea:adhocPhrasKomp*das} [\ldots{}] in ihrem Blick dieses unwiderstehliche Es-gibt-hier-noch-ein-kleines-Geheimnis-Lächeln [\ldots{}] -- Das wird niemand lüften.  
\z 

Besonders häufig ist das formale Muster mit einer NP aus Zahladjektiv und Maßsubstantiv wie in (\ref{ex:ZahlNP_Subst}) als Erstglied.\footnote{Cf. \citet[43]{Ortner_Komposita1991}.}

\ea \label{ex:ZahlNP_Subst} Für Ter Haigasun aber bedeutete das Fünf-Minuten-Glück dieser Frau gar nichts gegenüber einem reinen Eintritt in die Ewigkeit.\footnote{Franz Werfel: \textit{Die Vierzig Tage des Musa Dagh II}. Fischer 1947 [1933], S. 326.}
\z 

Die interne Struktur erinnert an syntaktische Maßausdrücke wie \textit{fünf Flaschen Milch}. Diese haben wir in Exkurs~\ref{subsection:Apposition} als enge Apposition bezeichnet. Im Unterschied zu engen Appositionen ist die Maßangabe \textit{fünf Minuten} jedoch in das Zweitglied des Kompositums \textit{Glück} eingebettet. Auf diesem Muster basieren die eingangs aufgeführten Komposita \textit{Einfamilienhaus} und \textit{Zweiraumwohnung}.

Nicht zu den Komposita zählen erstarrte, zu einem Wort zusammengerückte syntaktische Strukturen wie (\textit{eine}) \textit{Handvoll}, \textit{Zeitlang} oder (\textit{ein}) \textit{Spaltbreit}. Denn sie entsprechen nicht dem rechtsköpfigen und endozentrischen Bildungsmuster der Determinantivkomposita. Das ist auch bei syntaktisch geprägten Satzwörtern wie \textit{Taugenichts} oder \textit{Störenfried} der Fall. Denn hier geht es nicht um eine besondere Art von \textit{Nichts} oder \textit{Friede}, sondern um jemanden, der nichts taugt oder der den Frieden stört. Auf diese Bildungen kommen wir unter dem Stichwort \textit{syntaktische Konversion} im Abschnitt~\ref{subsection:Konversion} noch einmal zurück.  


Einen stabilen\is{Konfixkompositum} Randbereich der Komposition hin zur Derivation stellen Bildungen wie \textit{Philosoph} oder \textit{Velodrom} dar. Konfixe können typischerweise nicht alleine als syntaktisches Wort gebraucht werden. Anders als Derivationsaffixe verbinden sie sich jedoch untereinander und teilweise mit Derivationsaffixen (vgl. Abschnitt~\ref{subsection:Konfixe}). 

Wenige Konfixe können Erstglieder von Determinativkomposita wie in (\ref{ea:KonfixErstglied}) sein. Hierzu zählen auch heimische Konfixe wie \textit{stief-} oder \textit{zimper-}. Häufiger sind Konfixe die Zweitglieder von substantivischen Komposita. In diesen fungieren sie als morphologische Köpfe. Die Erstglieder sind ebenfalls Konfixe (\ref{ex:KonfixZweitglied1}), aber auch heimische Wortstämme (\ref{ex:KonfixZweitglied2}). Auch heimische Erstglieder werden dann in der Regel durch das Fugenelement \textit{-o} als Kompositionsstammformen ausgezeichnet.  

\ea \label{ea:Konfixkomposita}
\ea \label{ea:KonfixErstglied} Biomilch, Ökohaus, Thermojacke, Stiefschwester
\ex \label{ex:KonfixZweitglied1} Philosoph, Autodrom, Kosmonaut, Mikrophon.
\ex \label{ex:KonfixZweitglied2} Spielosoph, Schwimmodrom, Wahlomat
\z 
\z 

%Konfixkomposita sind besondere Determinativkomposita, da Konfixe typischerweise alleine nicht wortfähig sind.

Manche Determinativkomposita wie \textit{prunkvoll} und \textit{giftarm} sind metaphorisch motiviert. Die Semantik des morphologischen Kopfes wird durch das Erstglied stark modifiziert. Dinge, die wir als \textit{prunkvoll} bezeichnen, sind nicht einfach auf eine bestimmte Art und Weise voll, sondern eben voller Prunk. Wir können \textit{prunkvoll} aber auch mit dem semantischen Gewicht auf \textit{prunk} als \gq{prunkig} verstehen. Der metaphorische Bedeutungsbeitrag des Kopfes wird reihenbildend genutzt wie \zB in \textit{absichtsvoll}, \textit{anspruchsvoll}, \textit{erbarmungsvoll}, \textit{glanzvoll} und \textit{zweckvoll}. Das Kompositum \textit{giftarm} drückt weniger einen Untertyp des Armseins aus, sondern dass etwas arm an Giftstoffen oder eben wenig giftig ist. Auch \textit{arm} wird reihenbildend als Zweitglied genutzt, \zB in \textit{blutarm}, \textit{farbarm}, \textit{knitterarm}, \textit{rauscharm} und \textit{spracharm}. Bei diesen Bildungen haben die  Stammformen \textit{voll} und \textit{arm} eine gewisse Ähnlichkeit mit Affixen. Auch manche Erstglieder wie \textit{stock} in \textit{stocksteif}, \textit{tod} in \textit{todlangweilig}, \textit{brand} in \textit{brandeilig}, \textit{grund} in \textit{grundanständig}, \textit{hunde} in \textit{hundeelend} oder \textit{scheiß} in \textit{scheißegal} haben eine gewisse Ähnlichkeit mit Affixen. Die Vergleichslesarten beruhen jedoch metaphorisch auf der Bedeutung der Kompositionsstammformen im Sinne von steif wie ein Stock, so eilig wie man es bei einem Brand hat usw.\footnote{Aufgrund der Bedeutungsverschiebung und der Reihenbildung klassifiziert \citet[10--11]{Motsch2004} diese Erst- und Zweitglieder als \textit{Affixoide}, also als affixähnlich. Wir bestimmen sie wie \citet[310]{FleischerBarz2012} als metaphorisch motivierte Komposita, die wie andere Komposita aus Stammformen und nicht aus einer Stammform und einem Affixoid bestehen.} Tatsächlich gehen viele Derivationsaffixe historisch auf lexikalische Stammformen zurück. Um die Derivation geht es in dem folgenden Abschnitt \ref{section:Derivation}.
\end{Exkurs}


\section{Derivation}\label{section:Derivation}
    
Als\is{Derivation} Derivation bezeichnet man die Ableitung von lexikalischen Stämmen aus Derivationsstammformen. Anders als bei der Komposition, die ja auch eine Art von Abwandlung einer lexikalischen Stammform, eben des Determinatums ist, sind an der Derivation Affixe beteiligt, die nicht selbstständig als syntaktisches Wort gebraucht werden können, die also nicht wortfähig sind. 

Die Ableitung kann kombinatorisch wie von \textit{Kind} zu \textit{kind-lich} oder von \textit{arbeit}(\textit{en}) zu \textit{be-arbeit}(\textit{en})  erfolgen. In Abschnitt~\ref{subsection:Suffigierung} stellen wir die Suffigierung im\is{Suffigierung!nominal} nominalen, d.\,h. im substantiv- und adjektivbildenden Bereich dar. Nur am Rande gehen wir auf die Suffigierung im Bereich der Verbbildung ein. In Abschnitt~\ref{subsection:Präfigierung} geht es um die Präfigierung im verbalen Bereich. Ebenfalls nur am Rande weisen wir auf die Präfigierung im nominalen Bereich hin. 

Bei nicht kombinatorischer Ableitung durch den Wechsel des Stammvokals wie von \textit{trink}(\textit{en}) zu \textit{tränk}(\textit{en}) spricht man von \textit{impliziter Derivation}. Um diese geht es in Abschnitt~\ref{subsection:ImplDerivation}. 

Die Ableitung ohne Affix wie von dem Verbstamm \textit{lauf}(\textit{en}) zum Substantiv (\textit{der}) \textit{Lauf} heißt \textit{Konversion}, die wir in Abschnitt~\ref{subsection:Konversion} vorstellen. 

\begin{Definition}{Derivation}
    
\noindent Derivation ist ein Wortbildungsprozess, bei dem ein lexikalisches Morphem oder eine Morphemkette als Derivationsstammform meist durch ein Affix zu einem anderen lexikalischen Morphem abgewandelt wird. Affixlose Derivationen geschehen über den Wechsel des Stammvokals oder ohne Formsignal.  
\end{Definition}



\subsection{Suffigierung}\label{subsection:Suffigierung}

Lexikalische Stämme und Konfixe können durch ein\is{Suffigierung} Suffix zu einem neuen Stamm abgeleitet werden. Dies kann unter Beibehaltung der Wortart wie von \textit{grün} zu \textit{grün-lich} oder von \textit{Haus} zu \textit{Häus-chen} geschehen. Primär ist jedoch mit der Suffigierung ein Wortartenwechsel (\textit{Transposition}) verbunden. So werden aus der verbalen Derivationstammform \textit{les-} die Substantivstämme \textit{Les-er} oder \textit{Les-ung} abgeleitet. Manche Suffixe haben sich innerhalb von Komposita aus ursprünglichen lexikalischen Stammformen entwickelt. So wie\is{Rechtsköpfigkeit} letztere legen auch Suffixe die Wortart und substantivische Suffixe außerdem die Flexionsklasse der Derivata fest. Durch Suffigierung gebildete Wörter sind ebenfalls binär strukturiert. Sie bestehen aus einer Derivationsstammform als semantischer Basis und einem Suffix als morphologischem Kopf. 

Komposition und Suffigierung sind sich ähnlich. Der Strukturaufbau ist binär und die rechts stehende Einheit ist der Kopf. Ein wesentlicher Unterschied besteht aber darin, dass das Zweitglied eines Kompositums ein eigenständiges Zeichen ist und als solches einer Wortart angehört. Das Kompositum \textit{Freizeit} ist ein feminines Substantiv, weil \textit{Zeit} ein feminines Substantiv ist. Die Bedeutung von \textit{Freizeit} ist eine Art von Zeit. Auch das Derivatum \textit{Freiheit} ist ein feminines Substantiv. Das erkennen wir am Suffix \textit{-heit}. Aber wir würden nicht sagen, dass \textit{Feiheit} ein feminines Substantiv ist, weil \textit{-heit} ein feminines Substantiv ist. Auch bedeutet \textit{Freiheit} nicht eine bestimmte Art von \textit{-heit}.\footnote{Zur theoretischen Diskussion verweisen wir auf \citet{Hoehle1982DerivKomp}, für den Derivation und Komposition auf \textit{einem} Wortbildungsprozess beruhen sowie auf \citet{Reis1983} und \citet{Welke1995DerivKomp}, die für eine differenzierte Analyse plädieren, welche auch hier verfolgt wird.}

Einem Suffix wie \textit{-heit} begegnen wir nur als Bestandteil komplexer Wörter wie \textit{Freiheit}, \textit{Schönheit} oder \textit{Klarheit}. Deren Bestandteile \textit{frei}, \textit{schön} und \textit{klar} kennen wir auch als einfache Wörter. Erst wenn wir, wie \citet[543--544]{Booij2010} zeigt, diese einfachen Zeichen mit den komplexen Zeichen wie in (\ref{ea:einfach_komplex}) aufeinander beziehen, machen wir Strukturen wie in (\ref{ea:Stamm_Derivatum}) mit \textit{-heit} als Bestandteil aus.

\ea \label{ea:einfach_komplex}
\ea \label{ea:frei_freiheit} frei -- Freiheit
\ex \label{ex:schön_schönheit} schön -- Schönheit
\ex \label{ex:klar_klarheit} klar -- Klarheit
\z 
\z 

\ea \label{ea:Stamm_Derivatum}
\ea \label{ea:frei_heit} Frei[heit]
\ex \label{ex:schön_heit} Schön[heit]
\ex \label{ex:klar_heit} Klar[heit]
\z 
\z 

Die Segmentierung in (\ref{ea:Stamm_Derivatum}) ist Voraussetzung dafür, die Form \textit{-heit} aus entsprechenden Bildungen als Suffix wie in (\ref{ea:_heit}) zu abstrahieren. Zu dieser Abstraktion gehört die Wortart, über die das Suffix operiert, die resultierende Wortart und der Bedeutungsaspekt der Abwandlung. 

\ea \label{ea:_heit} [[X]\textsubscript{Adjektiv} \textit{-heit}]\textsubscript{Substantiv/Fem}: Name der Eigenschaft des Adjektivs
\z  

Das Suffix \textit{-heit} wandelt aus Vertretern der Wortart Adjektiv feminine Substantive ab. Das so entstandene Substantiv lässt uns über die vom Adjektiv bedeutete Eigenschaft sprechen, ohne ihren Träger zu benennen. Insofern ist das derivierte Substantiv der abstrakte Name der Eigenschaft. Auf der Grundlage von (\ref{ea:_heit}) bilden wir neue Substantive wie \textit{Grauheit} oder \textit{Hübschheit}. 

Mit \citet[34]{Eichinger2000} leisten Suffixe eine klassematische Einordnung der lexikalischen Basis, indem sie eine \gqq{Sehweise}, einen bestimmten Bedeutungsaspekt der Basis in Perspektive bringen. Am deutlichsten wird das, wenn wir von einem Verbszenario (vgl. Definition~\ref{Def:Szenario} in Abschnitt~\ref{section:KonzeptVerbvalenz}) ausgehen. So kann aus \textit{liefer}(\textit{n}) derjenige, der liefert, als \textit{Nomen agentis} (Handelnder) in Perspektive gebracht werden: \textit{Liefer-ant} oder \textit{Liefer-er}. Die Suffigierung \textit{Liefer-ung} bedeutet entweder das zum Redegegenstand gemachte Ereignis, das \gq{Geliefertwerden} oder das \gq{Liefern} als sogenanntes \textit{Nomen actionis} wie in (\ref{ea:Lieferung_actionis}) oder aber das Objekt, das Gelieferte, als sogenanntes \textit{Nomen acti} wie in (\ref{ex:Lieferung_acti}). 

\ea \label{ea:nomen_actionis_acti}
\ea \label{ea:Lieferung_actionis} Die Lieferung dauert einen Werktag. 
\ex \label{ex:Lieferung_acti} Die Lieferung liegt beim Nachbarn.
\z 
\z 

Ebenso kann das Adjektiv \textit{liefer-bar} abgeleitet werden, das die Eigenschaft \gq{geliefert werden zu können} ausdrückt. 

Der Gewinn von solchen Ableitungen ist die begriffliche Erweiterung einer Grundbedeutung mit entsprechender Veränderung der syntaktischen Verwendungsmöglichkeiten. Die Leistung einzelner Suffixe wird in dem folgenden Abschnitt \ref{subsubsection:Suffixe} dargestellt. 


\subsubsection{Suffixe}\label{subsubsection:Suffixe}

In den folgenden Unterabschnitten geht es um die Grammatik und Leistung der folgenden Suffixe zur Bildung von Substantiven und Adjektiven: \textit{-heit} wie in \textit{Frei-heit}, \textit{-ig} wie in \textit{staub-ig}, \textit{-ung} wie in \textit{Liefer-ung}, \textit{-erei} wie in \textit{Bäck-erei}, \textit{-bar} wie in \textit{liefer-bar}, \textit{-er} wie in \textit{Lehr-er}, \textit{-ling} wie in \textit{Lehr-ling}, \textit{-in} wie in \textit{Lehrer-in}, \textit{-schaft} und \textit{-tum} wie in \textit{Lehrer-schaft} und \textit{Streber-tum} sowie um \textit{-chen} und \textit{-lein} wie in \textit{Mäus-chen} und \textit{Mäus-lein}.


\subsubsubsection{\textit{-heit}}\label{subsubsection:heit}

Wie\is{Suffix!-\textit{heit}} einleitend erwähnt, werden durch \textit{-heit} einfache adjektivische Derivationsstammformen und das adjektivische Partizip II (\textit{Verliebt-heit}) zu Substantiven abgeleitet. Adjektive, die wie \textit{träge} oder \textit{fade} auf \textit{-e} enden, bilden ihre Derivationsstammform mittels Minus-Fuge wie in \textit{Träg-heit} und \textit{Fad-heit} ohne \textit{-e}. Adjektive, die selbst durch Suffixe abgeleitet sind wie \textit{glücklich} oder \textit{gruselig} werden durch \textit{-keit} wie in \textit{Glücklich-keit} und \textit{Gruselig-keit} zum Substantiv abgeleitet. 

Die Kombination von \textit{-ig} und \textit{-keit} wurde zu \textit{einem} Suffix \textit{-igkeit} wie in \textit{Haltlos-igkeit} (*\textit{haltlosig}) umgedeutet. So einen Prozess bezeichnet man als \textit{Reanalyse}. Mitunter kommt es zu Doppelbildungen wie \textit{Nett-heit} gegenüber dem gebräuchlicheren \textit{Nett-igkeit}. Auf diese Weise zum Substantiv abgeleitete Adjektive abstrahieren als Namen von Eigenschaften von ihrem jeweiligen Träger. Als Substantive bilden sie Nominalphrasen und können als Subjekt oder Objekt von einem übergeordneten Prädikat fungieren. Metonymisch können sie zur Bezeichnung des Eigenschaftsträgers gebraucht werden: \textit{Jemand ist eine Schönheit}.

\subsubsubsection{\textit{-ig}}\label{subsubsection:ig}

Durch\is{Suffix!-\textit{ig}} \textit{-ig} werden substantivische Derivationsstammformen wie \textit{stein} und \textit{staub} zu Adjektiven wie \textit{stein-ig} oder \textit{staub-ig} abgeleitet. Das Derivat bezeichnet eine Eigenschaft, die im weitesten Sinne durch die Bedeutung des Substantivs charakterisiert ist.\footnote{Die von Nomina actionis abgeleiteten Adjektive wie \textit{flücht-ig} von \textit{Flucht}, \textit{freud-ig} von \textit{Freude} sowie formal scheinbar von Verben abgeleitete Adjektive wie \textit{find-ig} (ursprünglich \textit{fünd-ig} von \textit{Fund}) bilden den Übergang zu deverbalen Ableitungen wie \textit{holp}(\textit{e})\textit{r-ig}, \textit{hör-ig} oder \textit{nörg}(\textit{e})\textit{l-ig}. Cf. \citet[95]{Paul1920Wb}.}

Die so\is{Suffigierung!Binarität} abgeleiteten Adjektive können wiederum durch \textit{-keit} zu Substantiven wie \textit{Steinig-keit} oder \textit{Staubig-keit} abgeleitet werden. Auch diese komplexen Derivata sind binär gebaut und rechtsköpfig, was die Konstituentenstruktur in der\is{Suffigierung!Baumdiagramm} Abbildung~\ref{Abb:Steinigkeit} zeigt.

	\begin{figure}
	\centering
	\begin{forest}
	[Stammform \\ Substantiv Fem \\ \textit{Steinigkeit} [Derivationsstammform \\ Adjektiv \\ \textit{steinig} [Stammform \\ Adjektiv \\ \textit{steinig} [Derivationsstammform \\ Substantiv Mask \\ \textit{stein}] [Suffix \\ bildet Adjektiv \\ \textit{-ig}]   ]  ] [Suffix \\ bildet Substantiv Fem \\ \textit{-keit}]  ]
	\end{forest}
	\caption{Binäre Konstituentenstruktur von \textit{Steinigkeit}}\label{Abb:Steinigkeit}
\end{figure}

Wenn es sich wie bei \textit{Haltlos-igkeit} um das reanalysierte komplexe Suffix \textit{-igkeit} handelt, so bildet dies eine Konstituente, die sich mit der adjektivischen Basis \textit{haltlos} verbindet.  

\subsubsubsection{\textit{-ung}}\label{subsubsection:ung}

Mit\is{Suffix!-\textit{ung}} \textit{-ung} werden Derivationsstammformen von meist telischen (vgl. Abschnitt~\ref{subsection:Verbklassen}) und oft komplexen Verben zu Substantiven abgeleitet. Ableitungen von Adjektiven sind wie *\textit{Schön-ung} nicht möglich. Der telische Verbstamm \textit{verschöner-} kann zu \textit{Verschöner-ung} abgeleitet werden. Von atelischen Verben wie \textit{arbeiten} werden kaum Substantive mit \textit{-ung} abgeleitet. Im Gegensatz dazu wurde aus dem telischen Präfixverb \textit{bearbeiten} das Substantiv \textit{Bearbeit-ung} gebildet.\footnote{Ableitungen von Substantiven wie \textit{Stall-ung}, \textit{Wald-ung} oder \textit{Zeit-ung} (ursprünglich: \gq{das zu einer Zeit Geschehene}) sind Relikte eines heute nicht mehr aktiven Bildungsmusters.}

Ableitungen mit \textit{-ung} dienen primär, wie bereits erwähnt, als Namen für Ereignisse (Nomen actionis) oder für deren Resultate (Nomen acti). 

Bei der\is{Nomen actionis} Lesart als Nomen actionis wird ein Geschehen als Redegegenstand ausgedrückt, das heißt zu einem abstrakten Referenten verdinglicht. Die Paraphrase eines Substantivs wie \textit{Lieferung} ist primär vorgangsbezogen als \gq{Geliefertwerden} wie in \textit{die Lieferung der Ware} oder \textit{die Lieferung durch den Kurier}. 

Manche Substantivierungen mit \textit{-ung} können auch als Ergebnis oder Resultat des vom Verb ausgedrückten Szenarios verstanden werden. Diese Lesart wird als\is{Nomen acti} \textit{Nomen acti} bezeichnet. So können wir \textit{Zerstörung} als Nachzustand eines Zerstörens im Sinne von \gq{Zerstörtsein} verstehen: \textit{die Zerstörung der Stadt lastet immer noch schwer auf den Menschen}. Da Nachzustände Eigenschaften sind, gibt es Konkurrenzbildungen mit \textit{-heit} wie \textit{Zerstörtheit}. Das Substantiv \textit{Lieferung} können wir auch als dingliches Resultat verstehen als das Gelieferte. Bei dem Rückbezug auf die Objektstelle des Verbs spricht man auch von \textit{Nomina patientis} (vgl. \textit{Patiens} Abschnitt~\ref{subsection:Verbklassen}). 

Keine der Lesarten ist grammatisch ausgewiesen. Viele deverbale Substantive weisen neben der Lesart als Nomen actionis auch die als Nomen acti/patientis auf. Erst innerhalb eines gegebenen Kontextes wird die entsprechende Lesart festgelegt. 


Die in\is{Valenzvererbung} der Verbvalenz angelegten Ergänzungen können an das Substantiv vererbt werden. Bei der Lesart als Nomen actionis kann primär das ursprüngliche (im Fokus stehende) Akkusativobjekt als Genitivattribut wie bei \textit{Lieferung der Ware} ausgedrückt werden. Wir verstehen die Zugehörigkeitsrelation der Ware zur Lieferung entweder als Handlungsgegenstand (Ware, die jemand liefert) oder als Vorgangs- oder Zustandsträger, d.\,h. als Ware, die geliefert wird. Das ursprüngliche Subjekt kann mit \textit{durch} wie bei \textit{Lieferung der Ware durch den Boten} als Verursacher des Vorgangs angeschlossen werden. Bei Eigennamen ist dies auch durch einen vorangestellten Genitiv möglich: \textit{Cäsars Eroberung des Westens}. Wenn das Verb atelisch ist, was bei \textit{ung}-Ableitungen selten der Fall ist, so kann auch das ursprüngliche Subjekt mit dem Genitiv angeschlossen werden, so wie bei \textit{Hoffnung des Kindes}. Die vom Verb regierte Präposition wird an das Substantiv vererbt: \textit{die Hoffnung des Kindes auf eine schöne Zukunft}. Ist das atelische Verb transitiv wie \textit{beobachten}, dann können entweder das ursprüngliche Objekt oder das ursprüngliche Subjekt als Genitivattribut an das Substantiv vererbt werden. Ohne Kontext kann \textit{Beobachtung des Diebes} sowohl als Vorgang verstanden werden, bei dem der Dieb beobachtet wird, als auch als Tätigkeit, die der Dieb ausführt.   

Bei Nomina acti können wir ein Genitivattribut auch auf die erste verbale Ergänzung des zugrunde liegenden Verbs beziehen: \textit{die Lieferung des Boten} oder \textit{die Lösung des Schülers}. Die Zugehörigkeitsrelation zwischen Bote und Lieferung bzw. zwischen Schüler und Lösung wird in der Regel unter Rückgriff auf unser Weltwissen und den entsprechenden Kontext verursachend verstanden, also als Lieferung, die durch den Boten geleistet wurde bzw. als Lösung, die durch den Schüler erbracht wurde.\footnote{Bei Interesse an der Vererbung der verbalen Ergänzungen an die Substantivkonstruktion verweisen wir auf \citet{Eichinger1995} und \citet{Welke2016}.}

Mit Nomina actionis wie \textit{Bedienung} oder \textit{Vertretung} können auch Personen bezeichnet werden. Diese Lesart verstehen wir nicht als Bedeutung der \textit{ung}-Ableitung, sondern als eine systematische Bedeutungsübertragung vom Geschehen auf einen am Geschehen Beteiligten.

\subsubsubsection{\textit{-erei}}\label{subsubsection:erei}

Ein\is{Suffix!-\textit{erei}} anderer Typ von Nomina actionis wird durch (\textit{-er})\textit{ei} wie in \textit{Raucherei} oder \textit{Telefoniererei} gebildet. Das Suffix \textit{-ei} ist eine lautliche Entwicklung des aus dem Lateinischen und Französischen entlehnten \textit{-ie} wie es in Lehnwörtern wie \textit{Abtei} oder \textit{Sakristei} vorkommt.\footnote{Cf. \citet[81--82]{Paul1920Wb}.} Analog zum Französischen \textit{-ie} dient \textit{-ei} der Ableitung von Nomina agentis wie \textit{Bäcker} zur Bezeichnung des professionellen Tätigkeitsortes wie \textit{Bäcker-ei}. Substantive auf \textit{-erei} dienen aber auch zur Kollektivbezeichnung von vielfach Hervorgebrachtem wie \textit{Mal-erei} oder \textit{Schreib-erei}. Die Suffixe \textit{-er} und \textit{-ei} wurden zu \textit{einem} komplexen Suffix \textit{-erei} reanalysiert und dienen heute direkt der Ableitung von Verben zu Substantiven wie \textit{Nasch-erei} oder \textit{Stehl-erei}. Aus der Kollektivbedeutung entwickelt sich die Iterativbedeutung, die häufige Wiederholung der Ereignisse. Dieser Bedeutungsaspekt liegt, wie \citet[352]{Olsen1991} bemerkt, der negativen Konnotation dieser Derivata zugrunde.


\subsubsubsection{\textit{-bar}}\label{subsubsection:bar}

Mit\is{Suffix!-\textit{bar}} \textit{-bar} werden Adjektive aus Verben abgeleitet. Die zentrale Domäne der \textit{bar}-Derivation sind heute die Derivationsstammformen von agentiven transitiven Handlungsverben (vgl. Abschnitt~\ref{subsection:Verbklassen}). Daher gibt es \textit{trink-bar}, \textit{ess-bar} und \textit{liefer-bar} aber nicht *\textit{sitz-bar}, *\textit{schwitz-bar} oder *\textit{bekomm-bar}.\footnote{Das produktive Muster muss von älteren Bildungen wie \textit{frucht-bar} oder \textit{schein-bar} unterschieden werden. Diese gehen auf die Herkunft von \textit{-bar} aus einem Verbaladjektiv mit der Bedeutung \gq{tragend} zurück. Der entsprechende Stamm ist heute noch in \textit{Bahre} und \textit{gebären} (\gq{austragen}) erhalten.} Den mittels \textit{-bar} gebildeten Eigenschaftsausdrücken wohnt die vom Verb ererbte Dynamik inne. Diese besteht innerhalb eines Möglichkeitsbegriffes: Was \textit{trinkbar} ist, kann getrunken werden. Der Träger jener Eigenschaft entspricht wie in (\ref{ea:lieferbar}) dem Vorgangsträger der sogenannten persönlichen Passivvariante (\ref{ex:wird_geliefert}) (vgl. Abschnitt~\ref{subsection:Passiv}).

\ea \label{ea:Eigenschaft_Vorgangsträger}
\ea \label{ea:lieferbar} eine lieferbare Ware
\ex \label{ex:wird_geliefert} eine Ware, die geliefert werden kann
\z 
\z 

Auf Bildungen wie \textit{brenn-bar} oder \textit{verrott-bar}, die nicht auf einem transitiven Verb als Basis beruhen, gehen wir im Abschnitt~\ref{subsection:Restriktionen} ein.

Die mit \textit{-bar} gebildeten Adjektive können durch \textit{-keit} weiter zu Substantiven wie \textit{Lieferbar-keit} abgeleitet werden. Wie bereits gesehen, wird damit die Eigenschaft von ihrem Träger abstrahiert und zum Redegegenstand.



\subsubsubsection{\textit{-er}}\label{subsubsection:er}

Mit dem\is{Suffix!-\textit{er}} Suffix \textit{-er} werden primär sogenannte\is{Nomen agentis} Nomina agentis als Individuenbezeichnungen von Verben abgeleitet. Es handelt sich um die in der Valenz von Handlungs- und Tätigkeitsverben angelegte Rolle des Agens: \textit{Les-er}, \textit{Schreib-er}, \textit{Hör-er}. Vorgangs- und Zustandsverben sind in der Regel von der \textit{er}-Ableitung ausgeschlossen: *\textit{Wächs-er} bzw. *\textit{Wachs-er}, ?\textit{Schwitz-er}. Hierauf kommen wir im Abschnitt \ref{subsection:Transparenz} zurück. 

Ähnlich weit\is{Nomen instrumenti} wie die Ausfächerung der semantischen Rolle Agens (vgl. Abschnitt~\ref{subsection:Verbklassen}) ist, verhält es sich auch mit dem Nomen agentis. Mit Bildungen wie \textit{Bohr-er} und \textit{Öffn-er} bezeichnen wir in der Regel Instrumente, weshalb jedoch nicht angenommen werden sollte, dass es sich um ein anderes Suffix und um ein anderes Bildungsmuster als bei \textit{Leser} und \textit{Schreiber} handelt. Von der Schablone des handelnden Menschen wird die Instrumentlesart agentiv abgeleitet: so wie der \textit{Leser} etwas liest, so bohrt auch der \textit{Bohrer} etwas. Immer wäre auch die Lesart als belebtes Nomen agentis möglich. 

Eine weitere\is{Nomen acti} Lesart von deverbalen \textit{er}-Ableitungen ist die als Nomen acti. Mit Bildungen wie \textit{Seufz-er}, \textit{Hüpf-er} oder \textit{Lach-er} wird das sequenzhafte Resultat des Geschehens ausgedrückt. Kontextlos ist auch hier die primäre Lesart als Nomen agentis möglich. Diese verweist auf Individuen, was im übertragenen Sinne auch die Sequenzlesart gegenüber dem nicht sequenzierten Seufzen, Hüpfen oder Lachen macht.

\subsubsubsection{\textit{-ling}}\label{subsubsection:ling}

Mit\is{Suffix!-\textit{ling}} dem Suffix \textit{-ling} wurden primär maskuline Substantive von Adjektiven abgeleitet. Derivata wie \textit{Schön-ling}, \textit{Fremd-ling}, \textit{Schwäch-ling}, \textit{Winz-ling} bezeichnen den Träger der Eigenschaft, die das Adjektiv bedeutet. 

Seltener sind Verbstämme durch \textit{-ling} abgeleitet wie in \textit{Prüf-ling}, \textit{Lehr-ling}, \textit{Säug-ling} und \textit{Impf-ling}. Anders als bei der \textit{er}-Ableitung wird der ursprüngliche Handlungsgegenstand als Vorgangsträger perspektiviert: ein \textit{Prüfling} wird geprüft, ein \textit{Impfling} wird geimpft und ein \textit{Säugling} gesäugt. Allerdings kann auch wie bei \textit{Lehr-ling}, \textit{Flücht-ling} oder \textit{Emporkömm-ling} eine agentive Lesart vorliegen. Diese ist typisch bei intransitiven Verbbasen.\footnote{Ein \textit{strebsamer Prüfling} kann nur Prüfling und strebsam zugleich sein. Primär gilt das auch bei \textit{strebsamer Lehrling}, wobei wir \textit{Lehrling} als jemanden verstehen, der gelehrt wird. Unter \textit{Lehrling} verstehen wir heute aber primär jemanden, der einen Beruf lernt. Diese agentive Bedeutung bedingt die Lesart von \textit{strebsamer Lehrling} als jemand, der strebsam lernt.}

\subsubsubsection{\textit{-in}}\label{subsubsection:in}

Zur femininen Sexuskennzeichnung können \textit{er}-Ableitungen, die Menschen bezeichnen, weiter durch das Suffix\is{Suffix!-\textit{in}} \textit{-in} deriviert werden. Aus \textit{Leser} wird \textit{Leser-in}. Diese Derivation bezeichnet man als\is{Movierung} \textit{Movierung}. Denn hier wird ein maskulines Substantiv ins feminine Genus mit Sexuswert verschoben (lat. \textit{movere} \gq{bewegen}). Die Movierung ist natürlich nicht auf \textit{er}-Ableitungen beschränkt, sondern gilt für Substantive, die Personen bezeichnen: \textit{Arzt} und \textit{Ärzt-in}, \textit{Dieb} und \textit{Dieb-in}. Nicht möglich ist die Movierung von Personenbezeichnungen, die wie die Oberbegriffe \textit{Mensch}, \textit{Person} oder \textit{Kind} nicht über das Merkmal Sexus verfügen. Eine Sexusdifferenzierung kann in diesen Fällen durch Unterbegriffe wie \textit{Mann}/\textit{Frau} oder \textit{Junge}/\textit{Mädchen} lexikalisch angelegt sein.

Rein morphologisch\is{generisches Maskulinum} betrachtet, spricht die Movierung dafür, ihre Basis, also \zB \textit{Arzt} und \textit{Leser} als sexusneutral zu verstehen. Dem liegt die Auffassung vom generischen Maskulinum zugrunde. Denn ein hinzugefügtes Suffix überschreibt nicht die Bedeutung der Basis, sondern fügt einen Bedeutungsaspekt hinzu. Demzufolge ist die Basis sexusneutral, die \textit{in}-Ableitung sexusmarkiert. In diesem Sinne ist die Bedeutung von \textit{Arzt} und \textit{Leser} als \gq{Person} zu beschreiben, die einen bestimmten Beruf hat oder eine bestimmte Tätigkeit ausübt.

Aber hier muss zwischen morphologischem System, Semantik und Sprachgebrauch unterschieden werden. In dem Moment, wo Minimalpaare wie \textit{Arzt} und \textit{Ärztin}, \textit{Leser} und \textit{Leserin} vermehrt gebraucht werden, verstehen wir die Opposition als männlich versus weiblich. Ein Indiz für die maskuline Sexusmarkierung der \textit{er}-Ableitung sind entsprechende Ableitungen von femininen sexusmarkierten Bildungen wie \textit{Witwer} von \textit{Witwe} und \textit{Hexer} von \textit{Hexe}.\footnote{Eine Übersicht zum Thema gendergerechte Sprache und generisches Maskulinum als linguistisches Thema bietet \citet{Diewald2018}.} Das Problem der Überschreibung wäre theoretisch zu lösen, wenn angenommen wird, dass auch \textit{-erin} analog zu anderen Reanalysen wie \textit{-igkeit} oder \textit{-erei} als \textit{ein} genus- und sexusmarkierendes Suffix reanalysiert wurde.

Teilweise scheint \textit{-ling} über eine maskuline Sexusmarkierung zu verfügen. Hierfür sprechen \zB die Bildungen \textit{Jüngling} (von \textit{jung}) oder \textit{Schönling} im Gegensatz zu sexusneutralen deverbalen Bildungen wie \textit{Säugling} und \textit{Impfling}. \citet[145, 152]{EisenbergSayatz2002} halten Movierungen von \textit{ling}-Derivata wie \textit{Feiglingin}, \textit{Lehrlingin} oder \textit{Neulingin} einfach für ungrammatisch. \citet[237]{FleischerBarz2012} werten sie als Gelegenheitsbildungen. Wir sehen hier eine durch den Gender-Diskurs motivierte Systemabweichung durch analogen Sprachgebrauch. Das ursprünglich sexusneutrale \textit{-ling} erhält durch die Movierung mit \textit{-in} einen weiblichen Sexuswert. Von diesem wird auf einen männlichen Sexuswert von \textit{-ling} geschlossen.

Ableitungen von Nomina agentis wie \textit{Schreiber} zu \textit{Schreiber-ling} oder \textit{Denker} zu \textit{Denker-ling} sind negativ konnotiert im Sinne von \gq{untalentiert}. Mögen \textit{-ling} und \textit{-in} ursprünglich gleichrangig auf das Suffix \textit{-er} folgen wie in \textit{Schreiber-ling} oder \textit{Schreiber-in}, so folgt im momentanen Sprachgebrauch das Suffix \textit{-in} in der Hierarchie auch auf \textit{-ling}, wofür die erwähnten Bildungen wie \textit{Neuling-in} oder \textit{Häftling-in} sprechen. Folglich bilden heute nicht \textit{Schreiber-ling} und \textit{Schreiberin} sexusinterpretierte Oppositionspaare, sondern \textit{Schrei-ber} und \textit{Schreiber-in} sowie \textit{Schreiber-ling} und selten gebrauchte Bildungen wie \textit{Schreiberling-in}.

\subsubsubsection{\textit{-schaft}, \textit{-tum}}\label{subsubsection:schaft}

Weitere Ableitungen sind durch die Kollektiva bildenden Suffixe\is{Suffix!-\textit{schaft}}  \textit{-schaft} und\is{Suffix!-\textit{tum}}  \textit{-tum} möglich: \textit{Lehrer-schaft}, \textit{Lehrerinnen-schaft} und \textit{Streber-tum}, \textit{Streberinnen-tum}. Die Abfolge der Suffixe zeigt ihre jeweilige hierarchische Position, die wir in der Abbildung~\ref{Abb:Lehrerinnenschaft} darstellen. 


	\begin{figure}
		\centering
	\begin{forest}
	[Stammform \\ Substantiv Fem \\ \textit{Lehrerinnenschaft} [Derivationsstammform \\ Substantiv Fem \\\textit{lehrerinnen} [Stammform \\ Substantiv Fem \\ \textit{Lehrerin} [Derivationsstammform \\ Substantiv Mask \\ \textit{lehrer} [Stammform \\ Substantiv Mask \\ \textit{Lehrer} [Derivationsstammform \\ Verb \\ \textit{lehr}] [Suffix \\ bildet Substantiv Mask \\ \textit{-er}] ] ] [Suffix \\ bildet Substantiv Fem \\ \textit{-in}]  ] ] [Suffix \\ bildet Substantiv Fem \\ \textit{-schaft}]  ]
	\end{forest}
	\caption{Lehrerinnenschaft}\label{Abb:Lehrerinnenschaft}
\end{figure}

Neben dem Kollektivausdruck wird mit \textit{schaft}-Ableitungen auch die Beziehung zwischen Personen (eines Kollektivs) ausgedrückt: \textit{Freund-schaft}, \textit{Feind-schaft}, \textit{Kamerad-schaft}. Die Ableitung von Herrscherbezeichnungen mit \textit{-tum} dient dem Ausdruck des Herrschaftsgebietes wie in \textit{Kaisertum} oder \textit{Fürstentum}. Die meisten bestehenden Derivata auf \textit{-schaft} und \textit{-tum} beruhen auf heute nicht mehr produktiven Mustern, so \zB \textit{Leidenschaft}, \textit{Liebschaft} und \textit{Schwangerschaft} sowie \textit{Altertum}, \textit{Reichtum} und\is{Suffigierung!Baumdiagramm} \textit{Siechtum}.


\subsubsubsection{\textit{-chen}, \textit{-lein}}\label{subsubsection:chen}

Von Substantiven\is{Diminutivbildung} lassen sich durch die Suffixe\is{Suffix!-\textit{chen}} \textit{-chen} und\is{Suffix!-\textit{lein}} \textit{-lein} Diminutivformen bilden. Diese Möglichkeit besteht bei Konkreta wie \textit{Mensch-lein}, \textit{Pferd-chen} oder \textit{Häus-chen} und bei Abstrakta wie \textit{Idee-chen} oder \textit{Ängst-chen}. Aus der verkleinernden Bedeutung haben sich wie bei \textit{Kätz-chen} und \textit{Freund-chen} bestimmte Konnotationen stabilisiert, das heißt emotional positive oder negative Mitbedeutungen entwickelt. 

Bei den sogenannten unzählbaren Stoff- und Substanzbezeichnungen ist die Diminutivbildung nur möglich, wenn eine zählbare Interpretationsmöglichkeit im Sinne von Mengen besteht. So bezeichnet man mit \textit{Bier-chen} ein kleines Glas oder eine kleine Flasche Bier und mit \textit{Hölz-chen} ein kleines Stück Holz. In der Suffixhierarchie stehen \textit{-chen} und \textit{-lein} ganz am Ende, weshalb weder *\textit{Leserchen-schaft} noch *\textit{Heldchen-tum} möglich sind.\footnote{Eine Ausnahme bilden \textit{Frauchen-schaft} und \textit{Herrchen-schaft} mit den lexikalisierten Bildungen \textit{Herrchen} und \textit{Frauchen} für Hundehalter und Hundehalterinnen.} Bildungen wie \textit{Heldentüm-chen} oder \textit{Leserschaft-chen} sind hingegen grundsätzlich möglich.

Bevor wir im nächsten Abschnitt~\ref{subsubsection:SuffigierungVerb} kurz auf den Randbereich der Suffigierung im verbalen Bereich eingehen, diskutieren wir im folgenden Exkurs~\ref{subsubsection:DerivationoderKomposition} den Status von Bildungen wie \textit{Dickhäuter} und \textit{grünäugig}. Handelt es sich um Komposita \textit{Dick-häuter} bzw. \textit{grün-äugig}? Oder geben wir die binäre Strukturiertheit auf und kombinieren Komposition und Derivation in einem Schritt \textit{Dick-häut-er} bzw. \textit{grün-äug-ig}? Oder haben wir es mit einer Art Derivation der Wortgruppen \textit{dicke Haut} bzw. \textit{grüne Augen} zu tun? Vor dem Hintergrund verschiedener Vorschläge werden wir einen anderen Lösungsansatz aufzeigen.

\begin{Exkurs}{Derivation oder Komposition?}\label{subsubsection:DerivationoderKomposition} 

Einen Sonderfall stellen Bildungen wie \textit{Dickhäuter} oder \textit{grünäugig} dar. \citet{Leser1990} analysiert sie als Komposita mit der Besonderheit, dass \textit{Häuter} (außer als Nomen agentis) und \textit{äugig} alleine nicht gebraucht werden. Bei \textit{Häuter} und \textit{äugig} handelt es sich um Derivata aus \textit{Haut} und \textit{Auge}. Dass die Derivata alleine nicht wortfähig sind, lässt sich pragmatisch begründen. Da Lebewesen und insbesondere Menschen Augen haben, mangelt es der Eigenschaft \textit{äugig} an Relevanz, um den entsprechenden Eigenschaftsträger zu beschreiben. Diese Eigenschaft kommt typischerweise allen Lebewesen zu. Ähnlich sieht es bei \textit{Häuter} als Träger von Haut aus. Erst die nähere Bestimmung erlaubt relevante Eigenschaftszuschreibungen wie \textit{grünäugig} und analog \textit{kurzbeinig}, \textit{langarmig} oder eben \textit{Dickhäuter}. 

\citet[61--62]{Meibauer_et_al_2015} schlagen die Aufgabe der Binärstruktur vor und analysieren \textit{grünäugig} und \textit{Dickhäuter} als gleichzeitige Verbindung von \textit{grün} mit \textit{äug-} und \textit{-ig}. Letztlich geschehen Komposition und Derivation in einem Schritt, was unbefriedigend ist. 

Die Aufgabe der binären Strukturiertheit ist keineswegs zwingend. Denn es kann, wie in der Konstruktionsgrammatik von \citet[128]{Booij2005} vorgeschlagen, ein dreigliedriges Bildungsschema mit interner Binarität wie in (\ref{ea:grünäugig}) angenommen werden: 

\ea \label{ea:grünäugig}[Adjektiv [Substantiv\textit{-ig}]\textsubscript{Adjektiv}]\textsubscript{Adjektiv}
\z 

Die Derivation eines Substantivs wie \textit{Auge} zum Adjektiv \textit{äugig} bildet das Zweitglied eines Kompositums, dessen Erstglied im nächsten Schritt durch ein Adjektiv besetzt wird. 

Anders analysiert \citet[93]{Donalies2005WbÜb} diese Bildungen. Sie geht nicht von Derivation und Komposition aus, sondern von der Derivation der syntaktischen Strukturen \textit{grüne Augen} zu \textit{grünäugig} und \textit{dicke Haut} zu \textit{Dickhäuter}. Hierfür spricht die Semantik. Denn ein \textit{Dickhäuter} trägt dicke Haut, ist aber anders als der \textit{Dickwanst} als dicker Wanst zur Bezeichnung seines Trägers kein dicker Häuter. Unbefriedigend an dieser Analyse ist, dass es in \textit{grünäugig} und \textit{Dickhäuter}, wie es für die Wortbildung typisch ist, keine Flexionsformen gibt. Diese sollten bei einer wirklich syntaktischen Basis jedoch vorhanden sein.

Wir schlagen deshalb im Rahmen unseres wortbasierten Ansatzes eine andere Analyse vor. Bildungen wie \textit{grünäugig}, \textit{Dickhäuter} und \textit{krummbeinig} bestimmen wir als Derivata von Komposita wie \textit{Grünauge} (\ref{ea:Grünauge}), \textit{Dickhaut} (\ref{ex:Dickhaut}) oder \textit{Krummbein} (\ref{ex:Krummbein}).

\ea \label{ea:KompositaAdjSubst}
\ea \label{ea:Grünauge} Was das 'Grünauge' (sie) und das 'Blauauge' (ihn) verbanden, drückt sich aus in etlichen Exponaten.\footnote{Berliner Woche, 29.12.2015.}
\ex \label{ex:Dickhaut} Der Outdoor-Hersteller mit der prähistorischen Dickhaut im Logo will näher an die Branchenführer heranrücken.\footnote{Tages-Anzeiger, 19.12.1998.}
\ex \label{ex:Krummbein} Hund, Hundi, Zamperl, Dackerl, Waldi [\ldots{}] Kann sich kein Mensch mehr vorstellen, daß so ein goldiges Krummbein vom Wolf abstammen soll.\footnote{Süddeutsche Zeitung, 14.04.1999.}
\z 
\z 

Diese meist als Possessivkomposita (vgl. Abschnitt~\ref{subsubsection:PossKomp}) gebrauchten Substantive, die auch eine Quelle von Familiennamen waren, werden ganz regulär mit dem entsprechenden Suffix abgeleitet. Die so entstandenen Bildungen \textit{grünäugig} oder \textit{Dickhäuter} können durchaus im Sinne des Bildungsschemas in (\ref{ea:grünäugig}) reanalysiert worden sein. In diesem Sinne muss die Komposition nicht (mehr) der Derivation vorausgehen, aber die Derivation bleibt an eine entsprechende Komposition gebunden.

\end{Exkurs}


\subsubsection{Suffigierung im verbalen Bereich}\label{subsubsection:SuffigierungVerb}

Im\is{Suffigierung!verbal} verbalen Bereich, d.\,h. zum Zwecke der Verbbildung, sind Suffixe heute marginal und haben sich grundsätzlich aus dem nominalen Bereich entwickelt.\footnote{Bis ins Althochdeutsche spielten Suffixe im verbalen Bereich eine wichtige Rolle. Allerdings sind die meisten dieser Suffixe wie \textit{-er} oder \textit{-el} nicht-verbalen Ursprungs: \textit{besser} zu \textit{bessern}, \textit{Acker} zu \textit{ackern} (Konversion; vgl. Abschnitt~\ref{subsection:Konversion}), aber \textit{Folge} zu \textit{folg-ern} (Suffigierung via Analogie), \textit{Nagel} zu \textit{nageln}, \textit{Angel} zu \textit{angeln} (Konversion), aber \textit{klug} zu (\textit{aus})\textit{klüg-eln} (Suffigierung). Cf. \citet[91--101]{Wilmanns1896Wb}.}

Das aus dem Französischen entlehnte Suffix\is{Suffix!-\textit{-ier}(\textit{en})} \textit{-ier}(\textit{en}) dient wie in \textit{telefon-ier}(\textit{en}), \textit{aktiv-ier}(\textit{en}) der Verbalisierung von entlehnten substantivischen und adjektivischen Stämmen oder von Konfixen wie in \textit{fris-ier}(\textit{en}). In Analogie kam es zur Ableitung heimischer Basen wie bei \textit{halb-ier}(\textit{en}) und \textit{gast-ier}(\textit{en}). Heute ist das Suffix -\textit{ier}(\textit{en}) nur noch im Sprachspiel produktiv bzw. reaktivierbar \textit{googlieren} statt \textit{googeln}.

Das Suffix\is{Suffix!-\textit{el}(\textit{en})} \textit{-el}(\textit{n}) dient der Bildung einer diminutiven oder iterativen Variante zu einem bestehenden Verb: \textit{läch-el}(\textit{n}) von \textit{lach-}(\textit{en}), \textit{hüst-el}(\textit{n}) von \textit{hust-}(\textit{en}), \textit{streich-el}(\textit{n}) von \textit{streich-}(\textit{en}) usw.\footnote{Ursprünglich ist das (\textit{e})\textit{l} Bestandteil von Substantiven wie \textit{Nagel}, \textit{Riegel}, die ohne Affix verbalisiert wurden: \textit{nageln}, (\textit{ver})\textit{riegeln}. Später wurde es als Verbalisierer empfunden und in Verbindung mit dem substantivischen Diminutiv-Suffix \textit{-l} gebracht: \textit{Brett} zu \textit{Brettl} oder \textit{Christkind} zu \textit{Christkindl} bzw. \textit{Christkindel}. Cf. \citet[96--98]{Wilmanns1896Wb}.} Eine weitere Funktion besteht in der lautmalerischen Verbbildung wie bei \textit{bimmel}(\textit{n}), \textit{lispel}(\textit{n}) oder \textit{nuschel}(\textit{n}).\footnote{Cf. \citet[431]{FleischerBarz2012}.} Viele dieser Bildungen lassen sich heute nicht mehr sinnvoll analysieren, da die Basis außerhalb dieser Verben nicht (mehr) existiert. Bei bestehenden substantivischen Stämmen wie bei \textit{witz-el}(\textit{n}), \textit{fröst-el}(\textit{n}) oder \textit{sächs-el}(\textit{n}) ist die Analyse problemlos möglich. Das gilt auch für deadjektivische Bildungen wie \textit{blöd-el}(\textit{n}) oder \textit{schwäch-el}(\textit{n}).

Ebenfalls als Ableitung von der nominalen Ebene\footnote{Das Suffix \textit{-er} bildet adjektivische Komparativformen, die ohne Affix verbalisiert wurden: \textit{besser} zu \textit{besser}(\textit{n}) oder \textit{mehr} zu \textit{mehr}(\textit{en}). Ebenso war es Bestandteil von Substantiven, die ohne Affix verbalisiert wurden: \textit{Acker} zu \textit{acker}(\textit{n}) sowie ein Pluralflexiv wie bei \textit{Löcher}, das ohne Affix zu \textit{löcher}(\textit{n}) konvertiert wurde. Ausgehend von solchen Konversionen wurde es als verbales Suffix empfunden und gebraucht. Cf. \citet[92--93]{Williams1981}.} wurde das Suffix\is{Suffix!-\textit{er}(\textit{n})} \textit{-er} verbbildend gebraucht, \zB für lautmalerische Iterativbildungen wie \textit{bibber}(\textit{n}), \textit{knabber}(\textit{n}) und \textit{schlitter}(\textit{en}) sowie für kausative Verben wie \textit{einschläf-er}(\textit{n}) von \textit{einschlaf-}(\textit{en}), \textit{steig-er}(\textit{n}) von \textit{steig-}(\textit{en}) und \textit{folg-er}(\textit{n}) von \textit{folg-}(\textit{en}).

Das Suffix\is{Suffix!-\textit{ig}(\textit{en})} \textit{-ig} bildet, wie in Abschnitt~\ref{subsubsection:ig} dargestellt, primär Adjektive aus Substantiven. Adjektive wie \textit{kräft-ig} oder \textit{stein-ig} wurden ohne Formsignal verbalisiert: \textit{kräftig}(\textit{en}), \textit{steinig}(\textit{en}). Um diese Art der Ableitung, um Konversion, geht es im Abschnitt~\ref{subsection:Konversion}. Bereits im Mittelhochdeutschen wurde das Suffix \textit{-ig} auch direkt als Verbalisierer gebraucht wie bei \textit{ängst-ig}(\textit{en}) und \textit{pein-ig}(\textit{en}), ohne dass es entsprechende Adjektive als Zwischenstufe gegeben hätte. Später verblasste die verbalisierende Kraft des Suffixes und es wurde nur noch zusammen mit Präfixen gebraucht, um die es im folgenden Abschnitt~\ref{subsection:Präfigierung} gehen wird. Es kam zu Parallelbildungen wie \textit{be-rein-ig}(\textit{en}) neben \textit{rein-ig}(\textit{en}), \textit{be-sänft-ig}(\textit{en}) neben dem heute nicht mehr gebräuchlichen \textit{sänft-ig}(\textit{en}), \textit{be-gnad-ig}(\textit{en}) eben \textit{be-gnad}(\textit{en}) (heute adjektivisch \textit{begnadet}). Es wurden aber auch ganz neue Verben wie \textit{be-grad-ig}(\textit{en}) ohne entsprechendes Verb *\textit{gradig}(\textit{en}) gebildet.\footnote{Die Tilgung des Schwa-Lautes (\textit{grad} statt \textit{gerade}) beruht auf einem phonologischen Prozess, den wir in Abschnitt~\ref{sec:Schwatilgung} beschreiben.} Das Präfix und das Suffix wurden als Zirkumfix \textit{be-ig} reanalysiert. Die Konstituentenstruktur\is{Zirkumfigierung!Baumdiagramm} stellen wir in der Abbildung~\ref{Abb:begradigen} dar.

\begin{figure}
	\centering
\begin{forest}
	[Stammform \\ Verb \\ \textit{begradig-} [Zirkumfix \\bildet Verb \\ \textit{be-}] [Derivationsstammform \\ Adjektiv \\ \textit{grad}] [Zirkumfix \\ bildet Verb \\ \textit{-ig}] ]
\end{forest}
\caption{Zirkumfixbildung \textit{begradigen}}\label{Abb:begradigen}
\end{figure}
 
         
\subsection{Präfigierung}\label{subsection:Präfigierung}

Im verbalen Bereich, d.\,h. zum Zwecke der Verbbildung, ist die\is{Präfigierung!verbal} Präfigierung so zentral wie die Suffigierung im nominalen. Präfixe (vgl. Abschnitt~\ref{subsection:Derivationsaffixe}) gehen historisch auf Adverbien und aus diesen entstandene Präpositionen zurück. Heute handelt es sich bei Präfixen wie bei Suffixen um unbetonte Affixe, also um Operatoren über Stammformen. Während wir einen Ursprung substantivischer Komposita auf eine Attribut-Kopf-Beziehung in der NP (\textit{des Diebes Gut} mit \textit{Gut} als Kopf) zurückgeführt haben, handelt es sich bei verbalen Präfixen ursprünglich um eine adverbiale Beziehung in der VP (\textit{hinauf steigen} mit \textit{steigen} als Kopf und \textit{auf-steigen}).   

Auch Präfixverben\is{Präfigierung!Binarität} sind binär strukturiert. Sie bestehen aus einer lexikalischen Basis und einem Präfix. Die lexikalische Basis geht als Derivationsstammform in die Präfigierung ein. Bei Präfixverben wie \textit{er-kält}(\textit{en}) oder \textit{ent-kräft}(\textit{en}) unterscheidet sich die Derivationsstammform von den lexikalischen Stammformen \textit{kalt} und \textit{Kraft}.  

Die Basen präfigierter Verben sind Verben wie bei \textit{be-lüg}(\textit{en}), Substantive wie bei \textit{ent-kern}(\textit{en}) und Adjektive wie bei \textit{er-stark}(\textit{en}). Im Zuge der Präfigierung entstehen Verben aus anderen Wortarten oder bestehende Verben werden in ihrer Bedeutung und meist in ihrer Valenz abgewandelt. Das zeigen wir in dem folgenden Abschnitt~\ref{subsubsection:be_Präfix} exemplarisch anhand von Präfigierungen mit \textit{be-} auf. Im Anschluss gehen wir auf die Telisierung von Verben durch Präfixe ein (Abschnitt~\ref{subsubsection:Perfektivierung_Präfix}), argumentieren dafür, Präfixe trotz möglicher Verbalisierungen nicht als morphologische Köpfe zu bestimmen (Abschnitt~\ref{subsubsection:Präfix_Kopf}), stellen eine wesentliche Eigenschaft einzelner Präfixe dar, ganze Reihen semantisch und valenziell ähnlicher Präfixverben zu bilden (Abschnitt~\ref{subsubsection:Reihe_Motivation_Präf}) und beleuchten ein Erbe der präpositionalen Herkunft vieler Präfixe (Abschnitt~\ref{subsubsection:Präfix_Präposition}). Abschließend werfen wir einen Blick auf einen Randbereich, nämlich auf Präfigierungen im nominalen und adjektivischen Bereich (Abschnitt~\ref{subsubsection:Präfix_Nomen}).   

 

\subsubsection{Exemplarisch:  \textit{be}-Präfigierung}\label{subsubsection:be_Präfix}
\largerpage
Die valenzschaffende und valenzverändernde\is{Präfigierung!Valenzdynamik} Kraft betrachten wir im Folgenden am Beispiel des\is{Präfix!\textit{be}-} Präfixes \textit{be-}. Dessen Leistung besteht darin, transitive Handlungsverben mit einem Akkusativobjekt als Handlungsgegenstand (Patiens) hervorzubringen. Die Basis kann wie bei \textit{lügen} in (\ref{ea:lügenitr}) intransitiv oder wie bei \textit{beruhigen} in (\ref{ex:beruhigen}) gar kein Verb (\ref{ea:ruhig}) sein.  

\ea \label{ea:belügen}
\ea [*] {Sie lügt ihn.}\label{ea:lügenitr}
\ex [] {Sie belügt ihn.}\label{ex:belügenAkk}
\z 
\z 

\ea \label{ea:ruhig_beruhigen}
\ea \label{ea:ruhig} ruhig\textsubscript{Adjektiv}
\ex \label{ex:beruhigen} Sie beruhigt ihn.
\z 
\z 

Auch mehrwertige Verben werden durch \textit{be-} zu transitiven Verben abgeleitet. Dies gilt für Verben mit Dativ wie \textit{dienen} in (\ref{ea:be_dienen}), für Verben mit einem Direktivum wie \textit{steigen} in (\ref{ea:be_steigen}), für ditransitive Verben mit Dativ- und Akkusativobjekt wie \textit{schenken} in (\ref{ea:be_schenken}) und für Verben mit Akkusativobjekt und Direktivum wie \textit{laden} in (\ref{ea:be_laden}). 

\ea \label{ea:be_dienen}
\ea \label{ea:dienenDat} Er dient dem Gast.
\ex \label{ex:bedienenAkk} Er bedient den Gast.
\z 
\z 

\ea \label{ea:be_steigen}
\ea \label{ea:steigenauf} Er steigt auf den Berg.
\ex \label{ex:besteigenAkk} Er besteigt den Berg.
\z 
\z 

\ea \label{ea:be_schenken}
\ea \label{ea:schenkenDatAkk} Er schenkt dem Jubilar Blumen.
\ex \label{ex:beschenkenAkk} Er beschenkt den Jubilar (mit Blumen).
\z 
\z 

\ea \label{ea:be_laden}
\ea \label{ea:laden} Er lädt Heu auf den Wagen.
\ex \label{ex:beladen} Er belädt den Wagen (mit Heu).
\z 
\z 

Durch die\is{Präfigierung!Valenzdynamik} \textit{be}-Präfigierung (b-Sätze) wird das ursprüngliche Ziel wie in (\ref{ea:steigenauf}) und (\ref{ea:laden}) oder der ursprüngliche Betroffene wie in (\ref{ea:dienenDat}) und (\ref{ea:schenkenDatAkk}) als Handlungsgegenstand in Perspektive gebracht. 

Diese Sicht ist durchaus umstritten. \citet[41--43]{Eroms1980} \ua nehmen an, dass es sich um reine Variation der Form und der Reihenfolge (Informationsstruktur) bei gleichbleibender Bedeutung handelt. Andere Autoren wie \citet[70--71, 189--190]{Welke1988} und \citet[268]{Eisenberg2020_Wort} gehen von einer Bedeutungsveränderung (gegenüber der Grundvalenz) aus.

Wir schließen uns letzterer Sicht an. Von wesentlicher Bedeutung ist die Unterscheidung zwischen einem Ereignis in der Welt und dessen sprachlicher Darstellung. Die a-Sätze und die b-Sätze beschreiben auf unterschiedliche Weise dasselbe Ereignis. Diese unterschiedliche Weise betrifft den Kern der Sprache. Die b-Sätze perspektivieren die entsprechenden Ereignisse als Handlungen mit einem Handlungsgegenstand als Akkusativobjekt. Reflexe dieser bedeutungsverändernden Umstrukturierung durch das Präfix \textit{be-} zeigen sich bei weiteren Wortbildungsprozessen. Die \textit{ung}-Ableitungen mit den entsprechenden Genitivattributen in (\ref{ea:be_Verb_ung}) funktionieren nur von den Präfixverben, aber nicht oder nur mit anderer Bedeutung von den einfachen Verben.

\ea \label{ea:be_Verb_ung}
\ea \label{ea:Bedienung} die Bedienung des Gastes -- *die Dienung
\ex \label{ex:Besteigung} die Besteigung des Bergs -- die Steigung des Bergs
\ex \label{ex:Beschenkung} die Beschenkung des Jubilars -- die Schenkung des Jubilars
\ex \label{ex:Beladung} die Beladung des Wagens -- die Ladung des Wagens
\z 
\z 

Auf\is{Präfigierung!Valenzdynamik} der semantischen und formalen Umstrukturierung beruht auch, dass nur die \textit{be}-Verben bezüglich ihres Akkusativobjektes mit \textit{-bar} zu Adjektiven abgeleitet werden können (\ref{ea:be_Verb_bar}), nicht aber die einfachen Verben oder nur mit anderer Bedeutung.

\ea \label{ea:be_Verb_bar}
\ea \label{ea:bedienbar} ein bedienbarer Gast -- *ein dienbarer Gast
\ex \label{ex:besteigbar} ein besteigbarer Berg -- *ein steigbarer Berg
\ex \label{ex:beschenkbar} ein beschenkbarer Jubilar -- ein (jemandem) schenkbarer Jubilar
\ex \label{ex:beladbar} ein beladbarer Wagen -- ein (irgendwohin) ladbarer Wagen
\z 
\z 

Aus dem gleichen Grund erlauben nur die \textit{be}-Verben durch das sogenannte persönliche Passiv (\ref{ea:be_Verb_Passiv}) die Veränderung der Handlungs- zur Vorgangsbedeutung und damit die Darstellung des ursprünglichen Handlungsgegenstands als Vorgangsträger (vgl. Abschnitt~\ref{subsection:Passiv}).

\ea \label{ea:be_Verb_Passiv}
\ea \label{ea:wirdbedient} der Gast wird bedient -- *der Gast wird gedient
\ex \label{ex:wirdbestiegen} der Berg wird bestiegen -- *der Berg wird gestiegen
\ex \label{ex:wirdbeschenkt} der Jubilar wird beschenkt -- der Jubilar wird (jemandem) geschenkt
\ex \label{ex:wirdbeladen} der Wagen wird beladen -- der Wagen wird (irgendwohin) geladen
\z 
\z 

Die \textit{be-}Verben drücken\is{Präfix!\textit{be}-} Handlungen aus, die auf eine Entität gerichtet sind. Die jeweiligen Paraphrasen können sich in Abhängigkeit von der lexikalischen Basis unterscheiden. Die Verben \textit{bedachen}, \textit{besohlen}, \textit{beflaggen} drücken aus, dass jemand eine Entität mit dem versieht, was das Substantiv bedeutet. Bildungen wie \textit{bemuttern} und \textit{bewirten} drücken aus, dass jemand jemanden in der Rolle des vom Substantiv Bedeuteten behandelt. \textit{Be}-Verben mit adjektivischer Basis wie \textit{begrünen}, \textit{beruhigen}, \textit{bestärken} oder \textit{befreien} drücken aus, dass jemand eine Entität in die vom Adjektiv bedeutete Eigenschaft versetzt.

In dem folgenden Abschnitt~\ref{subsubsection:Perfektivierung_Präfix} gehen wir auf die Telisierung von Verben ein, zu der es häufig durch Präfigierung kommt, und zeigen, warum eine Äußerung wie \textit{die verblühte Blume} problemlos möglich ist, nicht aber *\textit{die geblühte Blume}.


\subsubsection{Telisierung durch Präfigierung}\label{subsubsection:Perfektivierung_Präfix}

Eine andere, mit der Transitivierung verwandte Funktion von Präfixen ist die zeitliche Eingrenzung, die\is{Präfigierung!Telisierung} Telisierung einer Geschehensdarstellung. So wird mit dem einfachen Verb \textit{blühen} ein atelisches Geschehen beschrieben, d.\,h. ein Geschehen, das keinen Endpunkt impliziert. Im Unterschied dazu sind die präfigierten Verben \textit{erblühen} und \textit{verblühen} telisch. Mit \textit{erblühen} wird ein Geschehen beschrieben, das ein Blühen-Ereignis einleitet und mit dessen Eintreten endet. Wenn eine Blume erblüht ist, dann erblüht sie nicht mehr. Mit \textit{verblühen} wird ein Geschehen beschrieben, das ein Blühen-Ereignis beendet. Ist jenes Ende erreicht, so endet, was das Verb \textit{verblühen} beschreibt. Wenn eine Blume verblüht ist, dann verblüht sie nicht mehr. 

Durch die Telisierung wird ein Zustandswechsel impliziert, mit dem das Geschehen endet. Das hat weitreichende Konsequenzen. So wird das Perfekt eines atelischen intransitiven Verbs wie \textit{blühen} mit \textit{haben} gebildet (\ref{ea:hat_geblüht}), das von telischen intransitiven Verben wie \textit{ver-} und \textit{erblühen} aber mit \textit{sein} (\ref{ex:ist_erblüht}).

\ea \label{ea:Prefektivierung}
\ea \label{ea:hat_geblüht} Die Blume hat geblüht.
\ex \label{ex:ist_erblüht} Die Blume ist erblüht/verblüht.
\z 
\z 

Nur das Partizip II telischer intransitiver Verben kann als Adjektiv attributiv auf das Subjekt bezogen werden (\ref{ex:Pxblühte_Blume}), bei atelischen Verben ist das nicht möglich (\ref{ea:geblühte_Blume}).

\ea \label{ea:P2_attributiv}
\ea [*] {die geblühte Blume}\label{ea:geblühte_Blume}
\ex [] {die erblühte/verblühte Blume}\label{ex:Pxblühte_Blume}
\z 
\z 

Präfixe können bestehende Valenzen und/oder zeitliche Aspekte der Verbbedeutung verändern. Damit greifen sie in das Zentrum verbaler Bedeutung und Funktion ein. In dem folgenden Abschnitt~\ref{subsubsection:Präfix_Kopf} werden wir aber dafür argumentieren, Präfixe nicht als morphologische Köpfe zu bestimmen.


\subsubsection{Präfixe als Köpfe?}\label{subsubsection:Präfix_Kopf}

Mögliche Verbalisierungen\is{Präfigierung!Verbalisierung} wie bei \textit{bedachen} oder \textit{befeuchten} halten wir für einen nützlichen Nebeneffekt. Aus der valenzverändernden Kraft hat sich eine valenz- und damit verbbildende Kraft (möglicherweise über Zirkumfixe wie in \textit{be-grad-ig}(\textit{en}) vermittelt) entwickelt. Präfixe formen Verben.

Ein Indiz dafür, den Verbalisierungseffekt für sekundär zu halten, ist, dass die Stämme von Präfixverben ohne Formsignal zu Substantiven konvertiert werden können (vgl. Abschnitt~\ref{subsection:Konversion}). Aus \textit{besuch}(\textit{en}) wird \textit{Besuch}, aus \textit{erwerb}(\textit{en}) wird \textit{Erwerb}. Im Gegensatz zu Präfixen dienen Suffixe primär der Veränderung der Wortart. Deshalb lassen sich durch Suffigierung gebildete Adjektive wie \textit{ruhig} oder \textit{deutlich} anders als einfache Adjektive wie \textit{reif} nicht zu Verben konvertieren: *\textit{ruhigen}, *\textit{deutlichen} im Gegensatz zu \textit{reifen}. Mit Präfixen ist die Ableitung komplexer Adjektive zu Verben wie \textit{beruhigen} und \textit{verdeutlichen} allerdings möglich, was \citet[309]{Eisenberg2020_Wort} als Argument für ihren Kopfstatus aufführt. Weitere deadjektivische und desubstantivische Bildungen wie \textit{bewaffnen}, \textit{vernarben} und \textit{zermürben} sind auch für \citet[80--84]{Olsen1985} ein Indiz für ihren Kopfstatus. 


Präfixe haben eine andere Qualität als\is{Präfix!als Kopf?} Suffixe. Typischerweise bilden sie keine Verben, sondern formen diese durch ihre regelhaft valenzverändernde Kraft. Im Randbereich kann ihnen eine verbalisierende Kraft zukommen. Dennoch bestimmen wir sie nicht als morphologische Köpfe. Gegen die Kopfanalyse spricht vor allem die Herkunft der Präfixe von den Adverbien, die ursprünglich, d.\,h. syntaktisch einem verbalen Kopf untergeordnet waren. Präfixe haben keinen Einfluss auf die Flexionsklasse der Präfixverben. Starke Verben bleiben stark und schwache schwach. Die Konstituentenstruktur für das deadjektivische Verb \textit{erkälten} stellen wir in\is{Präfigierung!Baumdiagramm} Abbildung~\ref{Abb:erkälten} dar.

\begin{figure}
	\centering
\begin{forest}
	[Stammform \\ Verb \\ \textit{erkält-} [Präfix\\ formt Verb\\ \textit{er-}] [Derivationsstammform\\ Adjektiv \\ \textit{kält}] ]
\end{forest}
\caption{Konstituentenstruktur von \textit{erkälten}}\label{Abb:erkälten}
\end{figure}

Ist die\is{Präfigierung!Verbalisierung} Basis adjektivisch wie bei \textit{erkälten} oder substantivisch wie bei \textit{bewässern}, so kann sich die Derivationsstammform von der lexikalischen Stammform unterscheiden. Bei verbalen Basen entspricht die Derivations- der lexikalischen Stammform (Flexionsstammform). Somit erklärt sich der Erhalt der Flexionsklasse. Starke Verben wie \textit{laden} (\textit{lud}) bleiben wie \textit{beladen} (\textit{belud}) stark und schwache Verben wie \textit{packen} (\textit{packte}) bleiben wie \textit{bepacken} (\textit{bepackte}) schwach.  

In dem folgenden Abschnitt~\ref{subsubsection:Reihe_Motivation_Präf} gehen wir auf die Eigenschaft von Präfixen ein, ganze Reihen semantisch und valenziell ähnlicher Verben, sogenannte Reihen, zu bilden. 


\subsubsection{Reihenbildung}\label{subsubsection:Reihe_Motivation_Präf}

Für die systematische Beschreibung einzelner Präfixe und ihrer Leistung verweisen wir auf \citet[141--375]{Kuehnhold1973}. Eine kurze Übersicht bieten \citet[383--395]{FleischerBarz2012}. Uns geht es in der folgenden Übersicht (\ref{ea:Präfixreihen}) um Beispiele dafür, dass Präfixe häufig eine oder mehrere\is{Reihenbildung!Präfixverb} Reihen bilden, das heißt Präfixverben die trotz aller Unterschiede bezüglich der Basen, dieselbe Valenz haben und den gleichen Ereignistyp ausdrücken.

\ea \label{ea:Präfixreihen}
\ea \label{ea:be_V} jemand belügt, begafft, bespitzelt, berät, bespuckt\ldots{} jemanden
\ex \label{ex:ent_V} jemand entgrätet, entkalkt, entfärbt, entgiftet\ldots{} etwas
\ex \label{ex:er_V} jemand ertanzt, erschreibt, erlügt, erturnt\ldots{} sich etwas
\ex \label{ex:miss_V} jemand missdeutet, missversteht, missbraucht\ldots{} etwas oder jemanden
\ex \label{ex:ver_V} jemand verschreibt, verspricht, verhört, verläuft\ldots{} sich
\ex \label{ex:zer_V} etwas zerplatzt, zerspringt, zerrinnt, zerbröselt, zerbirst\ldots{} 
\z 
\z 

Einzelne Präfixe und ganze mit ihnen gebildete Reihen können systematisch durch Kontrastbildungen mit anderen Präfixen oder mit Partikeln (vgl. Abschnitt~\ref{section:Vpt Bildung}) motiviert sein, was exemplarisch in (\ref{ea:Kontrastmotivation}) dargestellt wird.

\ea \label{ea:Kontrastmotivation}
\ea \label{ea:be_an} jemanden belügen/anlügen
\ex \label{ex:be_ent} etwas bewässern/entwässern
\ex \label{ex:ver_ent} jemanden oder etwas vergiften/entgiften
\ex \label{ex:er_ver} erblühen/verblühen
\ex \label{ex:zer_be_an} etwas zerkratzen/bekratzen/ankratzen
\z 
\z 

Eingangs haben wir erwähnt, dass sich viele Präfixe aus früheren Präpositionen entwickelt haben. In dem folgenden Abschnitt~\ref{subsubsection:Präfix_Präposition} zeigen wir ein gewisses Erbe dieser Entwicklung auf, grenzen aber dennoch Präpositionen klar von Präfixen ab.

\subsubsection{Präfix versus Präposition}\label{subsubsection:Präfix_Präposition}

Eine\is{Präfix!versus Präposition} besondere Gruppe bilden die Präfixe \textit{unter-} wie in \textit{unterschreiben}, \textit{hinter-} wie in \textit{hintergehen}, \textit{über-} wie in \textit{übersetzen}, \textit{durch-} wie in \textit{durchleuchten} und \textit{um-} wie in \textit{umgarnen}. Ihre Herkunft von Präpositionen ist noch offensichtlich und sie sind abzugrenzen von den gleichförmigen betonten Verbpartikeln (vgl. Abschnitt~\ref{section:Vpt Bildung}). Auch diese Präfixe führen zu einer Begrenzung (Telisierung) der Verbbedeutung und häufig zu transitiven Handlungsverben. Wenn ein transparenter Bezug zur Präposition besteht, können wir von Vererbung des präpositionalen Arguments (\ref{ea:Präp_um}) an das Präfixverb (\ref{ex:Präfix_um}) sprechen. 


\ea \label{umfahren}
\ea \label{ea:Präp_um} Sie fährt um das Hindernis.
\ex \label{ex:Präfix_um} Sie umfährt das Hindernis.
\z 
\z 

Mitunter motiviert erst das Paraphrasemuster (\ref{ea:um_Garten}) die Verbalisierung (\ref{ex:umzäunt}). 

\ea \label{ea:umzäunen}
\ea \label{ea:um_Garten} Sie stellt/baut einen Zaun um den Garten.
\ex \label{ex:umzäunt} Sie umzäunt den Garten.
\z 
\z 

Ein grundsätzlicher Rückbezug auf eine Präpositionalphrase ist nicht möglich. Das wird anhand der Präfixverben in (\ref{ea:überschätzen}) deutlich. Eine präpositionale Paraphrase ist hier nicht möglich (\ref{ex:unterschätzen}).

\ea \label{ea:Pxschätzen}
\ea [] {Sie überschätzt/unterschätzt ihn.}\label{ea:überschätzen}
\ex [*] {Sie schätzt über/unter ihn.}\label{ex:unterschätzen}
\z 
\z 

Nachdem wir uns ausführlich mit Präfixen im Bereich der Verbbildung beschäftigt haben, weisen wir im folgenden Abschnitt~\ref{subsubsection:Präfix_Nomen} auf einen Randbereich hin, nämlich auf die Präfigierung im nominalen und adjektivischen Bereich. 


\subsubsection{Präfigierung im nominalen und adjektivischen Bereich}\label{subsubsection:Präfix_Nomen}

Im nominalen\is{Präfigierung!nominal} Bereich ist die Präfigierung marginal. Präfixe verändern im nominalen Bereich nicht die Wortart. Nominale Präfixe sind folglich keine morphologischen Köpfe. Wie \textit{erz-} in \textit{Erzbetrüger}, \textit{ur-} in \textit{Urzeit} oder \textit{ex-} in \textit{Exfreund} modifizieren sie die Basis oder dienen wie \textit{miss-} in \textit{Misserfolg} der negativen Umwertung der Basisbedeutung.\footnote{Bei \textit{miss-} ist zwischen nominalem Präfix wie in \textit{Misserfolg} oder \textit{Misslaune} und substantivierten Präfixverben wie \textit{Missachtung} von \textit{missachten} zu unterscheiden.}

Mit \textit{un-} können Substantive wie in \textit{Unwort}, \textit{Unart} oder \textit{Unmensch} präfigiert werden. Besonders produktiv ist aber die Ableitung von komplexen Adjektiven, die selbst auf produktiven Bildungen basieren wie bei \textit{untrinkbar}, \textit{unfreundlich} oder \textit{ungeistig}. Dadurch kann das Präfix \textit{un-} auch als Indikator für den adjektivischen Charakter vom Partizip II gelten: \textit{unbewiesen}, \textit{ungewaschen} versus *\textit{ungeschlafen} oder *\textit{ungearbeitet}. Aufgrund der starken Produktivität im Adjektivbereich wird \textit{un-} in einer Bildung wie \textit{Untrinkbarkeit} als Adjektivpräfix von \textit{trinkbar} und nicht als Substantivpräfix von \textit{Trinkbarkeit} analysiert. Ein anderes Indiz für diese Analyse ist die Tatsache, dass manche Substantive mit \textit{-heit} oder \textit{-keit} präferiert von dem mit \textit{un-} abgeleiteten Adjektiv gebildet werden. So ist \textit{Angenehmheit} deutlich seltener und ungewöhnlicher als \textit{Unangenehmheit}. 


\subsection{Implizite Derivation}\label{subsection:ImplDerivation}

Bei\is{Derivation!implizit} Bildungen wie \textit{fällen} aus \textit{fallen}, \textit{senken} aus \textit{sinken} oder \textit{tränken} aus \textit{trinken} wird die lexikalische Basis nicht durch Kombination mit einem Affix abgewandelt, sondern durch den musterhaften Wechsel des Stammvokals durch Umlaut oder durch Ablaut. Diese nicht-kombinatorische Ableitung bezeichnet man als \textit{implizite Derivation}, da die Formveränderung intern stattfindet. Sie ist nicht mehr produktiv.

Aus resultativen\is{Resultativ}, meist intransitiven Verben wurden kausative\is{Kausativ} (\textit{faktitive}) transitive Verben abgeleitet. Besonders deutlich tritt das Muster bei den sogenannten Positionsverben zutage. Aus den positionsbeschreibenden starken Verben wurden kausative schwache Verben abgeleitet, mit denen die Positionsveränderung einer Entität ausgedrückt wird: aus \textit{sitzen} wurde \textit{setzen} abgeleitet, aus \textit{liegen} wurde \textit{legen} abgeleitet und aus \textit{hängen} (\textit{hing}, \textit{gehangen}) wurde \textit{hängen} (\textit{hängte}, \textit{gehängt}) abgeleitet.

Über die Positionsverben hinaus finden sich Resultativa wie \textit{erschrecken} (\textit{erschrak}, \textit{erschrocken}) oder \textit{saugen} und durch implizite Derivation aus ihnen abgeleitete Kausativa wie \textit{erschrecken} (\textit{erschreckte}, \textit{erschreckt}) und \textit{säugen}.

Der Vokalwechsel ist kein originäres Mittel zur Bildung der Kausativa, sondern ein ursprünglich phonologischer Nebeneffekt durch Anpassung an ein folgendes Suffix (vgl. Abschnitt~\ref{subsection:Assimilation}).\footnote{Im Germanischen wurden kausative Verben durch das Suffix \textit{-jan} abgeleitet. Der Umlaut des Stammvokals war ursprünglich eine Anpassung (Hebung) an das folgende /j/. Cf. \citet[61--62]{Wilmanns1896Wb}.} Dieser Nebeneffekt wurde dann analog morphologisch genutzt.

Ebenfalls nicht mehr produktiv ist die\is{Derivation!implizit} implizite Derivation von Adjektiven wie \textit{kurz} oder \textit{krumm} zu kausativen Verben wie \textit{kürzen} oder \textit{krümmen}. Wie \citet[93]{Echenlohr1999} zeigt, lässt sich die umgelautete Form teilweise auf bestehende adjektivische Stammformen wie bei \textit{kurz} auf \textit{kürz-}(\textit{er}) beziehen. Bei \textit{krümm-} ist das jedoch nicht der Fall. Die Bildung ist jedoch gleichermaßen transparent.

Da die\is{Derivation!implizit} implizite Derivation heute nicht mehr produktiv ist, gehen wir im Rahmen des wortbasierten Ansatzes von aufeinander bezogenen Paaren wie \textit{liegen} und \textit{legen} oder \textit{kurz} und \textit{kürzen} aus, ohne dass die Ableitungsrichtung noch eine Rolle spielt. Folglich bezeichnen wir auch die um- oder abgelauteten Stämme nicht als Derivationsstammformen der historisch zugrunde liegenden Formen, sondern als eigenständige lexikalische Stammformen.

Substantive wie \textit{Biss}, \textit{Trank} oder \textit{Wurf} wurden nicht implizit von den verbalen Stammformen \textit{beiß-}, \textit{trink-} oder \textit{werf-} abgeleitet. Die Basis der Ableitung waren entweder die Präteritumstämme wie \textit{biss-} und \textit{trank-} oder heute nicht mehr bestehende Präteritumstämme wie \textit{wurf-}. Wie bei \textit{Trunk} kommt auch der Stamm vom Partizip II \textit{getrunken} als Ableitungsbasis in Frage. Ein Substantiv wie \textit{Tritt} lässt sich auf die 3. Person Singular Präsens Indikativ \textit{tritt} beziehen. In diesem Sinne fand, wie \citet[132]{Donalies2005WbÜb} bemerkt, gar keine implizite Derivation statt, sondern eine bestehende Form wurde ohne Formsignal in eine andere Wortart überführt. Ob es sich um eine implizite Derivation oder aber um Konversion handelt, ist letztlich eine Frage der historischen Perspektive. Aus heutiger Sicht ist durch den Wechsel des Stammvokals das Kriterium der impliziten Derivation erfüllt. Ursprünglich kann der Vokalwechsel durch ein Suffix, also durch explizite Derivation motiviert worden sein. Bereits im Urgermanischen wurden umgelautete Stammformen ohne weiteres Formsignal in andere Wortarten überführt, also konvertiert.\footnote{Cf. \citet[78]{Paul1920Wb}.} Wichtig ist, dass wir die implizite Derivation genauso wie die Konversion, um die es in dem folgenden Abschnitt~\ref{subsection:Konversion} geht, zur Derivation zählen.\footnote{\citet[305]{Eisenberg2020_Wort} bestimmt \textit{Schuss} als implizite Derivation von \textit{schießen} und zählt diese genauso wie die Konversion zur Derivation. \citet[89--90]{FleischerBarz2012} zählen die implizite Derivation selbst zur Konversion.}


 \subsection{Konversion}\label{subsection:Konversion}
 
 Genauso wie die implizite Derivation ist die\is{Konversion} Konversion nicht kombinatorisch. Sie dient dem Wortartenwechsel, das heißt der Umsetzung in andere syntaktische Umgebungen mit der damit verbundenen Integration in andere Flexionsparadigmen. Anders als bei der impliziten Derivation wird der Wechsel der Wortart formal nicht ausgezeichnet. Es handelt sich somit um die einfachste Art der Wortbildung, um aus einfachen und präfigierten Verbstämmen (\ref{ea:V_S}) und aus Adjektiven (\ref{ex:A_S}) Substantive, aus Substantiven (\ref{ex:S_V}) und aus Adjektiven (\ref{ex:A_V}) Verbstämme sowie aus infiniten Verbformen (\ref{ex:Inf_S}) und aus flektierten Adjektiven (\ref{ex:AFx_S}) Substantive abzuleiten. 

\ea \label{ea:Konversionen}
\ea \label{ea:V_S} schlaf(en)/(der) Schlaf, besuch(en)/(der) Besuch
\ex \label{ex:A_S} grün/(das) Grün
\ex \label{ex:S_V} Fisch/fisch(en)
\ex \label{ex:A_V} grün/grün(en)
\ex \label{ex:Inf_S} fischen/das Fischen
\ex \label{ex:AFx_S} (der) grüne (Ball)/(der) Grüne
\z 
\z 


Mit \citet[31]{Erben2006_1975} unterscheiden wir zwischen der Stammableitung (\ref{ea:V_S})--(\ref{ex:A_V}) als morphologischer\is{Konversion!morphologisch} Konversion und der Umsetzung flektierter Formen (\ref{ex:Inf_S})--(\ref{ex:AFx_S}) als syntaktischer\is{Konversion!syntaktisch} Konversion.\footnote{Es sei daran erinnert, dass die Infinitivendung \textit{-en} eine Flexionsform und kein verbales Wortbildungssuffix ist. Wäre es das, sollte es wie \textit{-ier} in \textit{studier}(\textit{en}) in allen Flexionsformen stabil die Wortart kennzeichnen.} Im Folgenden geht es zuerst um die morphologische und anschließend um die syntaktische\is{Konversion} Konversion.

\begin{Definition}{Konversion}
    
\noindent Konversion ist eine besondere Art der Derivation, bei der ein lexikalisches Morphem ohne Affix oder Stammänderung in eine andere Wortart überführt wird. Werden flektierte Formen eines Morphems in eine andere Wortart überführt, so handelt es sich um syntaktische Konversion.
\end{Definition}



\subsubsection{Morphologische Konversion}\label{subsubsection:MorphKonv}

Uneingeschränkt möglich ist die\is{Konversion!morphologisch} Konversion eines Partizip I wie \textit{trinkend} zum\is{Partizip I!als Adjektiv} attributiv gebrauchten Adjektiv wie in (\ref{ex:trinkend}). Das Partizip I drückt aus, dass das vom Verb bezeichnete Geschehen andauert. Als Adjektiv überträgt es die Vorgangsbedeutung als Eigenschaft auf das Subjekt des ursprünglichen Verbs (\ref{ea:Verbszene_trinken}). Die übrigen verbalen Ergänzungen und Angaben können vom derivierten Adjektiv wie in (\ref{ex:Erg_trinkend}) kopiert werden.

\ea \label{ea:Partizip1_Adjektiv}
\ea \label{ea:Verbszene_trinken} Das Kind trinkt nach der Schule einen Kakao.
\ex \label{ex:trinkend} das trinkende Kind
\ex \label{ex:Erg_trinkend} das nach der Schule einen Kakao trinkende Kind
\z 
\z 

Anders als typische originäre Adjektive sind die adjektivischen Partizipien nur dann prädikativ verwendbar, wenn sie lexikalisiert sind (\ref{ea:Part1lexprädikativ}). Ansonsten ist das nicht möglich (\ref{ex:Part1prädikativ}). Sie sind nicht steigerbar (\ref{ex:SteigerungPart1}) und nicht mit \textit{un}- präfigierbar (\ref{ex:un_Part1}).

\ea \label{ea_Partizip1_untypAdjektiv}
\ea [] {das Kind ist wütend}\label{ea:Part1lexprädikativ}
\ex [*] {das Kind ist trinkend}\label{ex:Part1prädikativ}
\ex [*] {das trinkendere Kind}\label{ex:SteigerungPart1}
\ex [*] {das untrinkende Kind}\label{ex:un_Part1}
\z 
\z  

Das Partizip II von transitiven telischen Verben ist durchgängig zum\is{Partizip II!als Adjektiv} Adjektiv konvertierbar (\ref{ea:Part2_trinken}), wenn es attributiv auf ein ursprüngliches Akkusativobjekt (\ref{ea:Verbszene_trinken}) bezogen wird. Das ursprüngliche Subjekt kann, analog zum Passiv, als Vorgangsverursacher mit \textit{durch} oder \textit{von} ausgedrückt werden (\ref{ea:gebtrunkener Kakao}). Andere Angaben oder Ergänzungen können von der Prädikats- auf die Attributebene kopiert werden (\ref{ex:geschenkte}). 

\ea \label{ea:Part2_trinken}
\ea \label{ea:gebtrunkener Kakao} der (vom Kind) getrunkene Kakao
\ex \label{ex:geschenkte} der seinem Bruder zum Geburtstag geschenkte Ball
\z 
\z 

Wie anhand von (\ref{ea:gebtrunkener Kakao}) deutlich wird, beruht der Bezug auf das ursprüngliche Akkusativobjekt auf der passivischen Vorgangslesart. Denn \textit{getrunkener Kakao} ist Kakao, der getrunken wurde. Das Partizip II wird adjektivisch auf den Träger eines telischen Vorgangs bezogen. Es drückt den Nachzustand aus, den das Verb impliziert und überträgt ihn als Eigenschaft. 

Ganz\is{Partizip II!als Adjektiv} analog kann das Partizip II von intransitiven telischen Verben zum Adjektiv konvertiert werden, wenn es attributiv auf den ursprünglichen Vorgangsträger bezogen wird (\ref{ea:gewachsenes Kind}). Dieser ist allerdings das ursprüngliche Subjekt (\ref{ex:Kind_ist_gewachsen}). 

\ea \label{ea:Part2_gewachsen_geschlafen}
\ea \label{ea:gewachsenes Kind} das gewachsene Kind
\ex \label{ex:Kind_ist_gewachsen} Das Kind ist gewachsen.
\z 
\z 

Das Partizip II von intransitiven atelischen Verben kann nicht zum Adjektiv konvertiert werden (\ref{ea:*geschlafene}), da es keinen Nachzustand bedeutet. Der Gebrauch ist auf komplexe Prädikate wie (\textit{geschlafen haben}) beschränkt.

\ea \label{ea:*geschlafene} *das geschlafene Kind
\z 

Produktiv ist auch\is{Konversion!Produktivität} die morphologische Konversion von entlehnten Substantiven zu Verben: \textit{shopp}(\textit{en}), \textit{flirt}(\textit{en}), \textit{mail}(\textit{en}), \textit{facebook}(\textit{en}) usw. Das funktioniert(e) natürlich genauso bei heimischen Substantiven von \textit{Flöte} zu \textit{flöt}(\textit{en}), von \textit{Öl} zu \textit{öl}(\textit{en}), von \textit{Haus} zu \textit{haus}(\textit{en}) oder von \textit{Dampf} zu \textit{dampf}(\textit{en}). Mitunter gibt es parallel zu den Konvertaten implizite Derivata zum Ausdruck kausativer Pendants: \textit{dampf}(\textit{en}) und \textit{dämpf}(\textit{en}), \textit{staub}(\textit{en}) und \textit{stäub}(\textit{en}) sowie \textit{tafel}(\textit{n}) und \textit{täfel}(\textit{n}). Durch Konversion gebildete Verben gehören der schwachen Flexionsklasse an.

\citet[107--114]{Plank1981} zeigt anhand der Verbkonversion von Tierbezeichnungen, wie spezifisch die Restriktionen innerhalb einzelner Lesarten sein können. Verben wie \textit{bock}(\textit{en}), \textit{hamster}(\textit{n}) oder \textit{tiger}(\textit{n}) setzen mit der Paraphrase \gq{so wie das Tier handeln} auffällige und kulturell interpretierte Verhaltensmuster der Tiere voraus. Das Geburtsschema, dem \textit{kalb}(\textit{en}), \textit{ferkel}(\textit{n}) und \textit{lamm}(\textit{en}) zugrunde liegen, setzt ein Substantiv für die Bezeichnung des Nachwuchses voraus. Bildungen wie \textit{laus}(\textit{en}) oder \textit{wanz}(\textit{en}) mit privativer Bedeutung (Fehlen oder Entfernung der Basis von lat. \textit{privare} \gq{berauben}) setzen die Klassifizierung eines Tiers als Schädling voraus und stehen in Konkurrenz zu dem Privativmuster mit dem Präfix \textit{ent-} wie in \textit{entlaus}(\textit{en}) und \textit{entwurm}(\textit{en}). Das Beuteschema, welches Konvertaten wie \textit{fisch}(\textit{en}) und \textit{maus}(\textit{en}) zugrunde liegt, ist nicht mehr produktiv.

Die durch Konversion entstandenen Verben sind semantisch genauso unterspezifiziert wie Komposita. Das ursprüngliche Substantiv muss eine sinnvolle Interpretation im Rahmen eines Geschehens erlauben. Vorgeprägt sind die Lesarten durch entsprechende Rollen innerhalb des Verbszenarios. Hierbei geht es um die Subjektrolle wie bei \textit{spitzel}(\textit{n}) und um die Objektrolle wie bei \textit{bündel}(\textit{n}). Bei \textit{spachtel}(\textit{n}) und \textit{pinsel}(\textit{n}) liegt eine instrumentale Interpretation vor, bei \textit{zelt}(\textit{en}) eine lokale.

Bei der Konversion von Verben zu Substantiven sind nur noch Bildungen von starken Maskulina mit \textit{s}-Plural wie \textit{Treff}, \textit{Dreh}, \textit{Schwenk} oder umgangssprachlich \textit{Knutsch} oder \textit{Drück} produktiv.\footnote{Cf. \citet[128]{Reis1983} und \citet[137--138]{Vogel1996}.} Diese Konvertate drücken bis auf das lexikalisierte \textit{Treff} als Ort eine Sequenz des Verbalgeschehens aus.\footnote{Wird die Sequenz als Resultat verstanden, dann dient das Suffix \textit{-er} der Ableitung: \textit{Treffer}, \textit{Dreher}, \textit{Hüpfer} usw.} Ältere Bildungen wie \textit{Dank}, \textit{Raub} oder \textit{Schlag} fungieren als Nomina actionis und/oder sind auf eine Rolle des zugrunde liegenden Verbs festgelegt, das heißt lexikalisiert: \textit{Besuch} \gq{wer jemanden besucht}, \textit{Aufschnitt} \gq{was aufgeschnitten wird}, \textit{Sitz} \gq{worauf gesessen wird}.\footnote{Häufig besteht Mehrdeutigkeit, die erst im Gebrauchskontext auf eine Lesart festgelegt wird. So wird unter \textit{Bau} zusammen mit einem Verb wie \textit{dauern} das Ereignis als Nomen actionis ausgedrückt, mit einem Positionsverb wie \textit{stehen} das gebaute Objekt als Nomen acti.}

Die\is{Konversion!morphologisch} morphologische Konversion von Adjektiven zu Substantiven ist nach \citet[398]{Wilmanns1896Wb} nur bei Farb- und Sprachbezeichnungen produktiv. Aus \textit{grün} leiten wir das Substantiv \textit{Grün} ab und aus \textit{russisch} das Substantiv \textit{Russisch} wie in \textit{mein Russisch}. Konvertate wie *\textit{Schön} oder *\textit{Trocken} gibt es nicht. Dies übernimmt die syntaktische Konversion wie bei \textit{das Schöne} oder die explizite Suffixderivation wie bei \textit{Schönheit}.

Die Konversion von Substantiven zu Adjektiven ist äußerst eingeschränkt. Das Adjektiv \textit{orange} kann formal als Konvertat des entlehnten Substantivs \textit{Orange} analysiert werden. Tatsächlich wurde \textit{orange} aber im 17. Jahrhundert als Adjektiv aus dem Französischen (wo es ein Konvertat ist) entlehnt. Es bleiben Konvertate wie \textit{ernst} aus \textit{Ernst} oder \textit{schmuck} aus \textit{Schmuck}. Diese Adjektive sind jedoch wie \textit{schrott} und \textit{klasse} (noch) häufig auf den prädikativen Gebrauch mit einem Kopulaverb eingeschränkt. 

Die Konversion von Verben zu Adjektiven war und ist nicht produktiv. Deshalb weist \citet[414]{Wilmanns1896Wb} auf die Singularität von \textit{wirr} aus \textit{wirren} und \textit{wach} aus \textit{wachen} hin.

Morphologisch einfache Adjektive wurden zu Verben wie \textit{grün}(\textit{en}), \textit{faul}(\textit{en}), \textit{krank}(\textit{en}), \textit{lahm}(\textit{en}) oder \textit{reif}(\textit{en}) konvertiert. Diese Verben drücken Vorgänge oder ihren Beginn und Zustände aus. Selten sind unter den Konvertaten kausative Verben wie \textit{trockn}(\textit{en}), \textit{leer}(\textit{en}) oder \textit{locker}(\textit{n}). Kausativa wurden, wie im Abschnitt~\ref{subsection:ImplDerivation} dargestellt, durch implizite Derivata wie \textit{lähm}(\textit{en}), \textit{kränk}(\textit{en}) oder durch Präfigierungen wie \textit{ver-dünn}(\textit{en}), \textit{be-grün}(\textit{en}) und \textit{er-höh}(\textit{en}) gebildet. Letzteres ist der heute produktive Weg der Verbbildung aus Adjektiven.

Morphologisch komplexe Adjektive, die durch Suffigierung gebildet wurden, lassen sich nur mittels Präfix zu Verben wie \textit{be-ruhig}(\textit{en}), \textit{ver-deutlich}(\textit{en}) ableiten. \citet[83]{Echenlohr1999} hält die Konversion von Verben aus Adjektiven heute für nicht produktiv. Wir halten eine Wiederbelebung für möglich.\footnote{Auch \citet[201]{Olsen1990Head} hält das Muster für produktiv. Es mangele nur an potentiellen morphologisch einfachen Adjektivbasen.} Derzeit überwiegt aber die Verbalisierung einfacher Adjektive durch Präfigierung wie bei \textit{ver-doof}(\textit{en}) und durch Verbpartikeln wie bei \textit{auf-cool}(\textit{en}). Bildungen wie ?\textit{cool}(\textit{en}), ?\textit{hipp}(\textit{en}) oder ?\textit{doof}(\textit{en}) sind nicht belegt.

Alle bisher besprochenen\is{Konversion!Restriktionen} Konversionen haben eine Gemeinsamkeit. Morphologisch konvertierbar sind nur nicht-suffigierte Stämme. Offenbar kann durch Konversion die explizite Wortartfestlegung durch ein Suffix nicht überschrieben werden. Konvertate können nicht weiter konvertiert werden. Das sind wesentliche qualitative Unterschiede zur expliziten Derivation, insbesondere zur Suffigierung. 

Ausnahmen von der Regel, dass nur nicht suffigierte Stämme konvertierbar sind, scheinen Verben wie \textit{schneider}(\textit{n}) und \textit{fußballer}(\textit{n}) zu sein. Der Schein trügt jedoch, wenn wir eine Bildung wie \textit{doktern} betrachten. Denn hier handelt es sich nicht um Konversion von \textit{Doktor}, sondern um die Suffigierung von \textit{dokt-} mit \textit{-er}(\textit{n}) analog zu den bereits angesprochenen iterativen Bildungen wie \textit{bibber}(\textit{n}) oder \textit{knabber}(\textit{n}). Insofern ist auch bei \textit{fußballer}(\textit{n}) und \textit{schneider}(\textit{n}) davon auszugehen, dass das substantivbildende Suffix \textit{-er} im Zuge der Konversion verbal umgedeutet wird.\footnote{\citet[174]{Marchand1963} schreibt von einer \gqq{zweite[n] Motivierung durch den Ausgang -ern}.} Dies wird durch die Semantik begünstigt. Denn das Suffix \textit{-er} dient primär der Bildung von Nomina agentis, also von Handlungs- und Tätigkeitsträgern aus Verben. Gibt es für die entsprechende Tätigkeit kein Verb, so wird dieses aus dem Nomen agentis gebildet. Aus \textit{Gauner} wird \textit{gauner}(\textit{n}), aus \textit{Kellner} wird \textit{kellner}(\textit{n}) und aus \textit{Gärtner} wird \textit{gärnter}(\textit{n}). Im Detail wird die Analyse von \textit{fußballert} im Abschnitt~\ref{section:BspAnalysenWB} vorgestellt.

In der Literatur werden verschiedene\is{Wortstamm!Wortart} Analysevorschläge für Konversionen gemacht. \citet[347]{Bergenholtz_Mugdan1979} schlagen vor, dass einfache Stämme vom System her überhaupt nicht bezüglich einer Wortart bestimmt sind, sondern \gqq{grundsätzlich als Substantiv-, Adjektiv- oder Verbstamm dienen} können. Damit erübrigt es sich, von Konversion zu sprechen. Folglich sollten die meisten Stämme wie \textit{grün} als Adjektiv wie in \textit{die grüne Wiese}, als Substantiv wie in \textit{sattes Grün} und als Verb wie in \textit{die Heide grünt} gebraucht werden können. Das ist aber gerade nicht der Fall. Wie gezeigt wurde, sind die Möglichkeiten der Konversion stark eingeschränkt: *(\textit{das}) \textit{Schön}, ?\textit{doofen}, *\textit{ein hauser Mensch}. Die Idee unklassifizierter und in alle Richtungen der Inhaltswörter nutzbarer Stämme ist unplausibel. Außerdem wurde im Bereich der Suffigierung im Abschnitt~\ref{subsection:Suffigierung} und im Bereich der Präfigierung im Abschnitt~\ref{subsection:Präfigierung} gezeigt, dass sich Suffixe und Präfixe häufig nur mit Stämmen ganz bestimmter Wortarten verbinden, weshalb heute Bildungen wie *\textit{Baumung} oder *\textit{baumbar} nicht möglich sind. Hierin sieht \citet[188]{Olsen1990Head} ein starkes Argument gegen die Annahme unklassifizierter Stämme.   

Als Alternative wird seit \citet[293--296]{Marchand1969Nullaffix} versucht, die Kombinatorik und Rechtsköpfigkeit der Suffigierung durch Annahme eines formal\is{Nullmorphem} leeren Affixes, eines sogenannten Nullmorphems (\O) auch für die Konversion geltend zu machen. Eine differenzierte Analyse mit zwei verschiedenen Nullmorphemen hat \citet[202--213]{Olsen1990Head} für das Deutsche ausgearbeitet.\footnote{Eine knappe Übersicht und Kritik bietet \citet[57--63]{Echenlohr1999}.} Problematisch ist, dass somit die beschriebenen qualitativen Unterschiede zwischen Suffigierung und Konversion weggewischt würden. Außerdem ist problematisch, dass ein explizit (von dem französischen Verb \textit{camper}) abgeleitetes Verb wie \textit{kampier}(\textit{en}) die gleiche morphologische Struktur hätte wie das (von dem englischen Substantiv \textit{camping}) abgeleitete Verb \textit{camp}(\textit{en}). Die binär-rechtsköpfige Analyse von \textit{kamp-ier}(\textit{en}) wird auf \textit{camp-}\O(\textit{en}) übertragen. Plausibel wäre die Annahme eines unsichtbaren Null-Elements nur, wenn es zwischen beiden Verben eine Ableitungsbeziehung gäbe, was gerade nicht der Fall ist. Es bleibt mit \citet[48]{Eichinger2000} pointiert zu fragen, wie bezüglich eines Nullmorphems der Unterschied zwischen Null und Nichts zu bestimmen sei.

Wir gehen\is{Wortstamm!Wortart} im Rahmen des wortbasierten Ansatzes davon aus, dass Stämme mehrfach kategorisiert sein können.\footnote{Ebenso \citet[17]{Motsch2004}, \citet[73--75]{Echenlohr1999} \ua.} Das Phänomen kennen wir aus der Flexionsmorphologie. Die Form \textit{schläft} drückt die Kategorien 3. Person Singular Indikativ Präsens oder aber 2. Person Plural Indikativ Präsens aus. Ein Stamm wie \textit{schlaf-} ist analog als Verbstamm \textit{und} als Substantivstamm ausgezeichnet. Als Verbstamm verfügt \textit{schlaf} über die von diesem abweichende zweite Stammform \textit{schlief}. Der Substantivstamm ist defektiv, da keine Pluralbildung möglich ist. In konkreten Gebrauchskontexten wird die eine \textit{oder} die andere Stammform im jeweiligen Flexionsparadigma genutzt. Nur als Erstglied in Komposita kann kategorielle Ambiguität herrschen. Bei \textit{Schlafforschung} kann es sich analog zu \textit{Sprachforschung} um den substantivischen Stamm handeln, aber auch analog zu \textit{Sprechforschung} um den verbalen Stamm.

Ohne klare Indizien für eine gerichtete Konversion gehen wir von einer ungerichteten Beziehung zwischen mehrfach kategorisierten Stämmen aus. Teilweise ist selbst sprachhistorisch unklar, ob ein Verbstamm wie \textit{arbeit-} aus dem Substantivstamm \textit{Arbeit} konvertiert wurde oder umgekehrt.  

Wenn es jedoch klare Indizien für eine Konversionsrichtung gibt, so berücksichtigen wir diese in der morphologischen Analyse. So sind Substantive wie \textit{Befehl}, \textit{Entscheid} oder \textit{Verrat} deverbale Konvertate, weil sich die Präfixe \textit{be-}, \textit{ent-} und \textit{ver-} mit Verben, aber nicht mit Substantiven verbinden. Bei Verben wie \textit{pinsel}(\textit{n}) und \textit{spitzel}(\textit{n}) ist vom Primat der Substantive \textit{Pinsel} und \textit{Spitzel} auszugehen. Denn beim Primat der Verbstämme wären diese eher durch die \textit{er}-Ableitung gebildet worden: \textit{Pinsler} analog zu \textit{Bohrer} und \textit{Spitzler} analog zu \textit{Sammler}. In anderen Fällen ist die Produktivität ein Indiz für die Ableitungsrichtung. So stammt das Verb \textit{flöt}(\textit{en}) von \textit{Flöte}, da das Konversionsmuster vom Substantiv zum Verb teilweise produktiv ist. Manchmal lässt sich semantisch argumentieren. Ein Verb wie \textit{fisch}(\textit{en}) ist ein Konvertat des Substantivs \textit{Fisch}, da es dieses in seiner Bedeutungsstruktur \gq{Fische fangen} voraussetzt. Umgekehrt lässt sich eine solche Beziehung nicht herstellen. 


\subsubsection{Syntaktische Konversion}\label{subsubsection:SyntKonv}

Die\is{Konversion!syntaktisch} syntaktische Konversion bildet einen Übergangsbereich zwischen Syntax und Wortbildung. Durch sie werden syntaktische, das heißt flektierte Wortformen in eine andere Wortart umgesetzt. 
Diese Technik lässt sich gut ausgehend von dem Szenario (vgl. Definition~\ref{Def:Szenario} in Abschnitt~\ref{section:KonzeptVerbvalenz}), das eine Verbform wie \textit{liebt} ausdrückt, illustrieren. Dieses Szenario verlangt nach \textit{einem} oder \textit{einer Liebenden} sowie nach \textit{einem} oder \textit{einer Geliebten}. Das Ereignis, das beide Rollen konstituieren, benennen wir als (\textit{das}) \textit{Lieben}. \citet[22]{Eichinger2000} bezeichnet die syntaktische Konversion als eine \gqq{unaufwendige Technik [\ldots{}], um einen Namen aus verschiedenen Teilen einer verbalen Szene zu destillieren}.

Das adjektivisch\is{Partizip I!als Substantiv} flektierte Partizip I \textit{liebend-} (\ref{ea:Adj_liebend}) setzen wir ohne Derivationssignal als \textit{der Liebende} oder \textit{die Liebende}, sexusneutral im Plural als \textit{die Liebenden} in die Wortart Substantiv um. Die flektierte Form des Adjektivs aus der Nominalphrase (\ref{ea:Adj_liebend}) wird zum substantivischen Kopf der Nominalphrase (\ref{ex:Subst_liebender}).

\ea \label{ea:PIalsSubstantiv}
\ea \label{ea:Adj_liebend} der liebende Bruder/ein liebender Bruder, die liebende Schwester/eine liebende Schwester
\ex \label{ex:Subst_liebender} der Liebende/ein Liebender, die Liebende/eine Liebende
\z 
\z 

Diese Struktur ist elliptisch motiviert. Durch die Auslassung (Ellipse) des substantivischen Kopfs (\ref{ea:Adj_liebend}) rutscht das flektierte Adjektiv in diese Funktion (\ref{ex:Subst_liebender}). Formal wird diese syntaktische Konversion durch Großschreibung gekennzeichnet. Von der Kategorie des Adjektivs erbt das substantivische Konvertat sowohl die ursprüngliche Adjektivflexion wie bei \textit{der Liebend-e} als auch die eigentliche Substantivflexion wie bei \textit{ein Liebend-er}, welche normalerweise am Artikel manifest wird. Ebenso ererbt ist die für Adjektive typische Genusvariabilität.

Die Substantivierung\is{Partizip II!als Substantiv} eines Partizip II wie \textit{geliebt} verläuft über den flektierten adjektivischen Gebrauch wie in (\ref{ea:geliebter_Adj}). Das flektierte Adjektivkonvertat kann genauso wie das flektierte adjektivische Partizip I zum Kopf der Nominalphrase und somit ad hoc zum Substantiv werden (\ref{ex:geliebt_Subst}).

\ea \label{ea:PartizipII_Verb_Adj_Subst}
\ea \label{ea:geliebter_Adj} der geliebte Bruder
\ex \label{ex:geliebt_Subst} der Geliebte/ein Geliebter
\z 
\z 

Der adjektivische Gebrauch bildet die Brücke zum substantivischen Gebrauch. Dieser ist systematisch bei allen flektierten Adjektiven zur Bezeichnung des Eigenschaftsträgers angelegt: \textit{der Große}/\textit{ein Großer}, \textit{das Große}/\textit{ein Großes}, \textit{die Große}/\textit{eine Große}. Die bestehende Flexionsform aus der Nominalphrase übernimmt bei Auslassung des ursprünglichen Kopfes direkt die Kopffunktion.

Verbale Infinitive\is{Infinitiv!als Substantiv} wie \textit{tun} können zu substantivischen Infinitiven wie (\textit{das}) \textit{Tun} konvertiert werden. Hierdurch wird von der prädikativen und satzbildenden Funktion (wer tut was) abstrahiert. Substantivierte Infinitive sind Neutra und nur dann pluralfähig, wenn sie lexikalisiert sind und auf Konkreta verweisen: \textit{zwei Schreiben} versus *\textit{zwei Tun}. \citet[307]{Eisenberg2020_Wort} bemerkt, dass der syntaktische Konversionszyklus \gqq{dem Prinzip \textit{Endstation Hauptwort} unterworfen} ist. Der eine Weg führt vom Infinitiv zum Substantiv, der andere Weg führt vom adjektivischen Partizip zum Substantiv.  

Was wird durch die syntaktische Konversion gewonnen? Die deverbale Konversion erlaubt es, die prädikativen Leistungen von Verben kompakt in Wortformen zu gießen und sie als solche im Text (wieder) zu benennen. Das kann als Gegenstand der Rede in einer neuen Prädikation erfolgen wie in (\ref{ea:Schlafender_SG}) und (\ref{ex:Schlaf_SG}) oder der Modifikation eines Redegegenstandes als Attribut wie in (\ref{ex:Schlafender_Attr}) dienen.  

\ea \label{ea:Konverate_Redegegenstände}
\ea \label{ea:Schlafender_SG} Am anderen Morgen ging er leise vor die Rohrhütte, in der sein Gast schlief, setzte sich auf die dunklen Steine des alten Feuerplatzes und sah den Schlafenden an.\footnote{Wiechert, Ernst: \textit{Das einfache Leben}. München: Ullstein, 2000 [1946], S. 384.}
\ex \label{ex:Schlaf_SG} Sein verbrauchter erschöpfter Körper wünschte sich Alleinsein, Schlafen.\footnote{Seghers, Anna: \textit{Die Wellblech-Hütte}. In: Erzählungen 1926--1944. Berlin: Aufbau-Verlag, 1981, S. 141.}
\ex \label{ex:Schlafender_Attr} Erlöst, ohne Scham lag sie neben ihm, während er eingeschlafen war. Das hübsche, verderbte Gesicht des Schlafenden sah kindlich aus [\ldots{}]\footnote{Feuchtwanger, Lion: \textit{Erfolg}. In: Gesammelte Werke in Einzelbänden. Bd. 6. Berlin: Aufbau-Verlag, 1993 [1930], S. 433. }
\z 
\z 

Das Substantivkonvertat \textit{Schlafender} in (\ref{ea:Schlafender_SG}) erfüllt die Rolle desjenigen, der angesehen wird. Dass er dabei schläft, wird nicht auf Satzebene prädiziert, sondern quasi lexikalisch vorausgesetzt. In (\ref{ex:Schlaf_SG}) fungiert das (von einem zu einer bestimmten Zeit Schlafenden abstrahierte) \textit{Schlafen}-Ereignis als eine gewünschte Entität. In (\ref{ex:Schlafender_Attr}) dient das Konvertat \textit{Schlafender} der näheren Beschreibung von \textit{Gesicht}.     

Die syntaktische\is{Konversion!Restriktion} Konversion unterliegt anders als die morphologische und anders als die eigentliche Wortbildung kaum Restriktionen. Infinitive und das Partizip I können grundsätzlich syntaktisch zum Substantiv konvertiert werden. Die syntaktische Konversion vom Partizip II zum Substantiv ist an die vorhergehende morphologische Konversion zum Adjektiv gebunden. Und diese ist klar restringiert. Sie funktioniert nur, wenn das Partizip II den Nachzustand eines Geschehens ausdrückt, der sinnvoll auf einen Eigenschaftsträger bezogen werden kann. Bei telischen intransitiven und transitiven Verben ist das problemlos möglich. Deshalb gibt es \textit{ein eingeschlafenes Kind}, was zu \textit{das Eingeschlafene} konvertiert werden kann. Bei \textit{geschlafen} ist das nur innerhalb einer besonderen Konstruktion wie \textit{der geschlafene Schlaf der Gerechten} möglich. Die syntaktische Konversion zu \textit{der Geschlafene} ist zumindest sehr auffällig. 

Eine besondere Art der\is{Phrasenkonvertat} syntaktischen Konversion liegt vor, wenn ganze Wortgruppen mit den internen Flexionsformen zu einem Substantiv konvertiert werden. Substantive wie in (\ref{ea:Phrasenkonvertat}) werden deshalb als Phrasenkonvertate bezeichnet.  

\ea \label{ea:Phrasenkonvertat}
\ea \label{ea:Taugenichts} der Taugenichts, der Nimmersatt
\ex \label{ex:Vergissmeinnicht} das Vergissmeinnicht, das Vaterunser
\ex \label{ex:Zitatkonv} das Ich-hab-keine-Lust-mehr
\z 
\z 

Die Letztglieder sind nicht die Köpfe der Phrasenkonvertate. Die Genusfestlegung geschieht semantisch. Bezeichnungen für Personen wie in (\ref{ea:Taugenichts}) sind maskulin. Bezeichnungen für Nicht-Personen wie in (\ref{ex:Vergissmeinnicht}) sind Neutra. Uneingeschränkt möglich sind entsprechende Konversionen von Phrasen zu Substantiven als zitierte Anführungen wie in (\ref{ex:Zitatkonv}).

Phrasenkonvertate werden, wie in Exkurs~\ref{subsection:Übergänge} dargestellt, auch als Erstglieder in Komposita gebraucht \zB bei \textit{das Ich-hab-keine-Lust-mehr-Gejammer}. Dies gilt auch für (noch) nicht wortfähige Phrasenkonvertate wie in \textit{Gute-Kita-Gesetz} oder \textit{Rotes-Kreuz-Hilfe}. 
 

Um eine\is{Inkorporation} Verbindung zwischen Phrasenkonversion und Komposition handelt es sich, wenn ein Infinitiv mit seiner Objekt-NP (\ref{ea:Haare_schneiden}) wie in (\ref{ex:Haareschneiden}) substantiviert wird.

\ea \label{ea:Phrasenkonv_Haareschneiden}
\ea \label{ea:Haare_schneiden} Haare schneiden
\ex \label{ex:Haareschneiden} das Haareschneiden
\z 
\z 

Die Pluralform des Erstgliedes \textit{Haare} in (\ref{ex:Haareschneiden}) ist syntaktisch durch (\ref{ea:Haare_schneiden}) motiviert. Das Konvertat entspricht gleichzeitig dem Muster der Determinativkomposita (vgl. Abschnitt~\ref{subsection:DetKomp}). Das Zweitglied \textit{Schneiden} ist der formale und semantische Kopf des Konvertats. Deshalb ist \textit{Haareschneiden} ein nicht pluralfähiges Neutrum. Nicht nur Objekt-NPs werden auf diese Weise inkorporiert, sondern auch PPs wie in \textit{das Auf-den-Bus-Warten}. Voraussetzung hierfür ist, dass die ursprünglichen Objekt-NPs indefinit wie \textit{Haare} in \textit{Haare schneiden} oder schwach definit wie \textit{auf den Bus} in \textit{auf den Bus warten} bereits ihre Referenz eingebüßt haben und Teile eines Tätigkeitsausdrucks (vgl. Abschnitt~\ref{subsection:IdiomeNVGFVG}) geworden sind.


Verbale Infinitive wie \textit{schneiden} können restriktionslos zu (\textit{das}) \textit{Schneiden} substantiviert werden, aber es handelt sich nur um Ad-hoc-Wörter. Sie können nicht wie etablierte Wörter in der Wortbildung gebraucht werden. Als rechts stehende Köpfe von Komposita betten sie keine Kompositionsstammformen als Erstglieder ein, sondern inkorporieren die ursprünglichen Objekt-NPs. Der Unterschied zeigt sich deutlich beim Vergleich von \textit{Haareschneiden} mit \textit{Haarschnitt}, von \textit{Händewaschen} mit \textit{Handwaschmittel} oder von \textit{Äpfel-Ernten} mit \textit{Apfelernte}. 

Nur wenn der substantivierte Infinitiv lexikalisiert ist wie \textit{Essen} im Sinne von \gq{Gericht}, kann ein ganz gewöhnliches Kompositum wie \textit{Kloßessen} (statt \textit{Klöße-Essen}) gebildet werden. Dieser Unterschied zeigt sich auch, wenn substantivierte Infinitive als Erstglied gebraucht werden. Nur die lexikalisierten Bildungen wie \textit{Leiden}, \textit{Leben} oder \textit{Essen} bilden substantivische Kompositionsstammformen mit der \textit{s}-Fuge wie in \textit{Leidensdruck}, \textit{Lebensglück} und \textit{Essensrest}. Nicht lexikalisierte Ad-hoc-Substantive wie \textit{Schneiden} oder \textit{Waschen} sind nicht erstgliedfähig. Hier werden verbale Kompositionsstammformen wie in \textit{Schneidebrett} oder \textit{Waschmittel} gebildet.     


\section{Bildung von Partikelverben}\label{section:Vpt Bildung}

Unter\is{Partikelverbbildung} semantischer und funktionaler Perspektive hätte die Bildung von Partikelverben wie \textit{einschlafen} der Bildung von Präfixverben wie \textit{verschlafen} nachgeschaltet werden müssen. Strukturell aber gibt es wesentliche Unterschiede, die die Abkopplung der Partikelverbbildung in diesem Abschnitt von der Präfigierung in Abschnitt~\ref{subsection:Präfigierung} motivieren.

Zu den Verbpartikeln (vgl. Abschnitt~\ref{subsection:VPartikeln}) gibt es bis auf \textit{ein-} wie in \textit{ein-kaufen} gleichlautende\is{Verbpartikel!versus Präposition} Präpositionen, \zB die Präposition \textit{an} und die Verbpartikel \textit{an-} in \textit{an-schließen}. Deshalb hat \citet[870]{Grimm1826Gr2} Partikelverben als \textit{Partikelkomposita} klassifiziert. Dagegen spricht jedoch, was auch Grimm sieht, dass sich Präpositionen anders verhalten als Verbpartikeln. Erstere regieren einen bestimmten Kasus (\ref{ea:schließen_an}), letztere tun das nicht (\ref{ex:anschließen}). 

\ea \label{ea:VPart_Präp}
\ea \label{ea:schließen_an} Sie schließt das Rad an den Baum.
\ex \label{ex:anschließen} Sie schließt das Rad an.
\z 
\z 

Außerdem wird in (\ref{ex:anschließen}) ersichtlich, dass Verbpartikeln anders als Erstglieder von Komposita syntaktisch getrennt stehen können. Partikelverben sind folglich nicht als Komposita mit präpositionalem Erstglied zu analysieren.

Semantisch lassen sich viele Verbpartikeln auf die entsprechenden Präpositionen im Sinne von räumlichen und davon abgeleiteten Beziehungen zurückführen. In (\ref{ex:anschließen}) richtet sich die Handlung des \textit{Anschließens} auf das Objekt \textit{Rad}, welches in eine gewisse Kontaktposition gebracht wird. Dies lässt sich wie in (\ref{ea:schließen_an}) präpositional paraphrasieren. Das Partikelverb \textit{anschließen} erlaubt, anders als \textit{schließen}, die Abstraktion von dem Kontaktort.

In (\ref{ea:ansteigen}) verläuft der vom Verb \textit{an-steigen} beschriebene Vorgang aus Sicht des Sprechers in einer bestimmten Richtung, was durch ein Adverb wie \textit{hinauf} in (\ref{ex:hinauf_steigen}) paraphrasierbar ist.\footnote{Diese beiden Grundrelationen beschreibt \citet[12--19]{Eichinger1989} detailliert.}

\ea \label{ea:an_hinauf_steigen}
\ea \label{ea:ansteigen} Die Preise steigen an.
\ex \label{ex:hinauf_steigen} Die Preise steigen hinauf.
\z 
\z 

Allerdings führt die Verbpartikel im Gegensatz zu dem entsprechenden Adverb häufig zu einer semantischen Spezialisierung. So bedeutet \textit{ausgehen} etwas anderes als \textit{herausgehen}, \textit{einwandern} etwas anderes als \textit{hineinwandern} und \textit{ankommen} etwas anderes als \textit{herankommen}.

Im Gegensatz zu Präfixen (\ref{ea:umfährt}) tragen Verbpartikeln (\ref{ex:fährt_um}) den Hauptakzent (\ref{ex:Betonung_Px_Pt}).

\ea \label{ea:um_umfahren}
\ea \label{ea:umfährt} Sie umfährt das Hindernis.
\ex \label{ex:fährt_um} Sie fährt das Hindernis um.
\ex \label{ex:Betonung_Px_Pt} UMfahren (\ref{ex:fährt_um}) versus umFAHRen (\ref{ea:umfährt})
\z 
\z 

Anders als\is{Verbpartikel!Stellung} Präfixe führen Verbpartikeln ein syntaktisches Eigenleben, indem sie in Verbzweit- (\ref{ea:Trennung1}) und Verberstsätzen (\ref{ex:Trennung2}) genauso wie andere Prädikatsteile (\ref{ex:Trennung3}) die rechte Satzklammer bilden (vgl. Abschnitt~\ref{section:Verbzweitsätze}). Bei Verbletztstellung (\ref{ex:Adjazenz1}) und beim Infinitiv (\ref{ex:Adjazenz2}) stehen Verbpartikel und Verb nachbarschaftlich (adjazent) und werden zusammengeschrieben.

\ea \label{ea:Stellung_VPt}
\ea \label{ea:Trennung1} Er wischt den Tisch ab.
\ex \label{ex:Trennung2} Wisch den Tisch ab!/Wischt er den Tisch ab?
\ex \label{ex:Trennung3} Er wischt den Tisch sauber.
\ex \label{ex:Adjazenz1} dass er den Tisch abwischt
\ex \label{ex:Adjazenz2} Abwischen will er den Tisch.
\z 
\z 

Während Präfixverben wie \textit{belügen} das Partizip II wie \textit{belogen} ohne \textit{ge-} bilden, steht es bei Partikelverben wie in \textit{an-ge-logen} zwischen Partikel und Verbbasis, also verbnäher. Das Gleiche gilt für die Infinitivpartikel \textit{zu}. Sie steht wie bei \textit{an-zu-lügen} zwischen Verbpartikel und Verb.

Insbesondere aus schulgrammatischer Sicht regen wir an, die Grundform von Partikelverben nicht nur durch Infinitive wie \textit{anlügen}, \textit{aufwachen} oder \textit{einladen} anzugeben, sondern durch eine flektierte klammerbildende Form wie \zB \textit{lügt \ldots{} an}, \textit{wacht \ldots{} auf} oder \textit{lädt \ldots{} ein}. Somit wird der Normalfall der syntaktischen prädikativen Klammerbildung unterstrichen. 

In dem folgenden Abschnitt~\ref{subsection:Vpt_Syntax_Wortbildung} gehen wir der Frage nach, ob es sich bei der Bildung von Partikelverben eher um ein syntaktisches oder morphologisches Phänomen handelt. Danach argumentieren wir in Abschnitt~\ref{subsection:Pt_Kopf} dafür, dass Partikeln analog zu den Präfixen Verben formen, aber nicht als Kopf bestimmt werden sollten. Abschließend gehen wir in Abschnitt~\ref{subsection:Vpt_Reihenbildung} auf den reihenbildenden Charakter von Partikeln ein.


\subsection{Syntax oder Wortbildung?}\label{subsection:Vpt_Syntax_Wortbildung}

Partikeln sind\is{Partikelverbbildung} auf eine andere Weise als Präfixe mit ihrer Basis verbunden, nämlich syntaktisch. Deshalb zählen \ua \citet{Luedeling2001} und \citet{Mueller2002} die produktive Partikelverbbildung zur Syntax und nicht zur Wortbildung. Denn resultative Partikeln (\ref{ea:austrinken}) verhalten sich analog zu Objektsprädikativen (\ref{ex:leer_trinken}) (vgl. Abschnitt \ref{subsection:SP und OP}), weshalb beide auch nicht miteinander kombiniert werden können (\ref{ex:leer_austrinken}).

\ea \label{ea:OP_VPart}
\ea [] {Sie trinkt die Flasche aus.}\label{ea:austrinken}
\ex [] {Sie trinkt die Flasche leer.}\label{ex:leer_trinken}
\ex [*] {Sie trinkt die Flasche leer aus.}\label{ex:leer_austrinken}
\z 
\z 

Der wesentliche Unterschied ist aber, dass als Objektsprädikativ gebrauchte Adjektive über eine eigene lexikalische Bedeutung verfügen, weshalb sie auch verbunabhängig wie in (\ref{ea:leere_Flasche}) gebraucht werden können. Das ist bei den meisten Partikeln nicht der Fall (\ref{ex:ause_Flasche}). Sie dienen wie Präfixe der Bildung komplexer Verben.

\ea \label{ea:OP_VPart2}
\ea [] {Die Flasche ist leer. -- die leere Flasche}\label{ea:leere_Flasche}
\ex [*] {Die Flasche ist aus -- *die ause Flasche}\label{ex:ause_Flasche}
\z
\z 

Entsprechende Gemeinsamkeiten und Unterschiede gibt es auch zwischen Adverbien bzw. adverbial gebrauchten Adjektiven und PPs (\ref{ea:hinauf_blicken}) auf der einen (syntaktischen) Seite und Partikeln (\ref{ex:auf_blicken}) auf der anderen Seite. Letztere sind im Gegensatz zu ihren syntaktischen Pendants (\ref{ex:Blick_nach_oben}) nicht außerhalb der Verbbildungsmuster wortfähig (\ref{ex:*Blick_auf}), d.\,h. sie können nicht als syntaktische Wörter gebraucht werden. 

\ea \label{ea:Adv_VPart}
\ea [] {Sie blickt hinauf/hoch/nach oben.}\label{ea:hinauf_blicken}
\ex [] {Sie blickt auf.}\label{ex:auf_blicken}
\ex [] {der Blick hinauf/hoch/nach oben}\label{ex:Blick_nach_oben}
\ex [*] {der Blick auf}\label{ex:*Blick_auf}
\z 
\z 

Allerdings können\is{Partikelverbbildung} sich die Partikeln bestimmter Gruppen durch Bedeutungsübertragung vom Partikelverb auf die Partikel zum Adjektiv entwickeln, das prädikativ und wenigstens umgangssprachlich auch attributiv gebraucht wird.

\ea \label{ea:Part_zu_Adj}
\ea \label{ea:an_aus_sein} Das Licht ist an/aus. -- das ane/ause Licht
\ex \label{ex:auf_zu_sein} Die Tür ist auf/zu. -- die aufe/zue Tür
\z 
\z 

Anders als bei typischen syntaktischen Strukturen wie \textit{leer trinken} und \textit{hinauf steigen} ergibt sich die Bedeutung von Partikelverben ähnlich wie die Bedeutung von Präfixverben häufig nicht aus der Partikel und dem Basisverb. Die Bedeutung von \textit{anfangen} und \textit{aufhören} (im Sinne von \gq{endigen}) liegt analog zu der von \textit{verstehen} nur in der Gesamtbildung. Sie ist nicht (mehr) die Summe der Bedeutungen ihrer Bestandteile bzw. die Summe aus Verbbedeutung und Wortbildungsoperation. Das Gleiche ist aber auch der Fall bei lexikalisierten Verbindungen aus Objektsprädikativ und Verb wie bei \textit{krankschreiben} oder \textit{freisprechen}.

Eine feste Basis bilden Verbpartikeln und Verbbasis erst bei der expliziten Derivation mittels Suffix (\ref{ea:Einladung}), bei der impliziten Derivation (\ref{ex:Abflug}), bei der morphologischen Konversion (\ref{ex:Einkauf}) und bei der syntaktischen Konversion (\ref{ex:Mitmachen}).

\ea \label{ea:VPt_Derivation}
\ea \label{ea:Einladung} einlad(en) $\rightarrow$ Einladung
\ex \label{ex:Abflug} abflieg(en) $\rightarrow$ Abflug
\ex \label{ex:Einkauf} einkauf(en) $\rightarrow$ Einkauf
\ex \label{ex:Mitmachen} mitmachen $\rightarrow$ Mitmachen
\z 
\z 

Verbpartikeln verbinden sich primär mit verbalen Stammformen, das heißt mit verbalen Flexionsstammformen. Denn auch Partikeln verändern nicht die Flexionsklasse eines Verbs. Starke Verben wie \textit{laden} (\textit{lud}) bleiben wie \textit{einladen} (\textit{lud \ldots{} ein}) stark und schwache Verben wie \textit{packen} (\textit{packte}) bleiben wie \textit{einpacken} (\textit{packte \ldots{} ein}) schwach. 

Von ureigentlicher Wortbildung kann eigentlich nur gesprochen werden, wenn Substantive wie \textit{Strang} und Adjektive wie \textit{blau} durch Verbindung mit einer Partikel verbalisiert werden wie in \textit{an-strängen} (\gq{ein Zugtier mit einem Strang vor einen Wagen spannen}) und \textit{ein-bläuen}. Dann handelt es sich bei der Basis aus heutiger Sicht um eine Derivationsstammform.\footnote{Häufig existierte oder existiert ein entsprechendes Simplexverb wie \textit{strängen} oder \textit{bläuen}. Hier handelt es sich um implizite Derivationen.} Dennoch werden wir in dem folgenden Abschnitt~\ref{subsection:Pt_Kopf} dafür argumentieren, Partikeln nicht als Köpfe zu bestimmen.


\subsection{Partikeln als Köpfe?}\label{subsection:Pt_Kopf}


Abgesehen von den wesentlichen Unterschieden zwischen Präfix- und Partikelverben, beeinflussen Partikeln ähnlich wie Präfixe die Valenz des Basisverbs und im Falle von Verbalisierung sind sie \gq{Valenzmacher}. Die\is{Partikelverbbildung!Valenzdynamik} Valenz eines Verbs kann reduziert werden, indem die Partikel ein ursprüngliches Direktivum sättigt bzw. inkorporiert (\ref{ea:aufhängen}). Genauso können analog zu Präfixen Valenzstellen eröffnet werden (\ref{ex:anlügen}). Die hinzugefügte Valenzstelle lässt sich teilweise direkt auf die Ergänzung der entsprechenden Präposition beziehen (\ref{ex:zuwerfen}). 

\ea \label{ea:VPt_Valenzmacher}
\ea \label{ea:aufhängen} Er hängt die Wäsche auf. (auf die Leine)
\ex \label{ex:anlügen} Er lügt sie an. (*Er lügt sie.)
\ex \label{ex:zuwerfen} Er wirft ihr den Ball zu. (zu ihr)
\z 
\z 

In Abschnitt~\ref{subsubsection:Präfix_Kopf} haben wir die Verbalisierung von Substantiven und Adjektiven durch Präfixe als einen Sekundäreffekt beschrieben. Auch Partikeln, die eigentlich der Modifikation von Verbbasen dienen, können wie in (\ref{ea:Vpt_Verbalisierung}) einen solchen verbalisierenden Sekundäreffekt\is{Partikelverbbildung!Verbalisierung} haben. 

\ea \label{ea:Vpt_Verbalisierung}
\ea \label{ea:ankreiden} Er feuchtet die Briefmarke an.
\ex \label{ex:auftischen} Er tischt ihr einen Braten auf.
\ex \label{ex:ausmisten} Wir misten den Stall aus.
\ex \label{ex:einschulen} Sie schulen die Kinder ein.
\z 
\z 

Ebenso können Partikeln, genauso wie Präfixe, bei gleichbleibender Valenz zur Telisierung von Verben führen (\ref{ea:Vpt_Perfektivierung}).

\ea \label{ea:Vpt_Perfektivierung}
\ea \label{ea:aufblühen} Am Wochenende baden wir (die Saison) an.
\ex \label{ex:aufblühen} Die Blume blüht auf.
\z 
\z 

Aufgrund der\is{Partikelverbbildung} zu Präfixen analogen valenziellen Kraft bis hin zur Verbalisierung zählen wir die Bildung von Partikelverben zur Wortbildung. Partikeln formen analog zu Präfixen Verben, sind aber keine morphologischen Köpfe. Der Bildung von Partikelverben räumen wir wie \citet[1032]{Weinrich2003} neben der Komposition und der Derivation einen ganz eigenen Status ein. Die Bildung von Partikelverben sehen wir als eine stabile Schnittstelle zwischen Syntax und Wortbildung. Die Nähe zur Wortbildung wird insbesondere dann sichtbar, wenn Substantive und Adjektive als Derivationsstammformen durch Partikeln verbalisiert werden. Ist die Basis verbal, so geht sie direkt als lexikalische Stammform in die Bildung ein.

Dass Partikeln ähnlich wie Präfixe Verben formen, zeigt sich besonders dann, wenn beide Bildungstypen kombiniert auftreten. In dem seltenen Fall,\footnote{Nach der Untersuchung von \citet[143]{Kuehnhold1973} machen solche Doppelbildungen nur 2\% der komplexen Verben aus.} dass ein Präfixverb wie \textit{erziehen} (\ref{ea:erziehen}) mit einer Partikel verbunden ist wie bei \textit{anerziehen} (\ref{ex:anerziehen}), hat der zuletzt stattgefundene Bildungsprozess, also die Partikelverbbildung, das valenzielle Sagen.

\ea \label{ea:an_erziehen}
\ea \label{ea:erziehen} Die Eltern erziehen ihr Kind.
\ex \label{ex:anerziehen} Die Eltern erziehen ihrem Kind Pünktlichkeit an.
\z 
\z 

Aus dem zweiwertigen Präfixverb \textit{erziehen} in (\ref{ea:erziehen}) ist in (\ref{ex:anerziehen}) das dreiwertige Partikelverb \textit{erziehen\ldots{} an} geworden.

Analog zu den Präfixen bilden auch Partikeln bzw. einzelne Untergruppen größere Gruppen semantisch und valenziell ähnlicher Verben, sogenannte Reihen, auf die wir in dem folgenden Abschnitt~\ref{subsection:Vpt_Reihenbildung} eingehen. 


\subsection{Reihenbildung}\label{subsection:Vpt_Reihenbildung}
\largerpage

Genauso wie bei Präfixverben bilden verschiedene Partikelverbmuster\is{Reihenbildung!Partikelverb} Reihen. Die Mitglieder einer Reihe drücken unabhängig von der Basis mit gleicher Valenz den gleichen Situationstyp aus, was anhand der Reihen in (\ref{ea:VPt_Reihen}) illustriert wird. Für die systematische Beschreibung einzelner Verbpartikeln und Bildungsmuster verweisen wir auf \citet[141--375]{Kuehnhold1973}. Eine kurze Übersicht bieten \citet[399--418]{FleischerBarz2012}.


\ea \label{ea:VPt_Reihen}
\ea \label{ea:ab_Reihe} jemand bricht/schließt/reißt \ldots{} etwas ab 
\ex \label{ex:an_Reihe} jemand lächelt/tanzt/spricht \ldots{} jemanden an
\ex \label{ex:aus_Reihe} jemand wandert/reist/bricht \ldots{} aus
\ex \label{ex:bei_Reihe} jemand fügt/mischt/rührt \ldots{} einer Sache etwas bei
\ex \label{ex:durch} etwas scheint/tropft/fließt \ldots{} durch
\ex \label{ex:ein_Reihe} jemand tütet/kerkert/schult \ldots{} etwas/jemanden ein
\ex \label{ex:gegen_Reihe} jemand liest/checkt/rechnet \ldots{} etwas gegen
\ex \label{ex:nach_Reihe} jemand läuft/strebt/ruft \ldots{} jemandem nach
\ex \label{ex:über_Reihe} etwas läuft/schäumt/quillt \ldots{} über
\ex \label{ex:um_Reihe} jemand füllt/färbt/deutet \ldots{} etwas um
\ex \label{ex:vor_Reihe} jemand sagt/spielt/tanzt \ldots{} jemandem etwas vor
\z 
\z 

Genauso wie bei Präfixverben sind einzelne Partikelverben oder ein gesamtes Bildungsmuster durch Kontrastformen mit anderen Partikeln und Präfixen motiviert, was in (\ref{ea:VPt_Kontrastformen}) exemplarisch dargestellt wird.

\ea \label{ea:VPt_Kontrastformen}
\ea \label{ea:an_aus} etwas anschalten/ausschalten
\ex \label{ex:an_be_ab} etwas ansprühen/besprühen/absprühen
\ex \label{ex:an_zu} jemanden ansehen/jemandem bei etwas zusehen
\ex \label{ex:an_ab} anbaden/abbaden
\ex \label{ex:an_be_durch} etwas ansprechen/besprechen/durchsprechen
\ex \label{ex:auf_zu} etwas aufschließen/zuschließen
\ex \label{ex:ein_ab} etwas einschalten/abschalten
\ex \label{ex:ein_aus_ver} etwas einpacken/auspacken/verpacken
\ex \label{ex:an_auf} etwas ankleben/aufkleben
\z 
\z 

Formal ähneln komplexe Verben mit substantivischem Erstglied wie \textit{eis-laufen}, \textit{brust-schwimmen} oder \textit{not-landen} den Partikelverben. Sie basieren aber nicht auf deren Bildungsmuster. Um solche Bildungen geht es im folgenden Abschnitt~\ref{section:andere Wb-Arten}. 


\section{Rückbildung, Kürzung, Dopplung oder Komposition und Derivation?}\label{section:andere Wb-Arten}

In diesem Abschnitt geht es um Bildungen, die auf den ersten Blick weder auf Komposition noch auf Derivation oder Partikelverbbildung beruhen. Wir zeigen, dass die bisher dargestellten Wortbildungsarten weitgehend ausreichen, um diese Bildungen systematisch zu analysieren. In dem folgenden Abschnitt~\ref{subsection:Rückbildung} betrachten wir sogenannte Rückbildungen wie \textit{not-land}(\textit{en}) aus \textit{Not-land-ung}. Im Anschluss beschäftigen wir uns in Abschnitt~\ref{subsection:Kurzwort} mit Kurzwortbildungen wie \textit{LKW} aus \textit{Lastkraftwagen} und mit \textit{i}-Derivationen wie \textit{Späti} für \textit{Spätkauf}. In Abschnitt~\ref{subsection:Resuplikation} geht es dann um sogenannte Reduplikationen, d.\,h. um Dopplungen von nicht-identischen Bestandteilen wie in \textit{Kuddel-muddel} sowie von identischen Bestandteilen wie in \textit{Milch-Milch} (statt \textit{Hafermilch}).


\subsection{Rückbildung als besondere Konversion}\label{subsection:Rückbildung}

Komplexe Verben wie \textit{not-land}(\textit{en}) (\ref{ea:Notland_en}) und \textit{bau-spar}(\textit{en}) (\ref{ex:Bauspar_en}) werden durch den Abzug eines Suffixes auf ein komplexes Ausgangswort bezogen.

\ea \label{ea:Rückbildung}
\ea \label{ea:Notland_en} Not-land-ung $\rightarrow$ not-land(en)
\ex \label{ex:Bauspar_en} Bau-spar-er $\rightarrow$ bau-spar(en)
\z 
\z 

Die\is{Rückbildung} Suffixkürzung mit anschließendem Wortartwechsel wie in (\ref{ea:Rückbildung}) wird als \textit{Rückbildung} bezeichnet. Die Annahme entsprechender Rückbildungen beruht meist auf Daten, in denen eine komplexe suffigierte Form wie \textit{Notlandung} vor einer rückgebildeten Form wie \textit{notland}(\textit{en}) belegt ist. 

\citet[95]{Erben2006_1975} hält Rückbildung für eine besondere, auf Analogie beruhende Wortbildung. \citet[144]{Echenlohr1999} hält Verbalisierungen von Komposita für ein ganz reguläres Bildungsmuster. Dagegen spricht allerdings, dass viele der so gebildeten Verben defektiv sind. Denn sie werden meist nur im Infinitiv (\ref{ea:notlanden}) oder als Partizip II (\ref{ex:notgelandet}) gebraucht, kaum aber flektiert wie in (\ref{ex:landet_not}) oder in (\ref{ex:notlandetet}).

\ea \label{notlanden}
\ea \label{ea:notlanden} Saint-Exupéry musste 1935 in der Nähe von Kairo notlanden.
\ex \label{ex:notgelandet} Saint-Exupéry ist 1935 in der Nähe von Kairo notgelandet.
\ex \label{ex:landet_not} Saint-Exupéry landete 1935 in der Nähe von Kairo not.
\ex \label{ex:notlandetet} Saint-Exupéry notlandete 1935 in der Nähe von Kairo.
\z 
\z 

Wir halten sogenannte Rückbildungen von Komposita zu Verben für ad hoc gebildete komplexe Stämme, die anschließend konvertiert werden. Dies wird dadurch begünstigt, dass der Kopf des ursprünglichen Kompositums deverbal ist und durch Reduktion des Suffixes wieder seinen Verbalcharakter annimmt. Dies trifft bei Verben wie \textit{video-überwach}(\textit{en}), \textit{kur-pfusch}(\textit{en}) oder eben \textit{not-land}(\textit{en}) und \textit{bau-spar}(\textit{en}) zu. Der Verbalcharakter des Zweitgliedes ist aber, wie das Verb \textit{ohr-feig}(\textit{en}) zeigt, keine Bedingung. Es handelt sich um morphologische Konversionen. 

Komplexe Verben wie \textit{eis-lauf}(\textit{en}), \textit{brust-schwimm}(\textit{en}) oder \textit{kunst-turn}(\textit{en}) analysieren wir mit \citet[245]{Eisenberg2020_Wort} als Rückkonversionen von Komposita, deren Kopf wie bei \textit{Brustschwimmen} ein zum Substantiv konvertierter Infinitiv ist. 

Bei Bildungen wie \textit{staub-saug}(\textit{en}) oder \textit{ehe-brech}(\textit{en}) kommen verschiedene Analysen in Frage. Es kann sich um syntaktische Konversionen via Inkorporation des Objekts handeln. Traditionellerweise wird diese Bildung als \textit{Univerbierung}\is{Univerbierung} bezeichnet. Aus syntaktischer Nachbarschaft wie in \textit{die Ehe brechen} oder \textit{den Staub saugen} entwickelt sich jeweils ein Wort. Das Ergebnis ist der Referenzverlust des Substantivs zugunsten einer rein prädikativen Funktion analog zu \textit{Rad fahren} oder \textit{Zeitung lesen}. Ebenso wäre die suffixlose Reanalyse komplexer Stämme möglich. Von \textit{Staubsauger} wird der komplexe Stamm \textit{staubsaug} kraft des verbalen Zweitgliedes zum komplexen Verbstamm konvertiert. Der Stamm \textit{Ehebruch} wird durch implizite Derivation zu \textit{ehebrech}(\textit{en}) abgeleitet.

Die bisher besprochenen Bildungen gehen auf Komposita zurück, sind aber Derivata, was die Konversion als Bildungstyp einschließt. Es handelt sich weder um verbale Komposita noch um Partikelverben.


\subsection{Kurzwortbildung und \textit{i}-Derivation}\label{subsection:Kurzwort}

Anderer\is{Kurzwortbildung} Art sind die Bildungen in (\ref{ea:Kurzwortbildung}). Es handelt sich um Kurzwörter, die aus einer längeren Vollform abgeleitet sind.

\ea \label{ea:Kurzwortbildung}
\ea \label{ea:LKW} Lastkraftwagen $\rightarrow$ LKW
\ex \label{ex:Krimi} Kriminalfilm $\rightarrow$ Krimi
\ex \label{ex:Azubi} Auszubildender $\rightarrow$ Azubi
\ex \label{ex:U-Bahn} Untergrundbahn $\rightarrow$ U-Bahn
\ex \label{ex:Mikro} Mikrophon $\rightarrow$ Mikro
\z 
\z 

In (\ref{ea:LKW}) handelt es sich um ein sogenanntes Initialwort, denn es besteht aus den Anfangslauten der lexikalischen Stämme der Vollform. Kurzwörter dieser Art werden wie \textit{LKW} buchstabierend gesprochen und/oder wie \textit{FAZ} (Frankfurter Allgemeine Zeitung) zusammengezogen.

Bei \textit{Krimi} (\ref{ex:Krimi}) handelt es sich um das zweisilbige Anfangssegment der Vollform. So gebildete Kurzwörter bestehen wie \textit{Uni}, \textit{Demo}, \textit{Abi}, \textit{Abo} oder \textit{Zivi} aus einer betonten Silbe, auf die eine unbetonte Silbe folgt, was den für das Deutsche typischen trochäischen Versfuß (vgl. Abschnitt~\ref{section:betonte_unbetonte_Silben}) bildet.   

Das Kurzwort \textit{Azubi} (\ref{ex:Azubi}) enthält unterschiedlich große Segmente der Vollform. In (\ref{ex:U-Bahn}) wurde das Erstglied des Kompositums \textit{Untergrundbahn} auf den Anfangslaut gekürzt. Ursprüngliche Komposita sind aber keine Bedingung für Bildungen dieser Art. So beruht \textit{H-Milch} auf der Wortgruppe \gq{haltbare Milche}. Der Bildung \textit{Mikro} (\ref{ex:Mikro}) liegt ein Konfixkompositum zugrunde, welches um das Zweitglied \textit{phon} gekürzt wurde. 

Bis auf den Typ \textit{U-Bahn} (\ref{ex:U-Bahn}) mit intaktem morphologischen Kopf verhalten sich die Kurzwörter anders als die Vollformen. Die auf Vollvokal endenden Kurzwörter wie \textit{Krimi}, \textit{Azubi} und \textit{Mikro} bilden ihren Plural genusunabhängig wie andere auf Vollvokal endende Substantive mit \textit{-s}. Auf Initialen beruhende Kurzwörter wie \textit{LKW} werden im Plural auch flexionslos gebraucht: \textit{zwei LKW} oder \textit{zwei LKWs}. Diese Veränderungen sprechen dafür, die Kurzwortbildung als eigenen Typ anzusehen. Dieser zeichnet sich durch die Kürzung des ursprünglichen morphologischen Kopfes aus.

Nach dem trochäischen Muster sind auch Kurzwörter wie \textit{Späti}, \textit{Studi} und \textit{Ami} gebildet. Aber hier wird die erste Silbe der Stammform mit \textit{-i} kombiniert. Bei mehrsilbigen Stämmen kommt es zu einer Kürzung. Das Modell für jene produktiven Bildungen sind nach \citet[295]{Koepcke2003} zu Kosenamen verkürzte Rufnamen wie \textit{Franzi} von \textit{Franziska} und entsprechende Verlängerungen, wenn der Name einsilbig ist wie bei \textit{Hansi} von \textit{Hans}. Im weiteren Sinne handelt es sich um eine explizite Derivation, die auch von Adjektiven wie \textit{blöd} zu \textit{Blödi}, von Substantiven wie \textit{Gruft} zu \textit{Grufti} und von Verben wie \textit{brumm}(\textit{en}) zu \textit{Brummi} möglich ist. 

Bei verkürzten Stämmen wie \textit{Studi} handelt es sich nach \citet[471]{Fery1997} nicht um eine Ableitung von \textit{Student}. Die Kurzform korrespondiere mit der Stammform \gqq{in dem Sinne, dass sie der Vollform \textit{Student} soweit ähnelt, wie es einem Trochäus mit finalem \textit{i} nur möglich ist}. Das bedeutet, die \textit{i}-Bildungen basieren weniger auf einem regelhaft restringierten Input als vielmehr auf dem trochäischen Resultat der Bildung unter der Bedingung, dass der Bezug zur Vollform gegeben ist.\footnote{\citet[480]{Fery1997} zeigt, inwieweit sich bei Bildungen wie \textit{Studi} auf der einen Seite und \textit{Fundi} auf der anderen Seite verschiedene Prinzipien durchsetzen. Bei \textit{Stu-di} endet die erste Silbe genauso wie die zweite auf einen Vokal, was silbenstrukturell als optimal gilt (vgl. Abschnitt~\ref{subsec:Silbenpräferenz}). Bei \textit{Fun-di} (statt \textit{Fu-ni}) wird mehr von der Vollform \textit{Fundamentalist} übernommen, was zum konsonantischen Ende der ersten Silbe führt. Diese weniger optimale Silbenstruktur erleichtert aber den Rückbezug auf die Vollform.} Das mindert aber nicht den Ableitungscharakter der Basis. Anders als bei der Derivation ist der Input primär phonologisch motiviert, indem die erste Silbe der lexikalischen Stammform abgeleitet wird.

Die wesentliche Frage ist, inwiefern \textit{-i} als Suffix zu charakterisieren ist. Für den Suffix- und damit morphologischen Kopfstatus spricht, dass dieser Typ der \textit{i}-Bildung maskuline Substantive hervorbringt.\footnote{Bei Personenbezeichnungen wie \textit{Mutti}, \textit{Omi} und \textit{Franzi} setzt sich das Sexus gegenüber dem Genus durch. Bildungen wie \textit{Püppi} und \textit{Käppi} zählt \citet[304]{Koepcke2003} zu Diminutivbildungen und grenzt sie von Kosenamen oder pejorativ gebrauchten Personenbezeichnungen wie \textit{Nazi} oder \textit{As}(\textit{s})\textit{i} ab.} Die Funktion der \textit{i}-Ableitung liegt primär in der Personenbezeichnung. Es kann sich wie bei \textit{Franzi} um Eigennamen oder wie bei \textit{Blödi} und \textit{Studi} um Appellativa handeln. Von den Eigennamen für Personen sind Eigennamen für Orte wie \textit{Kotti} von \textit{Kottbusser Tor} motiviert. Von appellativen Personenbezeichnungen sind appellative Bezeichnungen wie \textit{Späti} von \textit{Spätkauf} motiviert. Formal ist die \textit{i}-Derivation blockiert, wenn an der Stelle von \textit{-i} ein heimisches Suffix zur Personenbezeichnung stehen würde. Deshalb kann von \textit{Lehrer} nicht *\textit{Lehri} und von \textit{Prüfling} nicht *\textit{Prüfi} abgeleitet werden.

\subsection{Reduplikation und Komposition}\label{subsection:Resuplikation}
\largerpage
Eine archaische Möglichkeit, neue Wörter zu bilden, war die Verdopplung, die\is{Reduplikation} Reduplikation von lautlichen und morphologischen Einheiten wie in (\ref{ea:Klimbim}) bis (\ref{ex:Tamtam}). Reduplikative Strukturen finden sich auch in der Syntax wie in (\ref{ex:fixfoxi}) und (\ref{ex:xundx}).\footnote{\citet[236]{Kentner2017} bietet eine Übersicht zu reduplizierten Strukturen von Interjektionen wie \textit{haha} bis zu syntaktischen Dopplungen wie \textit{fix und foxi}.}

\ea \label{ea:Reduplikation}
\ea \label{ea:Klimbim} Klimbim, Kuddelmuddel
\ex \label{ex:Singsang} Singsang, Schnickschnack
\ex \label{ex:Tamtam} Tamtam, Pinkepinke
\ex \label{ex:fixfoxi} fix und foxi
\ex \label{ex:xundx} naja, nett und nett (das sind zweierlei Dinge)\footnote{Eine Analyse dieser Struktur bietet \citet{Finkbeiner2012}.}
\z 
\z 

Bildungen wie in (\ref{ea:Klimbim}) bis (\ref{ex:Tamtam}) bezeichnet \citet[470]{Donalies2009Wortbildung} als \gqq{eine Art Echowort durch lautassoziative Dopplung}. Die Herkunft der einzelnen Bildungen ist oft unklar. Die lautassoziative Dopplung kann sowohl den ersten als auch den zweiten Bestandteil betreffen.

In (\ref{ea:Klimbim}) handelt es sich um die Reduplikation des Reims, das heißt um die Verdopplung des Vokals samt silbenschließendem Konsonanten (vgl. Abschnitt~\ref{subsec:Silbenkonstituenten}). Der Konsonant am Silbenanfang (Onset) wird verändert. Bei \textit{Klimbim} zur Bezeichnung von \gq{unbedeutendem Drum und Dran} könnte \textit{klim} durch \textit{klingeln} und \textit{bim} von \textit{bimmeln} motiviert sein. Bei \textit{Kuddelmuddel} für \gq{Durcheinander} könnte \textit{kuddel} vom niederdeutschen Verb \textit{kuddeln} mit der Bedeutung \gq{etwas nicht gründlich waschen} (nur \textit{durchdrücken}) motiviert sein. \textit{Muddel} könnte durch niederdeutsch \textit{Modder} für \gq{Schlamm} motiviert sein.

Die Bildungen in (\ref{ex:Singsang}) sehen wie Komposita aus, deren Zweitglied ein implizites Derivat des Erstglieds oder umgekehrt ist. Die Veränderung zwischen Erst- und Zweitglied betrifft nur den Vokal. Mit \textit{Schnickschnack} von \textit{schnacken} wird \gq{unnötiges Beiwerk} oder \gq{Geschwätz} bezeichnet.

Bei den Bildungen in (\ref{ex:Tamtam}) handelt es sich um vollständige Doppelungen. Das reduplizierte \textit{Tamtam} wurde als Bezeichnung für ein trommelähnliches Musikinstrument aus dem Französischen entlehnt und später metaphorisch für \gq{Lärm}, \gq{Reklame} oder \gq{große Betriebsamkeit} gebraucht. \textit{Pinkepinke} geht auf das niederdeutsche Wort \textit{Pinke} für \gq{Geld} zurück, was durch das klimpernde Geräusch von Münzen motiviert sein könnte. Die wohl bekanntesten Vertreter der vollständigen Reduplikation ohne identifizierbaren Stamm sind \textit{Mama} und \textit{Papa}.

Der Unterschied zur Komposition besteht darin, dass die Bestandteile der Reduplikationen nicht oder nicht mehr einer lexikalischen Stammform zugeordnet werden können. Ist das wie bei \textit{Klimperwimper} möglich, so handelt es sich um Komposita. Anders als bei der Derivation tragen die reduplizierten Elemente wie \textit{warr} in \textit{wirrwarr} keine Funktion im Sinne eines Operators. 

Auf zahlreiche Neubildungen weist \citet[9]{Kentner2017} bei Usernamen im Internet hin: mit Reduplikation des Reims (vgl. Abschnitt~\ref{subsec:Reim} zur Silbenstruktur) wie \textit{Heinzpeinz}, mit Ablaut wie \textit{Wiebkewabke} oder am häufigsten durch vollständige Dopplung wie \textit{Heinheinz}. \citet[912]{Freywald2015} bezeichnet solche Bildungen als ein altes, etabliertes Wortspiel und führt Reduplikationen von Eigennamen aus festen Bildungen seit Anfang des 20. Jahrhunderts an, \zB \textit{Ilsebilse} und \textit{Jochen Pochen}.

Im sogenannten Kiezdeutsch, einer multiethnischen urbanen Jugend-Varietät, finden sich produktive Reduplikationen mit \textit{m-} wie \textit{icke micke}, \textit{Cola Mola}. Nach \citet[325]{Wiese2015} werden diese durch \textit{m}-Reduplikationen im Türkischen gestützt, stellen letztlich aber analoge Bildungen zu \textit{Kuddelmuddel} oder \textit{heckmeck} dar. Solche produktiven, sprachspielerischen Nischen zählen wir nicht zur Wortbildung im engeren Sinne. Denn es handelt sich bei diesen Bildungen weder um Träger abgewandelter Bedeutungen noch um Veränderungen der Wortart. Ihre Funktion liegt in konkreten Kommunikationsbereichen und ist pragmatisch motiviert.

Ein besonderer Fall von Dopplung liegt bei \textit{Mann-Mann} in (\ref{ea:Mann_Mann}) im Gegensatz zu \textit{Männer-Mann} in (\ref{ex:Männer_Mann}) vor.

\ea \label{ea:Redupl_Mann}
\ea \label{ea:Mann_Mann} Bei der Betrachtung der Rodin'schen Plastiken [\ldots{}] lassen wir uns besonders vom Pathos einer Figur beeindrucken. Das ist ein Mann-Mann, heroisch leidend, das Werk heißt L’ombre, der Schatten.\footnote{ Süddeutsche Zeitung, 13.08.2005.}
\ex \label{ex:Männer_Mann} Du hast zu 30 Jahren Falter einen sehr ausführlichen Rückblick geschrieben, da werden 25 Namen erwähnt, 23 Männer, aber nur zwei Frauen, nämlich Doris Knecht und Elisabeth Loibl. Warum ist das so? Bist du eher ein Männer-Mann?\footnote{Falter, 20.02.2019.}
\z 
\z 

In (\ref{ex:Männer_Mann}) handelt es sich um ein Kompositum mit der Besonderheit, dass die Kompositionsstammform \textit{männer-} und die lexikalische Stammform \textit{Mann} auf das gleiche Lexem zurückgehen. Das Gleiche ist bei \textit{Freundes-freund}, \textit{Helfers-helfer} oder \textit{Kindes-kind} der Fall. Die Bedeutung von \textit{Männermann} kann aufgrund des unterspezifizierten Musters der Determinativkomposita (vgl. Abschnitt~\ref{subsection:DetKomp}) paraphrasiert werden als ein Typ von Mann, der in irgendeiner Beziehung zu Männern (als Klassenausdruck) steht. Das kann ein Mann sein, der für oder gegen Männer ist, diese liebt oder bekämpft usw. 

Bildungen wie \textit{Mann-Mann} in (\ref{ea:Mann_Mann}) bestimmen wir mit \citet[299--302]{Hohenhaus2004}, \citet[72]{Donalies2011BasisWB} und \citet[240]{Kentner2017} als einen besonderen Typ von Komposita.\footnote{Die Bezeichnungen sind vielfältig. Als \textit{identical constituent compounds} bezeichnet sie \citet[297]{Hohenhaus2004}. \citet{Ghomeshi_al_2004salad} schreiben von \textit{contrastive focus reduplication}. \citet[72]{Donalies2011BasisWB} nennt sie \textit{Selbstkomposita}. \citet[199]{Stolz_al_2011Redupl} schreiben von \textit{REAL-X-Total-Reduplication}, was auf die Bedeutung \gq{ein echtes, richtiges X} abzielt.} Analog zu typischen Komposita steht auch bei \textit{Mann-Mann} der morphologische und semantische Kopf rechts und links das den Hauptakzent tragende Determinans. Die Bedeutung scheint aber stärker eingeschränkt als bei den unterspezifizierten Determinativkomposita.

Die Bedeutung von \textit{Mann} ist \gq{erwachsener Mensch männlichen Geschlechts}. Die Bedeutung von \textit{Mann-Mann} lässt sich umschreiben als \gq{echter, richtiger Mann}.\footnote{Cf. \citet[312--314]{Hohenhaus2004}.} Wie bei den Determinativkomposita kommt es zu einer Bedeutungsspezialisierung. Nicht jeder erwachsene Mensch männlichen Geschlechts ist ein \textit{Mann-Mann}, wohl aber ein \textit{Mann}. Die enger gefasste Bedeutung bildet eine Unterkategorie von Mann, nämlich den Prototyp. Das ist die Funktion dieses besonderen Musters innerhalb der Determinativkomposita.\footnote{Eine mögliche Paraphrase der Bedeutung von \textit{Mann-Mann}, nämlich \gq{ein Mann von einem Mann} entspricht einer der möglichen Interpretationen eines gewöhnlichen Kompositums mit metaphorischer Abwandlung.} Die Bedeutungsspezialisierung geschieht dadurch, dass typische Werte für Eigenschaften wie Gestalt, Alter, Charakter der Kategorie \gq{Mann} festgelegt werden. Wie die konkrete Wertfestlegung geschieht und welche Eigenschaften sie überhaupt betrifft, bleibt strukturell offen. Diese Offenheit ähnelt der generellen semantischen Unterspezifiziertheit der Determinativkomposita.

Die Festlegung bestimmter Eigenschaften auf bestimmte Werte geschieht erst im entsprechenden Kontext. Im Kontext von (\ref{ea:Mann_Mann}) verstehen wir unter \textit{Mann-Mann} einen Vertreter des Prototyps mit den Merkmalen stark, tapfer, heldenhaft leidend. In (\ref{ex:Mann_Mann:Macho}) hingegen werden dem Vertreter des Prototyps machohafte Merkmale zugeschrieben.

\ea \label{ex:Mann_Mann:Macho} Bernd Buchholz, der jetzt für die FDP in den Bundestag will, war der Mann-Mann, der breitbeinig, selbstsicher und laut durch den Salon geht.\footnote{Süddeutsche Zeitung, 01.12.2012.}
\z 

Eine wesentliche Bedingung für das sinnvolle Verständnis dieser reduplizierten Komposita ist, dass sich der Prototyp von Alternativen, das heißt von anderen (echten oder unechten) Unterkategorien abhebt. Diese Alternativen werden mitunter selbst durch Komposita ausgedrückt. So hebt sich die Bedeutung von \textit{Mutter-Mutter} je nach Kontext von \textit{Stiefmutter}, \textit{Spendermutter} oder \textit{Pflegemutter} ab und fokussiert die leibliche Mutterschaft als Prototyp im Sinne der wörtlichen Bedeutung von \textit{Mutter}. Analog hierzu kann mit \textit{Schlange-Schlange}, die vor einem Supermarkt ist, auf ein Tier als Vertreter der wörtlichen Bedeutung im Gegensatz zur metaphorischen \textit{Warteschlange} verwiesen werden. 

Die Rolle der Alternativen wird auch deutlich, wenn Stammreduplikationen mit gewöhnlichen Determinativkomposita koordiniert (\ref{ea:Lebens_Kunst_Kunst}) und mit diesen kontrastiert (\ref{ex:Tanz_Rock_Rock}) werden. Dies ist ein wichtiges Indiz dafür, diesen Typ von Reduplikation als Komposition zu analysieren.

\ea \label{ea:DetKomp_ReplKomp}
\ea \label{ea:Lebens_Kunst_Kunst} Von anderer Seite kommt der Vorschlag, Lebens-, Liebes- und Kunstkunst getrennt zu diskutieren.\footnote{Die Zeit, 13.11.2008.}
\ex \label{ex:Tanz_Rock_Rock} Wenn ich Rock mache, mit meiner Disko-No.1-Band, ist es nicht Rock-Rock sondern Tanz-Rock.\footnote{ Die Zeit, 11.04.2014.}
\z 
\z 

Auf drei strukturelle Unterschiede zu den Determinativkomposita macht \citet[924--926]{Freywald2015} aufmerksam. Erstens weisen die Erstglieder in reduplizierten Bildungen keine Fugenelemente auf, und zwar auch dann nicht, wenn der entsprechende Stamm eigentlich über eine Kompositionsstammform mit Fugenelement verfügt. Dies wurde bereits bei \textit{Männer-Mann} im Gegensatz zu \textit{Mann-Mann} deutlich. Ein anderes Beispiel ist \textit{Freundesfreund} im Gegensatz zu \textit{Freund-Freund}. 

Zweitens kommen in reduplizierten Bildungen nicht nur lexikalische Stammformen als Erstglieder vor, sondern auch flektierte Pluralformen, die nicht mit entsprechenden Kompositionsstammformen identisch sind. Die Bedingung hierfür ist, dass das Kopfsubstantiv im Plural steht. Ein Beispiel ist \textit{Freunde-Freunde} in (\ref{ea:Freunde_Freunde}).

\ea \label{ea:Freunde_Freunde} Aber meine Freunde-Freunde, also meine wirklich wahren Freunde, mit denen ist das total entspannt.\footnote{Braunschweiger Zeitung, 15.07.2011}
\z 

Drittens werden auch Adverbien (\ref{ea:jetzt_jetzt}) und selten auch Verben (\ref{ex:lieben_lieben}) redupliziert, obwohl sie nur selten Köpfe von Determinativkomposita sind.

\ea \label{ea:Adv_V_Redupl}
\ea \label{ea:jetzt_jetzt} Wenn sie es für die Lösung hielten, sich 15 Jahre lang hochzunetzwerken, bevor sie in der CDU mal als junge Wilde gelten, würden sie ja womöglich nicht fordern, dass jetzt -- also jetztjetzt! -- etwas geschieht.\footnote{der Freitag 14/2019.}
\ex \label{ex:lieben_lieben} Aber die Deutschen investieren ganz besonders in die Ausstattung von Garten, Balkon und Terrasse. Sie liebenlieben ihr Freiluft-Wohnzimmer.\footnote{Der Westen, 08.04.2013.}
\z 
\z 

Aufgrund dieser Unterschiede bestimmten wir die Reduplikation nicht als eigenen Wortbildungstyp, sondern als eine eigenständige Nische der Determinantivkomposita. Die wesentliche Rechtsköpfigkeit und die Bedeutungsspezialisierung bleiben von den Unterschieden unberührt.

Wenn die lexikalische Stammform der Kompositionsstammform entspricht, wie bei \textit{Mutter-Mutter} oder \textit{Kunst-Kunst}, so beinhaltet der Interpretationsspielraum im Rahmen der unterspezifizierten Determinativkomposita auch das prototypische Verständnis. Letzteres ist als Nische von Ersterem zu verstehen. Unter \textit{Kunst-Kunst} kann Kunst verstanden werden, die \zB Kunst zu ihrem Thema hat, aber auch Kunst, die als richtige, echte Kunst gilt. Ausgehend von der gemeinsamen Basis hat sich der reduplizierende Typ formal durch Identität von Zweit- und Erstglied stabilisiert.   


\section{Transparenz, Produktivität, Regel/Analogie, Restriktionen}\label{section:Produktivität}

Bisher haben wir uns mit der morphologischen Analyse bestehender komplexer Wörter beschäftigt. In diesem Abschnitt geht es nun um wesentliche Voraussetzungen für die Bildung neuer komplexer Wörter und um Erklärungen dafür, wie entsprechende Neubildungen geregelt sind. Exemplarisch werden wir uns in den folgenden Abschnitten hauptsächlich mit den Suffixen \textit{-er} wie in \textit{Lehr-er}, mit \textit{-bar} wie in \textit{trink-bar} und mit \textit{-ung} wie in \textit{Plan-ung} beschäftigen. 

Zuerst werden wir in Abschnitt~\ref{subsection:Transparenz} den Transparenzbegriff klären. Denn je transparenter bestehende komplexe Wörter sind, umso eher können wir neue Wörter auf ähnliche Weise, d.\,h. auf der Grundlage desselben Musters bilden. 

Unter dem Stichwort Produktivität werden wir in Abschnitt~\ref{subsection:Produktivität} den Zusammenhang zwischen der Zahl bestehender komplexer Wörter eines Musters und der Wahrscheinlichkeit, dass auf der Grundlage dieses Musters neue Wörter gebildet werden, thematisieren. 

Im Anschluss stellen wir in Abschnitt~\ref{subsection:RegelAnalogieMuster} zwei zentrale Erklärungen für die Neubildung von Wörtern vor und beziehen diese aufeinander, nämlich die Neubildung als analogen Nachbau bestehender komplexer Wörter und die regelbasierte, d.\,h. auf Anwendung einer Regel beruhende Neubildung. 

Im letzten Abschnitt~\ref{subsection:Restriktionen} liegt der Fokus auf Restriktionen der Wortbildung, d.\,h. auf Einschränkungen verschiedener Art, die die Bildung neuer Wörter begrenzen. 


\subsection{Transparenz}\label{subsection:Transparenz}

Morphologisch komplexe Wörter sind unterschiedlich\is{Transparenz} transparent. Das heißt, die Gesamtbedeutung ist unterschiedlich stark durch die Bedeutung der Stammformen bzw. bei Affixen durch die Funktion und durch das entsprechende Bildungsmuster bestimmt. Formal und funktional transparent sind Bildungen wie \textit{Les-er}, \textit{Schreib-er} oder \textit{Fahr-er}. Sie bestehen aus einem Verbstamm und dem Suffix \textit{-er}. Ihre Bedeutung zielt auf die Rolle des Handlungs- und Tätigkeitsträgers ab. Sie lässt sich als jemand, der das vom Verb Bedeutete tut, umschreiben, was ihre Bezeichnung als Nomen agentis begründet. Vollständig transparente Bildungen können als Einheiten im Lexikon gespeichert sein, wenn sie häufig gebraucht werden.

Teilweise transparent\is{Transparenz} sind Bildungen wie \textit{Öffn-er} und \textit{Koch-er}, wenn wir sie als Instrumente und damit als motivierte Ableitungen vom Prototyp des Nomen agentis verstehen. Die Instrumentlesart ergibt sich nicht strikt aus den Bestandteilen. Sie ist aber durch die (auch historisch) primäre Bildung von Nomina agentis motiviert.\footnote{Cf. \citet[524--526]{Dressler1986} und \citet[506]{Booij1986}.}

Derivata wie \textit{Füll-er}, \textit{Fehle-r} und \textit{Hock-er} sind morphologisch\is{Transparenz} transparent. Wir können sie problemlos in eine verbale Basis und das Suffix \textit{-er} zerlegen. Ihre Bedeutung ist jedoch nicht oder nicht mehr transparent. Ein \textit{Füller} (verkürzt aus \textit{Füllfederhalter}) füllt nichts, sondern ist ein Schreibgerät, das mit Tinte gefüllt wird. Ein \textit{Fehler} fehlt im übertragenen Sinne als eine Abweichung vom Richtigen oder allgemein als Mangel. Motiviert war seine Bildung in Analogie zu \textit{Treff-er}: ein Schuss, der trifft oder fehlt, wurde als \textit{Treffer} oder \textit{Fehler} bezeichnet. Ein \textit{Hock-er} ist alltagssprachlich eine Sitzgelegenheit für eine Person ohne Rückenlehne.\footnote{Sowohl umgangssprachlich als auch fachsprachlich kann \textit{Hocker} semantisch transparent(er) verstanden werden als Mensch, der müßig an einem Ort verharrt oder in der Archäologie als Skelett, welches mit angezogenen Beinen, also wie hockend, im Grab liegt.} Die heutigen Bedeutungen sind zwar motiviert, aber nicht aus den Bestandteilen erschließbar. Deshalb müssen diese Bildungen im Lexikon gespeichert sein.

Schließlich gibt es Substantive\is{Transparenz} wie \textit{Hammer}, \textit{Eimer} oder \textit{Teller}, die morphologisch komplex aussehen, aber, zumindest heute, nicht mehr morphologisch transparent sind. Diese Substantive können nicht oder nicht mehr in einen Stamm und das Suffix \textit{-er} zerlegt werden. Denn entsprechende Stämme gibt es nicht. Interessanterweise verhalten sich aber die morphologisch transparenten und die nicht transparenten in der Flexion gleich. So handelt es sich bei \textit{Leser} aber auch bei \textit{Eimer} und \textit{Hammer} um Maskulina, die ihren Genitiv mit \textit{-s} bilden, im Plural aber (anstelle des erwartbaren \textit{-e}) endungslos bleiben. Aufgrund dieses Parallelverhaltens bezeichnet \citet[228]{Eisenberg2020_Wort} die Endung \textit{-er} in morphologisch nicht transparenten Wörtern als \gqq{Pseudoaffix}.


\subsection{Produktivität}\label{subsection:Produktivität}

Bisher haben wir beiläufig bestimmte Bildungen als eher\is{Produktivität} produktiv oder eher unproduktiv bezeichnet. Was heißt das? Im wörtlichen Sinne produktiv sind letztlich nur Sprecherinnen und Schreiber. Im übertragenen Sinne ist ein Wortbildungsmuster oder, verkürzend gesprochen, ein Affix, produktiv, wenn es unbegrenzt viele neue Bildungen hervorbringen kann. Unbegrenzte Bildungsmöglichkeiten gibt es \zB bei den rekursiven Determinativkomposita mit substantivischen Bestandteilen wie \textit{Haus-tür} und bei den syntaktischen Konversionen vom verbalen zum substantivischen Infinitiv wie (\textit{das}) \textit{Lesen} oder vom adjektivischen zum substantivischen Partizip I wie (\textit{der}) \textit{Lesende}. Normalerweise sind Wortbildungsmuster nur unter bestimmten Bedingungen, also eingeschränkt produktiv. So werden mit dem Suffix \textit{-ung} heute vorrangig die Stämme telischer, meist komplexer Verben abgeleitet, nicht aber Adjektive und Substantive. Wir kommen hierauf zurück. 

Nach \citet[205--211]{Bauer2001} zielt\is{Produktivität} der Begriff der Produktivität sowohl auf die Zahl der in einem bestimmten Zeitraum belegten Neubildungen als auch auf die Wahrscheinlichkeit neuer Bildungen ab. Belegte Neubildungen und ihr Verhältnis zu bereits etablierten Bildungen desselben Musters sind ein wichtiges Indiz für die Produktivität eines Musters. Liegt der Fokus auf der Wahrscheinlichkeit neuer Bildungen, dann ist Produktivität nur graduell zu fassen.\footnote{Einen knappen Überblick zu den verschiedenen Perspektiven auf Produktivität zwischen Ist-Bestand und Potential gibt \citet{Plag2006}. Verschiedene Modelle zur Berechnung verschiedener Produktivitätsmaße diskutiert \citet{Baayen2009}.}


\subsection{Regel, Analogie und Muster}\label{subsection:RegelAnalogieMuster}

Produktivität beruht grundlegend auf zwei Bildungsprozessen. Entweder wird eine Regel angewendet oder aber ein neues Wort wird in Analogie\is{Analogie!versus Regel?} zu einem bestehenden Wort \gq{nachgebaut}. Als regelbasiert lässt sich die mögliche Bildung\is{Suffix!-\textit{er}} \textit{Grüß-er} beschreiben. Der Stamm eines agentiven Verbs wird mittels \textit{-er} zum Nomen agentis abgeleitet. Analogiebasiert lässt sich die Bildung von \textit{Grüß-er} auf die Beziehung zwischen dem etablierten \textit{Les-er} und \textit{les}(\textit{en}) beziehen und auf Stämme nicht agentiver Verben wie bei \textit{Stink-er} übertragen. Analoge Nachbildung kann auch Stämme anderer Wortarten betreffen. So lässt sich \textit{Schül-er} analog zu \textit{Lehr-er} analysieren, obwohl ersterer eine substantivische, letzterer eine verbale Basis hat.

Regel und Analogie verstehen wir wie \citet[23--25]{Motsch2004} als komplementäre Begriffe und nicht wie \citet[67]{FleischerBarz2012} als konträr. Analogiebildung ist die Fähigkeit, Ähnlichkeiten zwischen Wörtern wie \textit{Les-er} und \textit{les}(\textit{en}) zu erkennen, Unterschiede zu abstrahieren und das Verhältnis auf andere Entitäten zu übertragen. Die Abstraktion von lexikalischen Einzelfällen führt zu Regeln. Somit sind Regel und Analogie komplementäre Begriffe, was \citet[110]{Paul1880} erkannt und auch pädagogisch ausformuliert hat:

\begin{quote}
Kein Lehrer aber, der nicht ganz unpädagogisch verfährt, wird es versäumen zugleich Beispiele für die Regel [\ldots{}] zu geben. Regel und Muster ergänzen sich gegenseitig in ihrer Wirksamkeit; und man sieht aus diesem pädagogischen Verfahren, dass dem konkreten Muster gewisse Vorzüge zukommen müssen, die der abstrakten Regel abgehen.
\end{quote}

Wortbildungsregeln werden traditionell als\is{Wortbildungsmuster} Wortbildungsmuster bezeichnet.\footnote{Cf. \citet[67--68]{FleischerBarz2012} und \citet[15]{Motsch2004}. Je nach theoretischem Hintergrund werden unterschiedlich abstrakte Modelle und Subtypen angenommen. Dabei geht es \zB um die Frage, ob\is{Suffix!-\textit{er}} \textit{Leser} (Nomen agentis), \textit{Kocher} (Instrument) und \textit{Lacher} (Verbalsequenz) auf einem, zwei oder drei Modellen basieren.} Es handelt sich um Strukturschemata, welche bei der Derivation aus einer Derivationsstammform und aus einem positionsfesten Affix bestehen. Das Affix wandelt die Bedeutung der Basis ab und legt dabei die Wortart des komplexen Wortes fest. Ein Wortbildungsmuster dient primär der Analyse bestehender komplexer Wörter und kann unter bestimmten Bedingungen auch zur Neubildung produktiv genutzt werden. Dies setzt die Abstraktion konkreter Derivationsstammformen zu beliebigen oder bestimmten Vertretern einer Wortart voraus.

Semantische und\is{Produktivität!Voraussetzungen} morphologische Transparenz bestehender Bildungen sind die Voraussetzung dafür, dass das ihnen zugrunde liegende Muster produktiv sein kann. Ein wesentlicher Faktor für Produktivität ist, dass mit einem Muster bereits viele transparente Komplexe gebildet wurden. Nur aus bestehenden Bildungen kann durch Bezug auf die Basis das entsprechende Muster abstrahiert werden. So wie sich\is{Suffix!-\textit{er}} \textit{Les-er} zu \textit{les}(\textit{en}) in Beziehung setzen lässt, so verhalten sich auch \textit{Schreib-er} zu \textit{schreib}(\textit{en}), \textit{Fahr-er} zu \textit{fahr}(\textit{en}), \textit{Spiel-er} zu \textit{spiel}(\textit{en}) usw. Je mehr verschiedene lexikalische Basen mit einem Affix abgeleitet wurden, umso eher kann von der konkreten Basis zu einem Muster wie in (\ref{ea:NomenAgentisMuster}) abstrahiert werden.     

\ea \label{ea:NomenAgentisMuster} [[X]\textsubscript{Verb}\textit{-er}]\textsubscript{Subst. Mask} : Person, die tut, was das Verb bedeutet
\z 

Zwar können wir auch eine Bildung wie \textit{Fahrt} auf das Verb \textit{fahr}(\textit{en}) beziehen. Nur kommt es nicht zur Musterabstraktion, da es nur wenige Bildungen dieses Typs gibt, welche zudem wie \textit{nähen}/\textit{Naht}, \textit{sehen}/\textit{Sicht} verschiedene Derivationsstammformen aufweisen. Dadurch und durch Lexikalisierung der \textit{t}-Derivata ist \zB bei \textit{Bucht} von \textit{bieg}(\textit{en}), \textit{Sucht} von \textit{siech}(\textit{en}) oder \textit{Zucht} von \textit{ziehen} die entsprechende Transparenz nicht mehr oder nur schwach gegeben.\footnote{\citet[228]{Eisenberg2020_Wort} schlägt für nicht mehr produktive Affixe wie \textit{-t} die Bezeichnung \gqq{morphologischer Rest} vor.}

Auch die\is{Produktivität!Voraussetzungen} Entwicklung eines Affixes kann einen Einfluss auf die Produktivität des zugrunde liegenden Musters haben. So hat sich aus einigen Derivata mit\is{Suffix!-\textit{er}} \textit{-er} die Affixvariante \textit{-ler} herausgebildet. Das ursprünglich zur Basis gehörende \textit{-l} wurde als zum Affix gehörig reanalysiert. \citet[292--293]{Wilmanns1896Wb} führt Bildungen wie \textit{Gürtler} von \textit{Gürtel} sowie \textit{Säckler}\footnote{Mit \textit{Säckler} wird bis heute im Schwäbischen eine Person bezeichnet, die berufsmäßig Leder vernäht.} vom Diminutiv \textit{Säckel} auf. Spätere Bildungen wie \textit{Tisch-ler}, \textit{Wissenschaft-ler} oder \textit{Künst-ler} sind direkte Ableitungen mittels \textit{-ler}, welches sich im Gegensatz zu \textit{-er} auf substantivische Basen spezialisiert hat.


\subsection{Restriktionen}\label{subsection:Restriktionen}

In diesem Abschnitt geht es um Restriktionen, also um Einschränkungen der Produktivität eines Musters, die auf ganz verschiedenen Ebenen liegen. Wir werden uns der Reihe nach mit wortartbasierten, semantischen, valenziellen, pragmatischen und phonologischen Restriktionen beschäftigen.  

Ein zentraler Faktor, der die Produktivität\is{Produktivität!kategoriale Restriktion} eines Musters einschränkt, ist die Wortart der Ableitungsbasis.\footnote{\citet[128]{Bauer2001} hält diese kategoriale Festlegung für die wesentliche Produktivitätsrestriktion. Wortartspezifische Restriktionen werden nach \citet[47--48]{Aronoff1976} unter dem Stichwort \textit{unitary base hypothesis} diskutiert. Zur kritischen Diskussion verweisen wir auf \citet[44--45]{Plank1981}.} Die in (\ref{ea:NomenAgentisMuster}) aufgeführte Regel ist auf verbale Derivationsstammformen beschränkt. Verwandte\is{Suffix!-\textit{er}} Bildungen wie \textit{Kutsch-er}, \textit{Schäf-er}, \textit{Töpf-er} oder \textit{Schül-er} gehen auf Substantive zurück und basieren damit auf einer anderen Regel. Allerdings ist nicht zu klären, ob eine Bildung wie \textit{Schül-er} direkt in Analogie zu \textit{Lehr-er} mit verbalem Stamm oder als Output einer anderen, verwandten Regel zu analysieren ist. Aufgrund der semantischen Nähe zwischen \textit{Les-er} als \gq{jemand, der liest} und \textit{Kutsch-er} als \gq{jemand, der eine Kutsche führt} oder \textit{Töpf-er} als \gq{jemand, der Töpfe herstellt} kann mit \citet[48--49]{Scherer2005} auch von \textit{einer} Regel ausgegangen werden, die sowohl auf verbale als auch auf substantivische Stämme angewendet werden kann.\footnote{\citet[48--49]{Scherer2005} differenziert die Bildungen auf der Grundlage \textit{eines} Musters prototypisch, indem sie zwei Kriterien ansetzt. Substantivische \textit{er}-Ableitungen basieren auf \textit{einem} Muster, wenn sie Maskulina der starken Flexionsklasse sind und wenn sie Nomina agentis (Semantik) sind und/oder wenn ein Wortartwechsel (Syntax) vorliegt.}

Die wortartbasierte Beschränkung\is{Produktivität!Outputorientierung} wird auch relativiert, wenn wir Bildungsreihen wie \textit{an-kleben}, \textit{an-nähen}, \textit{an-schweißen}, \textit{an-seilen}, \textit{an-zwecken} und \textit{an-kleistern} betrachten. In allen Fällen handelt es sich um transparente Partikelverben zum Ausdruck einer Handlung, die der Kontaktherstellung bzw. davon abgeleitet der Befestigung dient. Die Basen aber sind verbaler und substantivischer Wortart. Mit \citet[44--45]{Plank1981} gehen wir von outputbasierten Bildungsmustern aus. Wichtiger als die Wortart der Ableitungsbasis ist das semantische und formale Resultat der Wortbildung. Die Ausweitung eines Musters auf Stammformen verschiedener Wortarten bezeichnet \citet[14]{Wilmanns1896Wb} als Eröffnen \gqq{neue[r] fruchtbare[r] Gebiete}, modern gesprochen geht es um die Ausdehnung eines Prototyps als Vorläufer.\footnote{Cf. \citet[295]{Langacker1991}: \gqq{Extension from the prototype occur [\ldots{}] because of our proclivity for interpreting the new or less familiar with reference to what is already well established; and from the pressure of adapting a limited inventory of conventional units to the unending, ever-varying parade of situations requiring linguistic expression.}} Wesentlich ist, dass die Bildung im Sinne der Zielkategorie des Musters interpretiert werden kann.

Die Interpretierbarkeit im Sinne der Zielkategorie\is{Produktivität!semantische Restriktion} führt zu einer semantischen Beschränkung innerhalb einer möglichen Wortartbeschränkung. Wir können nicht aus allen Verben mit Hilfe des Suffix\is{Suffix!-\textit{er}} \textit{-er} Substantive bilden. Bildungen wie *\textit{Sterb-er}, *\textit{Wächs-er} bzw. *\textit{Wachs-er} oder *\textit{Ertrink-er} sind aufgrund der Bedeutung ihrer Basen nicht möglich. Wer stirbt, wächst oder ertrinkt, tut sicher auch etwas, aber die Verben \textit{sterben}, \textit{wachsen} und \textit{ertrinken} beschreiben Vorgänge und keine Handlungen oder Tätigkeiten. Die vom Subjekt bezeichnete Entität kontrolliert den Vorgang nicht, in den sie involviert ist. In Abschnitt~\ref{subsection:Verbklassen} haben wir diese Rolle als Vorgangsträger im Gegensatz zum Handlungs- und Tätigkeitsträger (Agens) bezeichnet. Typischerweise ist die \textit{er}-Ableitung auf agentive Verben beschränkt, deren erste Ergänzung ein Handlungs- oder Tätigkeitsträger, aber kein Vorgangsträger ist. 

Restriktionen betreffen\is{Produktivität!valenzielle Restriktion} ebenfalls die Valenz der Basis. So ist der Prototyp des\is{Suffix!-\textit{bar}} \textit{bar}-Adjektivs von einem transitiven (mit Akkusativobjekt) Handlungsverb abgeleitet. Deshalb können wir regelkonform \textit{unterstütz-bar} aber nicht *\textit{helf-bar} bilden, \textit{wachklingel-bar} aber nicht \textit{klingel-bar}, \textit{verlänger-bar} aber nicht *\textit{dauer-bar}.     

\citet[54]{Riehemann1998} zeigt durch eine empirische Untersuchung, dass besonders häufig gebrauchte\is{Suffix!-\textit{bar}} \textit{bar}-Adjektive (7\% der untersuchten Typen) von dieser Restriktion abweichen und/oder nicht vorhersagbare Bedeutungsaspekte aufweisen wie \textit{zahl-bar} als \gq{muss gezahlt werden} statt \gq{kann gezahlt werden}. So gibt es \textit{unentrinn-bar}, obwohl das Basisverb den Dativ regiert (\textit{jemandem oder einer Sache entrinnen}), \textit{verfüg-bar} trotz Präpositionalrektion des zugrunde liegenden Verbs (\textit{über etwas verfügen}), \textit{brenn-bar}, dessen Basis intransitiv ist (\textit{etwas brennt)} und \textit{sicht-bar} mit substantivischer Basis. 

Solche, häufig älteren Bildungen setzen Restriktionen nicht außer Kraft, solange keine heute produktive Reihenbildung nachgewiesen werden kann. Die Bildungen verstoßen zwar gegen Restriktionen, aber sie teilen mit den vollkommen regulären Ableitungen den Ausdruck der Möglichkeit eines Vorgangs. Solange die Beziehung zwischen dem Stamm und der Ableitung transparent ist, können einzelne, von der Regel abweichende Bildungen, selber als Modelle für analoge Nachbildungen dienen. Sie eröffnen damit eine gewisse Produktivität an der Transitivitäts-Restriktion vorbei. Im Einzelfall kann das wie bei der werbesprachlichen (bewusst die Regel verletzenden) Bildung \textit{un-kaputt-bar} zu Nachbildungen wie \textit{un-platt-bar} führen. 

Die Bildung von \textit{bar}-Adjektiven aus intransitiven Verben wie\is{Suffix!-\textit{bar}} \textit{brenn-bar}, \textit{verrott-bar} oder \textit{verheil-bar} ist insofern motiviert, als diese Adjektive ausdrücken, dass eine Entität Träger des vom Basisverb beschriebenen Vorgangs sein kann. Etwas Brennbares kann brennen, ohne den Vorgang des Brennens zu kontrollieren. Das ist auch bei etwas Trinkbarem der Fall. Es kann getrunken werden, ohne den entsprechenden Vorgang zu kontrollieren. Adjektive, die mit \textit{-bar} aus intransitiven Vorgangsverben abgeleitet wurden, ähneln in ihrer Bedeutung der ebenfalls intransitiven passivischen Verbform von transitiven Verben. Ähnlich wie passivische Verbformen transitiver Verben beschreiben auch reflexive Medialverben (vgl. Abschnitt~\ref{subsection:Reflexiva}) Vorgänge. Selten werden auch diese mit \textit{-bar} abgeleitet: \textit{erneuer-bar} von \textit{sich erneuern}, \textit{verflüchtig-bar} von \textit{sich verflüchtigen}. Typische \textit{bar}-Adjektive wie \textit{trinkbar} beziehen sich auf verursachte Vorgänge wie (\textit{von jemandem}) \textit{getrunken werden}, weniger typische wie \textit{brennbar} oder \textit{erneuerbar} auf nicht verursachte, d.\,h. von alleine ablaufende Vorgänge wie \textit{brennen} und \textit{entrinnen}. 

Diese Erklärung gilt nicht für Verwendungen von \textit{flieg-bar} in (\ref{ea:fliegbar}) oder \textit{schlafbar} in (\ref{ex:schlafbar}).

 \ea \label{ea:schlafbar_fliegbar}
 \ea \label{ea:fliegbar} Man merkte den Piloten förmlich an, dass sie sich nach fliegbarem Wetter sehnten.\footnote{ Braunschweiger Zeitung, 05.06.2010.}
 \ex \label{ex:schlafbar} Wer sein Schlafzimmer in urbaner Optik gestalten möchte, dem bieten aus Paletten gefertigte Betten einen hippen schlafbaren Untersatz.\footnote{Nürnberger Nachrichten, 15.06.2019.}
 \z 
 \z 

 Hier handelt es sich (noch) um Regelverstöße. Denn diese \textit{bar}-Adjektive beziehen sich nicht auf einen Vorgangsträger, sondern auf Umstände entsprechender Vorgänge. Bei fliegbarem Wetter kann geflogen werden, aber das Wetter kann nicht geflogen werden. Auf einem schlafbaren Untersatz kann geschlafen werden, aber der Untersatz kann nicht geschlafen werden. Regelkonform wären Verwendungen wie \textit{ein fliegbares Flugzeug} als ein Flugzeug, das geflogen werden kann oder \textit{der schlafbare Schlaf der Gerechten} als der Schlaf der Gerechten, der geschlafen werden kann.  

Eine weitere, aber nicht grammatische Einschränkung der Produktivität eines Musters ist\is{Produktivität!pragmatische Restriktion} pragmatischer Natur. Appellative Wörter (also keine Eigennamen) dienen der sinnvollen Kategorisierung von Entitäten, die wir mit ihnen benennen. So sind Substantive mit\is{Suffix!-\textit{er}} (\textit{-n}/\textit{l})\textit{er} nur dann ableitbar, wenn es eine entsprechende Kategorie \zB im Sinne eines Berufsbildes gibt. Das ist bei \textit{Töpf-er}, \textit{Glas-er} und \textit{Gärtn-er} der Fall, aber nicht bei ?\textit{Teller-er} oder ?\textit{Park-ner}. Diese Bildungen sind potenziell möglich, aber ohne entsprechenden Kontext nicht sinnvoll als Kategorisierung von Personen zu verstehen.

Mit potenziell möglichen Derivata\is{Suffix!-\textit{er}} wie \textit{Aufsteh-er} oder \textit{Geh-er} können wir kaum sinnvoll auf Typen von Tuenden verweisen. Denn \textit{aufstehen} und \textit{gehen} bieten nur schwer kategoriale Abgrenzungsmerkmale. Dass es sich hierbei um pragmatische und nicht um grammatische Beschränkungen handelt, wird deutlich, wenn wir plausible Kontexte wie in (\ref{ea:Aufsteher}) und (\ref{ex:Geher}) betrachten. Mitunter leisten Komposita die entsprechende Kategorisierung wie in (\ref{ex:Frühaufsteher}). 

\ea \label{ea:pragmatische_Restriktion}
\ea \label{ea:Aufsteher} Aber Kinofans teilen sich untereinander nochmal auf, nach jeder Filmvorführung nämlich, in Aufsteher und Sitzenbleiber.\footnote{Berliner Zeitung, 14.02.2002.}
\ex \label{ex:Geher} Auf Grund der unterschiedlichen Geschwindigkeiten mussten die Geher bereits gestern um vier Uhr starten, während die Radler um 8 Uhr los rollten und die Läufer sich erst um 21 Uhr in Bewegung setzten.\footnote{Mannheimer Morgen, 28.06.2000.}
\ex \label{ex:Frühaufsteher} Einige Frühaufsteher trainieren sogar morgens um fünf Uhr vor der Arbeit.\footnote{
Süddeutsche Zeitung, 11.01.2003.}
\z 
\z 

Weitere Restriktionen\is{Produktivität!lautliche Restriktion} können phonologisch, das heißt lautsystematisch motiviert sein. So lassen sich heute \zB originäre Adjektive mit\is{Suffix!-\textit{heit}} \textit{-heit} zu Substantiven ableiten, wenn das Adjektiv auf der letzten (oder einzigen) Silbe betont wird. Das spiegeln Bildungen wie \textit{Gleich-heit} und \textit{Gesund-heit} im Gegensatz zu nicht möglichen Bildungen wie *\textit{Heiter-heit} oder *\textit{Brüderlich-heit} wieder. Wenn die Betonung nicht auf der letzten Silbe liegt, wird wie bei \textit{Heiter-keit} oder \textit{Brüderlich-keit} mit dem Suffix\is{Suffix!-\textit{keit}} \textit{-keit} abgeleitet.\footnote{Von dieser Regel weichen mit \textit{-ern} abgeleitete Adjektive ab, die heute durchgängig wie \textit{Albern-heit} oder \textit{Lüstern-heit} mit \textit{-heit} abgeleitet werden. Die phonologische Restriktion wird affixspezifisch außer Kraft gesetzt.}

Die hier exemplarisch aufgeführten Restriktionen schränken die Produktivität eines Bildungsmusters ein. Mit steigender Zahl an Restriktionen sinkt die Produktivität. Diese Entwicklung wird durch Konkurrenz zwischen verschiedenen Mustern befördert. Bei zwei konkurrierenden Mustern kann es zu einer semantischen Spezialisierung eines der Muster kommen. Eine Folge hiervon kann die Zunahme spezifischer Restriktionen sein. 

So zeigt\is{Produktivität!Konkurrenz} \citet[396--397]{Demske2000}, dass die Produktivität des deverbalen\is{Suffix!-\textit{ung}} \textit{ung}-Musters vom Frühneuhochdeutschen zum Gegenwartsdeutschen abgenommen hat. Hierfür macht sie die Abnahme der Transparenz entsprechender Bildungen verantwortlich. Diese zeigt sich \ua bei Bildungen wie \textit{Fest-ung} oder \textit{Los-ung}, deren verbale Basen heute nicht mehr bestehen und die wir als Objekte (ursprünglich als Resultate der von den Verben bedeuteten Handlungen) verstehen. Ursprünglich bestand die Transparenz zwischen Verb- und \textit{ung}-Ableitung darin, dass letztere auf die verbale Verlaufsbedeutung ohne zeitliche Eingrenzung (\ref{ea:Planung_Verlauf}) und/oder auf ein abgeschlossenes Ereignis (\ref{ex:Planung_Abschluss}) abzielte.   

\ea \label{ea:Planung}
\ea \label{ea:Planung_Verlauf} Engpässe gibt es immer mal, aber eigentlich ist da eine ständige Planung und Abstimmung.\footnote{Nürnberger Nachrichten, 08.07.2009.}
\ex \label{ex:Planung_Abschluss} Im Sommer soll die Planung fürs Gesundheitszentrum abgeschlossen sein.\footnote{Rhein-Zeitung, 09.03.2009.}
\z 
\z 

Die\is{Suffix!-\textit{ung}} zeitlich begrenzte Ereignislesart in (\ref{ex:Planung_Abschluss}) ist heute weitaus typischer als die in (\ref{ea:Planung_Verlauf}). Die zeitliche Begrenzung in (\ref{ex:Planung_Abschluss}) geht einher mit semantischen Restriktionen, welche heute bei durativen Verben wie \textit{leiden} und bei iterativen Verben wie \textit{gackern} die \textit{ung}-Ableitung erschweren. Die Präferenz für die zeitlich begrenzte Ereignislesart geht damit einher, dass vorwiegend transitive und/oder zielbezogene Präfix- und Partikelverben als Basis in Frage kommen. So besteht zu dem intransitiven Verb \textit{arbeiten} das stammkonvertierte Substantiv \textit{Arbeit}. Das transitive Präfixverb \textit{be-arbeiten} und das transitive Partikelverb \textit{ein-arbeiten} können durch \textit{-ung} zu \textit{Bearbeit-ung} und \textit{Einarbeit-ung} abgeleitet werden. Analog zu den Partikelverben werden verbale Basen mit Objektsprädikativ durch \textit{-ung} zu Substantiven wie \textit{Klarstell-ung} oder \textit{Freilass-ung} abgeleitet. 

Die der Verbbedeutung nächste Lesart als zeitlich unbegrenztes Ereignis (\ref{ea:Planung_Verlauf}) wird heute durchgängig und transparent durch die syntaktische Konversion des Infinitivs wie in (\ref{ea:Planen_Verlauf}) geleistet.

\ea \label{ea:Planen_Verlauf} Manchmal stört mich das ständige Planen.\footnote{Spiegel-Online, 21.06.2006.}
\z 

Ein\is{Suffix!-\textit{ung}} Vorteil konvertierter Infinitive liegt in der relativ eindeutigen zeitlich unbegrenzten Vorgangslesart. Mit dieser geht eine generische (unzählbare) Lesart (\ref{ea:Irren}) gegenüber der partikularen (zählbaren) Lesart der \textit{ung}-Ableitung (\ref{ex:Irrungen}) einher.

\ea \label{ea:Irren_Irrungen}
\ea \label{ea:Irren} Irren ist menschlich.
\ex \label{ex:Irrungen} Irrungen, Wirrungen\footnote{Romantitel von Theodor Fontane, 1888.}
\z 
\z 

Auf Musterkonkurrenz\is{Produktivität!Konkurrenz} basierende Restriktionen können auch sprachnormativ motiviert sein. Hierauf weisen \citet[134]{FuhrhopWerner2016} in Bezug auf \textit{er}-Ableitungen in Konkurrenz zu entsprechenden substantivierten Partizipien hin: \textit{Leser} versus \textit{Lesender} und \textit{Leserin} versus \textit{Lesende}. Das Motiv ist insofern normativ, als die Pluralformen der ursprünglichen Adjektive \textit{Lesende} oder \textit{Schreibende} genus- und somit sexusneutral sind. Im Singular bieten die Konvertate gegenüber den Derivata den beabsichtigten Vorteil nicht.\footnote{Diese Musterkonkurrenz geht dem Genderdiskurs allerdings voraus. Es sei an Bildungen wie \textit{Reisender}, \textit{Reisende} oder \textit{Vorsitzende}, \textit{Vorsitzender} gedacht.}

Diese Motivation beschränkt sich nicht auf \textit{er}-Ableitungen an sich, sondern betrifft genusmarkierte Ableitungen von Nomina agentis im Allgemeinen: \textit{Student} und \textit{Studentin} versus \textit{Studierende}, \textit{Demonstrant} und \textit{Demonstrantin} versus \textit{Demonstrierende} sowie \textit{Publizist} und \textit{Publizistin} versus \textit{Publizierende}. Gewinn und Verlust sind gelassen zu beobachten. Können und wollen wir auf die Differenzierung zwischen \textit{Student} bzw. \textit{Studentin} als Inhaber und Inhaberin eines gewissen Status zugunsten von \textit{Studierende} als \gq{Menschen, die gerade studieren} verzichten? Von diesen Abwägungen unbetroffen sind lexikalisierte Bildungen wie \textit{Stürmer} und \textit{Stürmerin} oder \textit{Denkerin} und \textit{Denker}. Diese lassen sich nicht durch \textit{Stürmende} und \textit{Denkende} ersetzen.

Nicht grammatischer\is{Produktivität!lexikalische Blockierung}  oder diskursiver Natur ist die Blockierung ansonsten produktiver Muster durch lexikalisierte Konkurrenzformen. So ist \textit{Stehl-er} zwar die regelkonforme Ableitung von \textit{stehl}(\textit{en}), aber diese ist durch das Wort \textit{Dieb} blockiert. Das komplexe Wort \textit{Koch-er} kann durch die Konkurrenzformen \textit{Koch} und \textit{Köch-in} anders als \textit{Öffn-erin} und \textit{Öffn-er} nicht mehr als \gq{Person, die professionell kocht} verstanden werden. Die Bildung \textit{Wärt-er} ist eine Ableitung von \textit{wart}(\textit{en}) im Sinne von \gq{betreuen}. Die entsprechende Ableitung von \textit{wart}(\textit{en}) im Sinne von \gq{harren} ist dadurch blockiert. Das systematisch bildbare Substantiv \textit{Kalt-heit}\footnote{\textit{Kaltheit} ist eine seit dem Mittelhochdeutschen belegte Konkurrenzform zu der gängigeren Bildung \textit{Kälte}. Cf. Deutsches Wörterbuch von Jacob Grimm und Wilhelm Grimm. Digitalisierte Fassung im Wörterbuchnetz des Trier Center for Digital Humanities, Version 01/21, \url{https://www.woerterbuchnetz.de/DWB?lemid=K00767}, abgerufen am 29.03.2021.} ist durch die bestehende Ableitung \textit{Kält-e} mehr oder weniger blockiert. Die Ableitung \textit{Erlaub-ung} ist durch \textit{Erlaub-nis} blockiert.\footnote{Diese lexikonbasierte (vs. grammatische) Restriktion wird seit \citet[16]{Kiparsky1983} als \textit{Avoid Synonymy}-Prinzip diskutiert: \gqq{The output of a lexical rule may not be synonymous with an existing lexical item.}}

Dass diese Art der Einschränkung im Lexikon (und nicht in der Grammatik) begründet liegt, zeigt sich daran, dass sich nicht lexikalisierte Bildungen nicht gegenseitig blockieren. So wird \textit{Rauch-erei} neben \textit{Ge-rauch-e} gebildet und \textit{Häus-chen} neben \textit{Häus-lein}. Bildungen, die auf nahezu uneingeschränkt produktiven Mustern beruhen, blockieren sich nicht. 

Ein weiteres Indiz dafür, dass es sich um lexikalische und nicht um grammatische Restriktionen handelt, ist, dass mitunter neue Wörter trotz bestehender lexikalisierter Varianten gebildet werden. Das kann wie bei \textit{Stehler} anstelle von \textit{Dieb} stilistisch motiviert sein (\ref{ea:Stehler}) oder wie bei \textit{Kaltheit} anstelle von \textit{Kälte} durch metaphorischen oder einfach Aufmerksamkeit heischenden Gebrauch (\ref{ex:Kaltheit}).
 
 \ea \label{ea:Stehler_Kaltheit}
\ea \label{ea:Stehler} Der Hehler ist schlimmer als der Stehler.\footnote{Sprichwort.}
\ex \label{ex:Kaltheit} Clany wählt eine vom Geschehen distanzierte Perspektive und vermeidet bewusst moralische Fingerzeige und tragische Einzelschicksale, beschreibt dagegen den Krieg in seiner strategischen Kaltheit.\footnote{Belegposition 530616173 im deutschen Teil des \textit{Web as Corpus} (DeWaC). Durchsucht im CQP-Interface der Humboldt-Universität zu Berlin. \url{https://korpling.german.hu-berlin.de/cqpwi}. 29.03.21.}
\z 
\z 

Wie wir gesehen haben, wird die Möglichkeit der Neubildung von Wörtern durch verschiedene Faktoren beschränkt. Diese können, wie im letzten Fall, lexikalischer Natur sein, aber auch kategorial bezüglich der Wortart des abzuwandelnden Stamms  bestimmt sein. Hinzu kommen semantische, valenzielle, pragmatische und phonologische Beschränkungen. Je weniger Restriktionen für ein Wortbildungsmuster gelten, umso produktiver kann es für Neubildungen genutzt werden. Eine wesentliche Voraussetzung für die mögliche Produktivität eines Musters ist die Transparenz möglichst vieler bestehender Bildungen. Je transparenter bestehende Wortbildungen sind, umso eher können neue Wörter auf analoge Weise gebildet werden, was die Ausbildung und Festigung des Musters selbst zur Folge hat. 

\section{Zusammenfassung Wortbildung}\label{section:Zfsg_Wortbildung}

Bevor wir uns im folgenden Kapitel~\ref{kap:Phonetik} mit den Lauten der gesprochenen Sprache, ihrer Klassifizierung und Bildung beschäftigen, fassen wir das Wichtigste zum Thema Wortbildung in Stichpunkten zusammen. 
 
\begin{itemize}

\item Wortbildungsarten: Zur Bildung komplexer Wortstämme aus einfachen Wortstämmen bedienen wir uns primär der Komposition, der Derivation und der Partikelverbbildung. Die Wortbildung ist regelhaft, analogiebasiert aber weniger vorhersagbar als die Flexion und kann -- anders als die Flexion -- zur Veränderung der Wortart führen.

\item Morphem: Die kleinsten bedeutungs- und funktionstragenden Einheiten eines komplexen Wortes nennt man Morpheme. Man unterscheidet zwischen lexikalischen Morphemen, die durch Einbettung in andere lexikalische Morpheme oder mithilfe von grammatischen Morphemen (Affixe) abgewandelt werden.

\item Komposition ist die Einbettung einer Kompositionsstammform als Erstglied in eine lexikalische Stammform als Zweitglied. Kompositionsstammformen können über mehr oder weniger motivierte Fugenelemente verfügen. Die rechtsstehende lexikalische Stammform ist der Kopf des Kompositums. Er legt die Wortart des Kompositums fest. Der produktive Prototyp ist das Determinativkompositum. Die Kompositionsstammform bestimmt (determiniert) den vom Zweitglied bezeichneten Typ näher.

\item Derivation ist die Ableitung einer lexikalischen Derivationsstammform durch Suffixe (Suffigierung), Präfixe (Präfigierung), durch den Wechsel des Stammvokals (implizite Derivation) sowie die Ableitung eines Stammes in eine andere Wortart ohne Formsignal (Konversion).

\item Suffigierung leitet eine lexikalische Derivationsstammform durch ein Suffix ab. Suffixe verbinden sich produktiv nur mit Stämmen bestimmter Wortarten (Selektion) und legen als morphologische Köpfe die Wortart der Gesamtbildung fest. Die Suffigierung dient primär der reihenhaften Bildung von Substantiven und Adjektiven.

\item Präfigierung dient primär der Formung von Verben durch Präfixe. Präfixe sind keine morphologischen Köpfe. Sie führen häufig zu einer vorhersagbaren, reihenbildenden Veränderung der Valenz und im Randbereich auch zur Verbalisierung nicht verbaler Stämme.

\item Implizite Derivation ist eine nicht mehr produktive Ableitung eines Wortstamms durch den musterhaften Wechsel des Stammvokals durch Um- und Ablaute. Sie diente der semantischen Modifikation und des Wortartwechsels.

\item Konversion ist die Ableitung eines Wortes in eine andere Wortart ohne Formsignal. Man unterscheidet zwischen der Ableitung lexikalischer Wortstämme (morphologische Konversion) und der Ableitung flektierter Stämme (syntaktische Konversion) in eine andere Wortart. Bei der morphologischen Konversion ist aus heutiger Sicht die Ableitungsrichtung oft unklar.

\item Partikelverbbildung dient ähnlich wie die verbale Präfigierung der Formung von Verben. Sie führt häufig zu einer vorhersagbaren, reihenbildenden Veränderung der Valenz und im Randbereich auch zur Verbalisierung nicht verbaler Stämme. Anders als Präfixe sind Verbpartikeln nicht fest mit dem Verbstamm verbunden. In Verberst- und Verbzweitsätzen besetzen sie die rechte Satzklammer.

\item Die Transparenz und Zahl bestehender Wortbildungen bedingt die Möglichkeit von Neubildungen auf analoge Weise und führt zur Ausbildung von Wortbildungsmustern bzw. Wortbildungsregeln. Die mögliche Produktivität eines Wortbildungsmusters wird durch verschiedene Faktoren eingeschränkt.

\end{itemize}

Der folgende letzte Abschnitt bietet Ihnen die Gelegenheit, durch Übungen zu lernen.

\section{Beispielanalysen}\label{section:BspAnalysenWB}

Im Folgenden analysieren wir Wortformen, die auf komplexen Stammformen beruhen. Dabei kommentieren wir unser Vorgehen und weisen auf Fehlerquellen hin. Gegebenenfalls zeigen wir Alternativen und ihre Konsequenzen auf. Wir raten Ihnen, die Wörter zuerst alleine zu analysieren und Ihr Ergebnis dann mit den Beispielanalysen zu vergleichen. 

\subsection{Beispiel 1}\label{subsection:Bsp1Wb}

Beginnen wir mit der Wortform \textit{Unannehmbarkeiten} in (\ref{ea:Unannehmbarkeit}).

\ea \label{ea:Unannehmbarkeit} Solche Unannehmbarkeiten trieben schon den Mathematiker Carl Friedrich Gauß: Egal wann wir geboren werden, es ist immer zur falschen Zeit.\footnote{Der Tagesspiegel, 23.10.2012.}
\z 

Zerlegen Sie zuerst die Wortform \textit{Unannehmbarkeiten} in seine morphologischen Bestandteile, benennen Sie diese und zeigen Sie danach die hierarchische Strukturiertheit in einer Baumdarstellung. Geben Sie dabei den jeweiligen Wortbildungsprozess an.


Kommen wir nun zu der Analyse: Es ist relativ einfach, \textit{Unannehmbarkeiten} in seine Bestandteile zu zerlegen. Es besteht aus dem Präfix \textit{un-}, der Partikel \textit{an}, dem Verbstamm \textit{nehm}, aus den Suffixen \textit{-bar} und \textit{-keit} sowie aus dem Flexiv \textit{-en}. Damit wissen wir aber noch nichts über die interne Strukturiertheit. Diese ergibt sich erst, wenn wir die Beziehungen der Bestandteile zueinander analysieren. Dabei stellen sich \zB die folgenden Fragen: Verbindet sich das Suffix \textit{-bar} mit dem Verbstamm \textit{nehm} oder mit dem komplexen Partikelverbstamm \textit{annehm}? Verbindet sich das Präfix \textit{un-} mit dem Adjektiv \textit{annehmbar} oder mit dem Substantiv \textit{Annehmbarkeit} oder gar mit  \textit{Annehmbarkeiten}?

Ohne Frage besteht der erste Bildungsschritt in der Verbindung der Verbpartikel \textit{an} mit dem Verbstamm \textit{nehm} zu dem komplexen Verbstamm \textit{annehm}. Denn Verbpartikeln verbinden sich, wie ihr Name sagt, mit Verben oder führen bei der Verbindung mit Adjektiven oder Substantiven zu deren Verbalisierung. Wie im Abschnitt \ref{subsection:Pt_Kopf} dargestellt, formen Partikeln Verben. Deshalb ist die Verbindung mit dem Adjektiv \textit{nehmbar} zu dem Adjektiv \textit{annehmbar} nicht möglich. 

Danach wird der Verbstamm \textit{annehm} als Derivationsstammform mittels \textit{-bar} zu dem Adjektiv \textit{annehmbar} deriviert. Im nächsten Schritt wird das Adjektiv durch das Präfix \textit{un-} zu \textit{unannehmbar} deriviert. Die Wortart wird nicht verändert, da nicht-verbale Präfixe keine Köpfe sind. Im nächsten Schritt erst wird das komplexe Adjektiv \textit{unannehmbar} als Derivationsstammform durch das Suffix \textit{-keit} abgeleitet. Das Suffix fungiert als morphologischer Kopf, weshalb das so gebildete Wort \textit{Unannehmbarkeit} ein feminines Substantiv ist. 

Gegen die Verbindung von \textit{un-} mit dem Substantiv \textit{Annehmbarkeit} spricht erstens, dass die Präfigierung komplexer Adjektive mit \textit{un-} weitaus produktiver ist als die von Substantiven und zweitens, dass manche Substantive überhaupt erst von der durch \textit{un-} abgeleiteten komplexen Adjektivbasis bildbar sind (vgl. Abschnitt~\ref{subsubsection:Präfix_Nomen}). 

Die Stammform \textit{Unannehmbarkeit} wird im letzten Schritt flektiert und damit für eine ganz bestimmte Funktion in einer Phrase vorbereitet. Das Flexiv \textit{-en} markiert den Plural. Die flektierte Form \textit{Unannehmbarkeiten} besetzt im Pluralparadigma die Plätze für alle Kasus. In (\ref{ea:Unannehmbarkeit}) handelt es sich um eine Nominativform.   
Die Analyse stellen wir in der Abbildung~\ref{caption:Unannehmbarkeit} dar.

\begin{figure}

	\resizebox{.62\textwidth}{!}{
	\begin{forest}
	[Bildung: Flexion\\Substantiv Fem/Pl/Nom\\\textit{Unannehmbarkeiten} [Bildung: Derivation\\Stammform\\Substantiv Fem\\ \textit{Unannehmbarkeit}  [Derivationsstammform\\Adjektiv\\\textit{unannehmbar} [Bildung: Derivation\\Stammform\\Adjektiv\\\textit{unannehmbar} [Präfix\\formt Adjektiv \\\textit{un-}] [Derivationsstammform\\Adjektiv\\ \textit{annehmbar} [Bildung: Derivation\\Stammfom\\Adjektiv\\\textit{annehmbar} [Derivationsstammform\\Verb\\\textit{annehm} [Bildung: Partikelverbbildung\\Stammform\\Verb\\\textit{annehm} [Verbpartikel\\ formt Verb\\\textit{an}] [Stammform\\Verb\\\textit{nehm}] ] ] [Suffix\\bildet Adjektiv\\\textit{-bar}]  ]  ]  ] ] [Suffix\\bildet Substantiv Fem\\ \textit{-keit} ]] [Flexiv\\ Fem/Pl/Nom\\\textit{-en}]] 
	\end{forest}
	}
	\caption{Konstituentenstruktur von \textit{Unannehmbarkeiten} }\label{caption:Unannehmbarkeit}
\end{figure}


\subsection{Beispiel 2}\label{subsection:Bsp2Wb}

Als nächstes analysieren wir die Struktur der Wortform \textit{Weltumseglungen} aus (\ref{ea:Weltumseglungen}).

\ea \label{ea:Weltumseglungen} In der Zeit der Weltumseglungen entdeckten europäische Forscher auf unbewohnten Inseln, dass die dort lebenden Tiere keine Scheu vor ihnen hatten.\footnote{die tageszeitung, 02.07.2018.}
\z 

Zerlegen Sie zuerst die Wortform \textit{Weltumseglungen} in seine morphologischen Bestandteile, benennen Sie diese und zeigen Sie danach die hierarchische Strukturiertheit in einer Baumdarstellung. Geben Sie dabei den jeweiligen Wortbildungsprozess an.


Kommen wir nun zu der Analyse: Zuerst zerlegen wir \textit{Weltumseglungen} in seine Bestandteile: in das Substantiv \textit{Welt}, das Präfix \textit{um-}, den Verbstamm \textit{segel}, das Suffix \textit{-ung} und das Flexiv \textit{-en}. Die Schwa-Tilgung zu \textit{Umseglung} (statt \textit{Umsegelung}) ist kein wortbildender, sondern ein phonologischer Prozess (vgl. Abschnitt~\ref{sec:Schwatilgung}). Eine Frage, die man hier stellen könnte, ist, ob der Verbstamm \textit{segel} eine Konversion des Substantivs \textit{Segel} ist. In der Tat wurde das Verb vom Substantiv abgeleitet.\footnote{Deutsches Wörterbuch von Jacob Grimm und Wilhelm Grimm, digitalisierte Fassung im Wörterbuchnetz des Trier Center for Digital Humanities, Version 01/21, \url{https://www.woerterbuchnetz.de/DWB?lemid=S24623}, abgerufen am 22.07.2021.} Wie im Abschnitt~\ref{subsubsection:MorphKonv} dargestellt, gehen wir bei etablierten Konversionen ohne klare Indizien für eine Konversionsrichtung von doppelt klassifizierten Stämmen aus, leiten also nicht mehr das Verb aus dem Substantiv ab.

Kommen wir nun zu den Beziehungen zwischen den Konstituenten. Hier stellt sich die Frage, ob sich \textit{um-} mit dem Verbstamm \textit{segel} oder mit dem abgeleiteten Substantiv \textit{Seglung} verbindet. Gegen letzteres spricht erstens, dass sich \textit{um-} als Verbpräfix mit Verben verbindet oder aber wie bei \textit{umgarnen} zur Ableitung von Verben dient. Zweitens spricht dagegen, dass \textit{Seglung} nicht (mehr) im Gebrauch ist.\footnote{Einige Belege finden sich in den Historischen Korpora (1465–1969) des Digitalen Wörterbuchs der deutschen Sprache: \url{https://www.dwds.de/r/?corpus=dtaxl&q=\%40Seglung}. Abgerufen am 22.07.2021.} Folglich beginnt die Analyse bei der Ableitung der verbalen Derivationsstammform \textit{segel} mit dem Präfix \textit{um-}. Der komplexe Verbstamm \textit{umsegel} wird als Derivationsstammform mittels \textit{-ung} zum Substantiv \textit{Umseglung} abgeleitet. Das komplexe Substantiv wird im Rahmen der Komposition mit der Kompositionsstammform \textit{welt} verbunden. Das Substantiv \textit{Weltumseglung} ist ein Determinativkompositum. Das Erstglied \textit{Welt} determiniert den Kopf \textit{Umseglung}. Wir interpretieren das Erstglied als Objekt des Verbs \textit{umsegeln}. Dieses Determinationsverständnis wird als Rektionskompositum bezeichnet.

Ganz zum Schluss wird das so gebildete Substantiv als Stammform (Flexionsstammform) mit \textit{-en} flektiert, wodurch der Plural markiert wird. Dass es sich um eine Genitivform handelt, geht erst aus dem Kontext in (\ref{ea:Weltumseglungen}) hervor. Die Flexionsform \textit{Weltumseglungen} besetzt im Pluralparadigma auch die Stellen für den Nominativ, Akkusativ und Dativ. 

Die Konstituentenstruktur stellen wir in der Abbildung~\ref{caption:Weltumseglungen} dar.

\begin{figure}
	\centering
	\begin{forest}
		[Bildung: Flexion\\Substantiv Fem/Pl/Gen\\ \textit{Weltumseglungen} [Bildung: Komposition\\Stammform\\Substantiv Fem\\ \textit{Weltumseglung} [Kompositionsstammform\\Substantiv Fem\\ \textit{welt}] [Stammform\\Substantiv Fem\\ \textit{Umseglung} [Derivationsstammform\\Verb\\ \textit{umsegel} [Bildung: Derivation\\Stammform\\ Verb\\ \textit{umsegel} [Präfix\\formt Verb\\ \textit{um-}] [Derivationsstammform\\Verb\\ \textit{segel}] ] ] [Suffix\\bildet Substantiv Fem\\ \textit{-ung}]  ]] [Flexiv\\ Fem/Pl/Gen \\ \textit{-en}] ]
	\end{forest}
	\caption{Konstituentenstruktur von \textit{Weltumseglungen} }\label{caption:Weltumseglungen}
\end{figure}



\subsection{Beispiel 3}\label{subsection:Bsp3Wb}

Als nächstes analysieren wir die Struktur von \textit{Kinderarmutsberichtes} aus (\ref{ea:Kinderarmutsberichtes}).

\ea \label{ea:Kinderarmutsberichtes} Kinderarmut ist ein Skandal in einer reichen Gesellschaft. So das Fazit nicht nur des letzten Kinderarmutsberichtes der Caritas.\footnote{Mannheimer Morgen, 22.08.2009.}
\z 

Zerlegen Sie zuerst die Wortform \textit{Kinderarmutsberichtes} in seine morphologischen Bestandteile, benennen Sie diese und zeigen Sie danach die hierarchische Strukturiertheit in einer Baumdarstellung. Geben Sie dabei den jeweiligen Wortbildungsprozess an.


Kommen wir nun zu der Analyse: Auch hier lassen sich die Konstituenten erst einmal recht leicht bestimmen. Das Wort besteht aus den Kompositionsstammformen  \textit{Kinder} und \textit{Armuts}, dem Präfix \textit{be-}, dem Verbstamm \textit{richt} und dem Flexiv \textit{-es}. Wichtig ist es, die Kompositionsstammformen (vgl. Abschnitt~\ref{subsection:Stämme}) nicht mit Flexionsformen zu verwechseln. Während das Fugenelement bei \textit{Kinder} durchaus durch die Flexion motiviert ist und deshalb als paradigmatisch (vgl. Abschnitt~\ref{subsection:Fuge}) bezeichnet wird, ist dies bei \textit{Armuts} nicht der Fall. Diese Form existiert im Flexionsparadigma des femininen Substantivs \textit{Armut} nicht.

Nun kommen wir zur internen Struktur. Die Wortform \textit{Kinderarmutsberichtes} ist primär ein Kompositum, das aus der Kompositionsstammform \textit{kinderarmuts} und dem Substantiv \textit{Bericht} besteht. Rein formal könnte es auch in die Kompositionsstammform \textit{kinder} und das Substantiv \textit{Armutsbericht} zerlegt werden. Dagegen spricht unsere Interpretation: Es geht um einen Bericht über Kinderarmut und nicht um einen Armutsbericht, der in irgendeiner Beziehung zu Kindern steht. Die Bildung ist ein Determinativkompositum. Dessen Interpretation ist verbal motiviert: \textit{über Kinderarmut berichten}, weshalb es auch ein Rektionskompositum ist.

Das Erstglied und das Zweitglied sind komplex. Der Kompositionsstammform \textit{kinderarmuts} liegt das Kompositum \textit{Kinderarmut} zugrunde. Dieses besteht aus der Kompositionsstammform \textit{kinder} und dem Substantiv \textit{Armut}. Durch die Kenntnis von \textit{arm} als Adjektiv könnte man versucht sein, \textit{Armut} in \textit{arm} und \textit{-ut} bzw. \textit{-mut} (analog zu \textit{Sanft-mut}) zu zerlegen. Aber dahinter steht kein produktives Wortbildungsschema, weshalb \textit{Armut} nicht sinnvoll in Konstituenten zerlegt werden kann.\footnote{Wahrscheinlich handelt es sich um die Substantivierung des althochdeutchen Adjektivs \textit{aramuoti}, wobei \textit{muoti} ein Adjektiv ist, welches als Zweitglied in Komposita gebraucht wird. Cf. \citet[104--105]{Paul1920Wb}.}


Der Kopf des Kompositums \textit{Bericht} beruht auf der morphologischen Konversion (vgl. Abschnitt~\ref{subsubsection:MorphKonv}) des Verbstamms \textit{bericht}. Für diese Konversionsrichtung gibt es ein klares morphologisches Indiz. Denn das Präfix \textit{be-} ist primär ein Bestandteil von Verben, nicht von Substantiven. Deshalb ist es plausibel, die Konversion in der Analyse zu berücksichtigen und das Substantiv \textit{Bericht} auf den komplexen Verbstamm \textit{bericht} zurückzuführen. Dieser Schritt führt zu der Bestimmung der Gesamtbildung als Rektionskompositum. 

Den komplexen Verbstamm \textit{bericht} zerlegen wir in das Präfix \textit{be-} und den Verbstamm \textit{richt}. Rein formal ist die Bildung transparent. Semantisch ist sie jedoch nicht mehr transparent. Denn die heutige Bedeutung von \textit{richten} als \gq{in gerade Form} bzw. \gq{Ausrichtung bringen} findet sich nicht mehr in dem Derivatum \textit{berichten}. Das liegt daran, dass \textit{berichten} ursprünglich (bei seiner Bildung) die Bedeutung \gq{in Ordnung bringen/korrigieren} hatte, die von \textit{berichtigen} übernommen wurde.\footnote{\textit{berichten} in: Wolfgang Pfeifer et al.: \textit{Etymologisches Wörterbuch des Deutschen}, 1993. Digitalisierte und von Wolfgang Pfeifer überarbeitete Version im Digitalen Wörterbuch der deutschen Sprache: \url{https://www.dwds.de/wb/berichten}, abgerufen am 26.07.2021.}

Das Kompositum \textit{Kinderarmutsbericht} wird im letzten Schritt mit \textit{-es} flektiert, das im Singularparadigma die Position für den Genitiv markiert.  

Die Konstituentenstruktur stellen wir in der Abbildung~\ref{caption:Kinderarmutsberichtes} dar.

\begin{figure}
	\centering
    \resizebox{\textwidth}{!}{
	\begin{forest}
	[Bildung: Flexion\\Substantiv Mask/Sg/Gen\\ \textit{Kinderarmutsberichtes} [Bildung: Komposition\\ Stammform\\ Substantiv Mask\\ \textit{Kinderarmutsbericht} [Kompositionsstammform\\Substantiv Fem\\ \textit{kinderarmuts} [Bildung: Komposition\\Stammform\\Substantiv Fem\\ \textit{Kinderarmut} [Kompositionsstammform\\Substantiv Neutr\\ \textit{kinder}] [Stammform\\Substantiv Fem\\ \textit{Armut}] ] ] [Bildung: Konversion\\Stammform\\Substantiv Mask\\\textit{Bericht} [Bildung: Derivation\\Stammform\\ Verb\\ \textit{bericht} [Präfix\\formt Verb\\ \textit{be-}] [Derivationsstammform\\Verb\\ \textit{richt}] ] ] ] [Flexiv\\ Mask/Sg/Gen\\ \textit{-es}] ]
	\end{forest}
    }
	\caption{Konstituentenstruktur von \textit{Kinderarmutsberichtes} }\label{caption:Kinderarmutsberichtes}
\end{figure}



\subsection{Beispiel 4}\label{subsection:Bsp4Wb}

Zum Schluß analysieren wir die Struktur von \textit{fußballert} in (\ref{ea:fußballert}).

\ea \label{ea:fußballert} Abzusehen ist bereits, wie sich der Herr Sportlehrer in die Herzen seiner Jungs fußballert [\ldots{}]\footnote{Berliner Zeitung, 19.10.1994.}
\z 

Zerlegen Sie zuerst die Wortform \textit{fußballert} in seine morphologischen Bestandteile, benennen Sie diese und zeigen Sie danach die hierarchische Strukturiertheit in einer Baumdarstellung. Geben Sie dabei den jeweiligen Wortbildungsprozess an.


Kommen wir nun zu der Analyse: Auch \textit{fußballert} lässt sich erst einmal recht einfach in seine Konstituenten zerlegen, nämlich in die Substantive \textit{Fuß} und \textit{Ball} sowie in das Suffix \textit{-er} und in das Flexiv \textit{-t}. 

Damit ist noch nichts über die Struktur der Bildung gesagt, um die es nun geht. Die Bildung beginnt beim Kompositum \textit{Fußball}, das aus \textit{fuß} als Kompositionsstammform und \textit{Ball} als Stammform besteht. Dieses Substantiv wird als Derivationsstammform mit dem  Suffix \textit{-er} zu \textit{Fußballer} abgeleitet. Wie im Abschnitt~\ref{subsubsection:er} dargestellt wurde, dienen primär Verbstämme als Basis der Ableitung. Die Ausweitung des Musters auf Substantive ist semantisch begründet, wenn entsprechende Substantive als Nomina agentis wie \textit{Gärtn-er}, \textit{Töpf-er} oder eben \textit{Fußball-er} deutbar sind (vgl. Abschnitt~\ref{section:Produktivität}). 

Auf den ersten Blick handelt es sich bei der Bildung des Verbs \textit{fußballer}(\textit{n}) von \textit{Fußballer} um eine reine Konversion. Diese wäre semantisch gestützt, da der kognitive Schritt von einem Nomen agentis hin zur Tätigkeit (\textit{actio}) klein ist.

Im Abschnitt~\ref{subsubsection:MorphKonv} haben wir jedoch gesehen, dass für die morphologische Konversion nur nicht suffigierte Stämme in Frage kommen. Denn die affixlose Konversion kann explizite Suffixköpfe nicht überschreiben. Aber Verbableitungen von desubstantivischen \textit{er}-Ableitungen sind grundsätzlich möglich: \textit{gärtnern}, \textit{töpfern} und \textit{ballern}. Deshalb schlagen wir Folgendes vor: Es handelt sich um eine besondere Konversion, die von der Reanalyse von \textit{-er} gestützt wird. Denn \textit{-er} war auch ein verbbildendes Suffix wie in \textit{schlittern}, \textit{bibbern} oder \textit{knabbern}, deren Basen nicht (mehr) bekannt sind. Bei kausativen Bildungen wie \textit{einschläfern}, \textit{steigern} und \textit{folgern} sowie beim Vergleich von \textit{Doktor} mit (\textit{herum})\textit{doktern} ist das Verbsuffix \textit{-er} allerdings noch zugänglich. Die Reanalyse von \textit{-er}, das maskuline Nomina agentis bildet, zu einem verbalen Suffix stützt die Konversion von \textit{Fußballer} zu \textit{fußballer}(\textit{n}).

Im letzten Schritt wird der auf die beschriebene Weise abgeleitete Verbstamm \textit{fußballer} flektiert. Das Flexiv \textit{-t} besetzt im Präsens Indikativ Aktiv die Position der 3. Person Singular und der 2. Person Plural. In dem Beleg (\ref{ea:fußballert}) handelt es sich um die 3. Person Singular. Denn das Subjekt \textit{der Herr Sportlehrer} steht im Singular und korrespondiert mit dem Pronomen der 3. Person Singular \textit{er}.

Die Konstituentenstruktur stellen wir in der Abbildung~\ref{Abb:fußballert} dar.

\begin{figure}
	\centering
	\begin{forest}
		[Bildung: Flexion\\Verb 3. Pers. Sg/Präs/Ind/Akt\\ \textit{fußballert} [Stammform\\Verb\\\textit{fußballer} [Konversion\\Nomen agentis auf \textit{-er} $\rightarrow$ Verbsuffix \textit{-er} [Stammform \\ Substantiv Mask\\ \textit{Fußballer} [Derivationsstammform \\ Substantiv Mask \\ \textit{fußball} [Stammform \\ Substantiv Mask \\\textit{Fußball} [Kompositionsstammform \\ Substantiv Mask \\\textit{fuß}] [Stammform \\ Substantiv Mask \\ \textit{Ball}] ] ] [Suffix \\ bildet Substantiv Mask \\\textit{-er}] ] ] ] [Flexiv\\3. Pers. Sg/Präs/Ind/Akt\\\textit{-t}]]
	\end{forest}
	\caption{Konstituentenstruktur von \textit{fußballert}}\label{Abb:fußballert}
\end{figure}

